\documentclass[11pt]{article}
\usepackage{amsmath,amssymb,amsthm,mathrsfs}
\usepackage[margin=1.1in]{geometry}
\usepackage{hyperref}
\usepackage{booktabs}
\usepackage{multirow}

\theoremstyle{plain}
\newtheorem{theorem}{Theorem}
\newtheorem{lemma}[theorem]{Lemma}
\newtheorem{proposition}[theorem]{Proposition}
\newtheorem{corollary}[theorem]{Corollary}
\theoremstyle{definition}
\newtheorem{definition}[theorem]{Definition}
\newtheorem{hypothesis}[theorem]{Hypothesis}
\theoremstyle{remark}
\newtheorem{remark}[theorem]{Remark}

\newcommand{\R}{\mathbb{R}}
\newcommand{\C}{\mathbb{C}}
\newcommand{\N}{\mathbb{N}}
\newcommand{\Z}{\mathbb{Z}}
\newcommand{\GR}{\Gamma_{\R}}
% Identifiers are \detokenize'd throughout.  Two reasons: they may be written with raw
% underscores, and -- crucially -- the resulting characters carry a correct ToUnicode
% mapping, so every identifier in this paper can be copied out of the PDF verbatim.
% (\texttt{a\_b} renders correctly but extracts as "a b", which silently corrupts any
% name a reader tries to paste into a source file or a search.)
\newcommand{\lean}[1]{\texttt{\detokenize{#1}}}
\newcommand{\leanref}[1]{%
  \par\nobreak\vspace{2pt}%
  {\footnotesize\raggedright\noindent\textsc{Lean:}~\texttt{\detokenize{#1}}\par}%
  \vspace{2pt}}
\DeclareMathOperator{\ord}{ord}

\title{A carrier readout for the zeros of $\zeta$,\\
and a machine-checked index-to-ordinate extractor}
\author{Samuel Lavery}
\date{Draft --- \today}

\begin{document}
\maketitle

\begin{abstract}
We work in a three-dimensional model and descend to the classical one-dimensional
setting only as a readout. The stage is a double-ended Archimedean helix with
closed-form arclength; a \emph{carrier} --- a numberline scaled by a harmonic constant
$\Delta>0$ --- is laid along it, its sites at equal arclength $x_n=\Delta n$; and the
$L$-function rides the carrier as a \emph{fiber}, whose phasors the conductor sorts into
channels that cancel \emph{focally}, at a height, leaving no residue. The
one-dimensional $L$-function is this object seen in a chart, reached by dropping the
radius (Cayley) and then the angle ($y=\log z$); the height survives, and the classical
ordinate is its registration.

In the chart the construction is a one-parameter family of \emph{carrier readouts}:
real-valued functions $u_\Delta$ on the critical line, the chart of the bank at each
height. In the chart that readout is Hardy's $Z$ up to a strictly positive factor, and
the isolation results below are what the chart reads: they are stated and proved as
properties of the readout, and they coincide with the classical argument, which is the
form the chart gives them. \S\ref{sec:loadbearing} sets out what lives only in three
dimensions. One item is not available in the chart at all: in the ordinate variable the distance off the critical
line is a \emph{radial} coordinate, so a certified contour count over a band
$\lvert\beta-\tfrac12\rvert\le h$, matched against certified sign changes, certifies
that no zero in a window has left the line by as much as $h$
(Theorem~\ref{thm:bandverify}), and an elementary certificate on the sliver near
$\operatorname{Re}s=1$ removes the band, so that every zero with ordinate in the window
lies on the line (Corollary~\ref{cor:bandfree}) --- a statement about $\beta$, which the
readout's sign provably cannot reach. We prove, unconditionally and
with machine-checked proofs, that the warped fiber on the scaled carrier \emph{is} the
one-dimensional function --- its ordered chart equals $\Delta^{-s}(1-2^{1-s})\zeta(s)$ on
the whole channel $\operatorname{Re}s>0$, on the three-dimensional phasors and at every
carrier scale, and determines $\zeta$ everywhere by the identity theorem
(Theorem~\ref{thm:fiberisfunction}); that $u_\Delta$ is a strictly positive multiple of
$\Lambda(\tfrac12+it)$, so its zeros are exactly the on-line zeros of $\zeta$ at every
carrier scale; that its sign changes by $(-1)^m$ across an event of multiplicity $m$,
with no simplicity hypothesis; that zeros off the line pair at a common ordinate, so the
parity of the \emph{full} strip count is the readout's sign change; and that a sign
change across an interval on which $\Xi'$ does not vanish forces exactly one event
there, simple. Chaining this yields the index: the $i$-th cell contains the $i$-th
event, uniquely, and a sign-change sub-bracket of that cell encloses a zero of $\zeta$
on the critical line. Every hypothesis is an evaluation at finitely many points; no
counting routine, zero-density estimate, or bound on $S(t)$ enters.

Separately, and off that path, the \emph{abscissa} of a zero proves recoverable from
data computed only on the line: the bank's log-jets are exactly a power sum over the
fiber's zeros, a zero $\beta+i\gamma$ becomes a node $\gamma-i(\beta-\tfrac12)$, and a
matrix pencil returns those nodes as eigenvalues. The reading is calibrated against
off-line zeros the fiber carries in closed form.

The development is formalised in Lean~4/Mathlib --- $100$ theorems in the chart chain,
with the geometry, its descent and ten supporting modules besides; zero \lean{sorry},
zero \lean{axiom}. An implementation is reported and validated at $297$ indices, and
its own hypotheses are discharged by computation against an explicit Euler--Maclaurin
remainder. Two propositions earlier carried as open hypotheses are \emph{disproved}
rather than assumed. \S\ref{sec:notclaimed} states precisely what is not claimed.

\medskip
\noindent\textbf{Contribution.} This is a geometric model for the Dirichlet
$L$-functions, and it computes. Its objects are a \emph{scaled carrier} --- a
numberline reparametrized by a harmonic constant on its own integer lattice, the
sites at arclength $x_n=\Delta n$ --- and a \emph{warp}, a reparametrization by a
function: dynamic, readout-preserving, of unit modulus. The warp's generator is a
self-adjoint operator on $\ell^2$ with the phasors as eigenbasis; the fiber is its flow,
$\zeta$ is its trace, and the critical line is the boundary of the Hilbert space in which
the fiber lives (\S\ref{sec:hamiltonian}). The fiber rides the carrier;
the lanes are the conductor's $\pm$ channels; the weld gauge is the fiber's own
argument, and it absorbs the root number, so the readout is real for every primitive
character and not only for the self-dual ones. The model is carrier-generic: reality,
the functional equation and the counts are properties of the involution and the
readout map, established once for every fiber that rides the carrier, with $\zeta$ the
case formalised here. From that geometry --- addressed by index, with no Gram point and
no asymptotic expansion needed to \emph{define} the readout --- one receives, given a
count satisfying the interface of \S\ref{sec:extractor}, a certified interval containing
the $i$-th on-line zero, and the chain down to $\zeta(\tfrac12+i\gamma)=0$ is
machine-checked; \S\ref{sec:notclaimed}(iv) says exactly how much that certifies. (The Riemann--Siegel formula does appear
later, but only as an interchangeable \emph{evaluator} of the readout's sign,
\S\ref{sec:obligations}(O4); nothing in the model or the theorems depends on it.)
Geometric readings of $\zeta$ have to this point been descriptive; this one is
operational.

\medskip\noindent\textbf{Status of the obligations.} A companion aim of this paper is
that it should carry no unproved hypothesis and no unbacked measurement. Every theorem
here is unconditional, no theorem cited takes an unproved proposition as an argument,
and the implementation's own hypotheses are now discharged by computation: the sign
evaluations are certified against an explicit Euler--Maclaurin remainder, the
non-vanishing of $\Xi'$ is certified on the intervals returned, and the zero count is
certified by contour, all from an explicit region bound on $\lvert\Xi\rvert$ derived
here rather than quoted. Two propositions that earlier drafts carried as open
hypotheses are \emph{disproved} instead: the $\mu_6$ sector balance law
(Theorem~\ref{thm:arcobstruction}) and the fixed-cutoff clock bound
(Proposition~\ref{prop:o2false}). \S\ref{sec:obligations} gives the status of each item
and what closed it.
\medskip
\noindent\textbf{Keywords.} Riemann zeta function; zeros on the critical line;
certified computation; formalized mathematics; Lean 4; Prony's problem.

\medskip
\noindent\textbf{2020 Mathematics Subject Classification.}
Primary 11M06, 11Y35; Secondary 11M26, 68V20.

\medskip
\noindent\textbf{Code availability.} The Lean sources (\lean{RequestProject/}), the
implementation and every numerical check named here live in the project repository, and
\S\ref{sec:repro} lists what each one runs and the decision rule it was given before it
ran. They are \emph{not} embedded in this PDF: a reader sent the PDF alone cannot run
any of it, and no claim here should be taken on the strength of a script the reader
cannot see.
\end{abstract}

\tableofcontents

%======================================================================
\section{Introduction}\label{sec:intro}

\subsection*{What is constructed}
The object is a three-dimensional one, and the order of exposition follows it. There
are three layers, and they are kept apart throughout because conflating them is what
makes the subject look one-dimensional.

\emph{The geometry} (\S\ref{sec:geometry}) is the stage: a double-ended Archimedean
helix
\[
  \gamma(k)\;=\;\bigl(rk\cos2\pi k,\;rk\sin2\pi k,\;pk\bigr)
\]
together with its conjugate anti-helix through the origin, both running
$0\to+\infty$. It has a closed-form arclength and no marks. Scaling does not act here.

\emph{The carrier} (\S\ref{sec:carriergeom}) is a numberline laid along the geometry
--- a ruler --- and it is where harmonic scaling acts: the sites sit at equal
\emph{arclength} $x_n=\Delta n$ for a harmonic constant $\Delta>0$. One geometry can
host several such rails at once.

\emph{The fiber} (\S\ref{sec:fiber}) is what rides the carrier: the $L$-function's
representation, carrying a phasor at each site consumed. Phasors are born at zero
magnitude and grow; they spin in the logarithm because the ruler they sit on turns.
The conductor sorts them into positive and negative channels --- for $\zeta$, the even
and odd sites, which is what a conductor of $2$ supplies. A \emph{warp} --- scaling by
a function, unit modulus --- moves the fiber through configuration space and sets the
height at which its channels cancel. That cancellation is \emph{focal}: the channels
point at a common height and close there, exactly, leaving no residue. It is an event
at a height, not a root of a formula.

Only then does the one-dimensional $L$-function appear, and it appears as a
\emph{readout} (\S\ref{sec:descent}): $3\text{-D}\to2\text{-D}$ by Cayley, dropping the
radius; $2\text{-D}\to1\text{-D}$ by $y=\log z$, dropping the angle. The height
survives both steps, and the classical ordinate is its registration --- the two
coincide, and are not the same quantity. The logarithm of the entire model enters once,
in the second step. That the geometric event and the analytic one coincide is proved,
as an \emph{iff}, for non-principal real characters (Theorem~\ref{thm:focal}); for
$\zeta$ that particular form degenerates and the licence is instead an exact
factorisation of the fiber on the line (Theorem~\ref{thm:weld}), with
Remark~\ref{rem:zetabridge} recording precisely how much of the geometric form is
proved for $\zeta$ and how much is a named hypothesis.

In the chart (\S\ref{sec:carrier} onward) the fiber is a complex number at each height
which turns as the height advances; the \emph{weld gauge} is its own argument, read off
from the archimedean local factor, and dividing it out leaves a real number --- the
\emph{carrier readout}. The rest of the paper is statements about where that readout
vanishes and how its sign changes there, all invariant under $\Delta$: the carrier
scale is a gauge, not a parameter of the answer.

\subsection*{Contribution}\label{sec:contribution}
\emph{A three-dimensional model, and the one-dimensional function is its chart.} The
layers are set out in \S\S\ref{sec:geometry}--\ref{sec:descent} and not repeated here.
What the ordering buys is that each classical difficulty can be sorted into one of two
boxes --- a fact about the object, or a fact about the projection --- and the sorting is
not a matter of taste. The abscissa of convergence, the strip, and the devices for
reconstructing the functional equation from inside the projection are all in the second
box; the double end, the focal cancellation and the height are in the first.

Two structures of the model are distinguished throughout because conflating them is the
characteristic error. The \emph{carrier scale} is scaling by a harmonic
\emph{constant} $\Delta$, acting on the ruler: the sites sit at arclength
$x_n=\Delta n$. The \emph{warp} is scaling by a \emph{function} --- dynamic, unit
modulus, readout-preserving --- acting on the fiber; its radial generator is the
log-weight operator $H$ multiplying the $n$-th phasor by $\log n$, and on the carrier of
scale $\Delta$ that generator is $H_\Delta=H+\log\Delta\cdot I$, the shift by a constant
being the whole of the carrier's effect on it. That sentence is a theorem and not a
description: $H$ is an unbounded self-adjoint operator on $\ell^2(\mathbb N,\mathbb C)$
on its maximal domain, the phasors are a Hilbert basis of its eigenvectors,
$H_\Delta=H+\log\Delta\cdot I$ holds with equal domains, and the warp $e^{-sH}$ is its
functional calculus on that basis (\S\ref{sec:hamiltonian}). The carrier scale is a gauge --- the
events and the readout's sign are invariant across the family $\Delta>0$. That
invariance is not a theorem with content: $\Delta^{-s}$ cancels between the fiber and the
gauge, so $u_\Delta/u_{\Delta'}=(\Delta'/\Delta)^{1/2}$, a positive constant independent
of $t$. It is worth stating precisely because it says the model has \emph{no free
parameter to tune} --- an honest design property, not a discovery. The warp is different:
it genuinely moves the height at which the cancellation completes. There is no inverse warp; the reciprocal of
a warp is again a warp.

\emph{It is a model for the Dirichlet $L$-functions, not for $\zeta$ alone.} A fiber
rides the carrier; the lanes are the conductor's $\pm$ channels, supplied by the
character. Reality, the functional equation and the counting laws are properties of
the weld involution and the readout map, and hold for every fiber that rides the
carrier --- $\zeta$ being the case whose formalisation is carried out here
(\S\ref{sec:general} records what changes for a general primitive $\chi$: the root
number enters as a fixed ray, which the weld gauge absorbs, and nothing else does).

\emph{The chart, and what does not fit in it.} The readout is the chart of the bank at a
height, and in the chart it is Hardy's $Z$ up to a strictly positive factor
(Proposition~\ref{prop:isZ}); the isolation chain of
\S\S\ref{sec:ordinate}--\ref{sec:capstone} is therefore what the chart reads, and it
coincides with the classical argument, which is not claimed to be otherwise. What lives
only in three dimensions is set out in \S\ref{sec:loadbearing}: the lane structure
supplies a series convergent \emph{on} the line, which is what the locator evaluates; the
same lane structure imposes a linear head, which is a cost and is reported as one; the
$\mu_6$ cell width is the constant in the refutation of Theorem~\ref{thm:arcobstruction};
the lane factor supplies a closed-form off-line calibration set for the abscissa reading;
the focal equivalence (Theorem~\ref{thm:focal}) is a theorem about the model for the
characters the $Z$-identification does not reach; the site weight is forced by the
carrier scale being a gauge and is the climber's registration of the site
(Theorem~\ref{thm:gauge}, \S\ref{sec:fiber}); the distance off the line is a radial
coordinate of the object, invisible to the chart (\S\ref{sec:bandverify}); and the fiber
has a self-adjoint generator whose Hilbert space ends at the critical line
(\S\ref{sec:hamiltonian}). Two of those are negative results about the carrier, which is
the part we would defend hardest.

\emph{And it computes.} On one side of the existing literature, geometric readings of
$\zeta$ --- the Dirichlet walk and its helical carrier, spirals of $\zeta$, phasor
accumulations in the Argand plane --- are descriptive: they supply pictures and
intuitions about where zeros can and cannot be, and no algorithm. On the other, the
methods that do take an index and return an ordinate --- Turing's method and its
modern form, Platt's isolation, Newton refinement in ball arithmetic,
Odlyzko--Sch\"onhage --- carry no geometric representation, working from the
Riemann--Siegel formula and its relatives. The model here is in both columns: one
addresses it by index and receives a certified interval --- certified relative to the
count supplied through the interface of \S\ref{sec:extractor}, as
\S\ref{sec:notclaimed}(iv) spells out --- and every step from the
representation down to $\zeta(\tfrac12+i\gamma)=0$ is machine-checked
(\S\ref{sec:formalization}). What is claimed is the model and its operation, not a
faster algorithm and not a new localization principle; \S\ref{sec:notclaimed} is
explicit about the difference.

\subsection*{What is proved}
First, that the object and its chart are the same thing: the lane-balanced chart of the
warped bank on the scaled carrier equals $\Delta^{-s}(1-2^{1-s})\zeta(s)$ on the whole
channel $\operatorname{Re}s>0$, on the three-dimensional phasors and at every carrier
scale, and any entire function agreeing with that chart on the channel is
$(1-2^{1-s})\zeta(s)$ everywhere (Theorem~\ref{thm:fiberisfunction}). The chain then runs
as follows, each link machine-checked. The weld gauge makes the readout
a strictly positive multiple of the completed zeta function on the line
(Theorem~\ref{thm:weld}), so the readout's zeros are exactly the on-line zeros of
$\zeta$, at every carrier scale (Proposition~\ref{prop:zeros}). Across an event of
multiplicity $m$ the readout's sign changes by $(-1)^m$, with no simplicity hypothesis
(Theorem~\ref{thm:signflip}): multiplicity is \emph{carried} through the argument
rather than assumed away. The zeros off the line occur in conjugate pairs at a common
ordinate with equal multiplicity, so the parity of the \emph{full} strip count --- not
merely the on-line count --- equals the readout's sign change
(Theorem~\ref{thm:stripparity}). On an interval across which the readout changes sign
and on which $\Xi'$ does not vanish, there is exactly one event and that event
is simple (Theorem~\ref{thm:countone}); chaining this across a ladder of ordinates
gives the index (Theorem~\ref{thm:chain}), and the capstone
(Theorem~\ref{thm:capstone}) converts a sign-change sub-bracket of the $i$-th cell into
an interval provably containing the $i$-th zero of $\zeta$ on the critical line. Every
hypothesis of the capstone is an evaluation at finitely many points: signs of the
readout at the ladder ordinates, and non-vanishing of $\Xi'$ between them
(Remark~\ref{rem:whichderiv} says why the second is stated on $\Xi$ and not on the
readout). No zero-counting routine, no zero-density estimate, and no bound on $S(t)$
appears anywhere in the chain.

Off that chain, the operator named in \S\ref{sec:contribution} is proved to be one
(\S\ref{sec:hamiltonian}): the log-weight operator is self-adjoint on $\ell^2$ with the
phasors as eigenbasis (Theorem~\ref{thm:logweight}); the warp is its functional calculus
and, on the line, its unitary group (Theorem~\ref{thm:generator}); $\zeta$ is its trace,
$-\zeta'/\zeta$ its von Mangoldt trace, and the Euler product its decomposition into
prime clocks (Theorem~\ref{thm:zetatrace}); and the fiber is its flow --- the bank
$B(s)=e^{-sH}\mathbb 1$, square-summable exactly to the right of the critical line, of
which every $L$-series is a chart, and whose entry into the kernel of the lane-balanced
chart is the chart's zero (Theorem~\ref{thm:bank}).

\subsection*{Relation to existing work}\label{sec:related}
\emph{Index to certified interval is not new, and we do not claim it.} Turing's method~\cite{Turing1953}
counts zeros on the line by bounding $\int_{t_1}^{t_2}S(t)\,dt$ over a suitable
interval, which converts a sign count into a certified total; its modern rigorous form
is due to Platt~\cite{Platt2017}, who isolated the non-trivial zeros with imaginary part up to
$3.061\times10^{10}$ to absolute precision $\pm2^{-102}$. That algorithm is implemented
and shipped. Arb~\cite{Johansson2017} provides
\lean{acb_dirichlet_hardy_z_zeros} and
\lean{acb_dirichlet_platt_hardy_z_zeros} take an index $n$ and a count and return
ball enclosures of consecutive zeros, the documented example retrieving zeros of index
$\approx1.048\times10^{9}$ near height $3.888\times10^{8}$ in about a quarter of a
second. A reader wanting the $i$-th ordinate fast should use that. What is offered here
is a different \emph{representation} of the same data, with the path from the
representation to $\zeta(\tfrac12+i\gamma)=0$ machine-checked.

\emph{Refinement.} Once a zero is isolated, Newton's method in ball arithmetic with
precision doubling refines it at the cost of roughly two full-precision evaluations at
any digit count. Johansson~\cite{Johansson2015} gives the rigorous high-precision
evaluation machinery this rests on; the record computation of the first ordinate to
over $3\times10^{5}$ digits is a separate artifact built on that machinery and on
Arb~\cite{Johansson2017}, not a result of~\cite{Johansson2015} itself, and is recorded
as the reference value in OEIS \texttt{A058303}. We make no speed claim against any of
it (\S\ref{sec:notclaimed}(i)).

\emph{Amortised evaluation.} Odlyzko and Sch\"onhage~\cite{OdlyzkoSchonhage1988} reduce
the cost of evaluating $\zeta$ at many equally spaced heights by transforming the
Dirichlet sum into a rational-function evaluation at roots of unity and applying the
FFT. That is an amortisation over a \emph{region}; it does not give random access to a
single index, which is what is addressed here.

\emph{The ancestor of the jet enclosure.} Gargantini and
Henrici~\cite{GargantiniHenrici1971} ask exactly the question of \S\ref{sec:jets}: if
the value of the logarithmic derivative $p'/p$, and of finitely many of its
derivatives, is known at a point, and the other zeros are known to lie in specified
circular regions, where is the remaining zero? They answer it in circular (complex
interval) arithmetic, connect it to a classical result of Laguerre, and --- this is the
relevant sentence --- apply the theory to locating zeros of certain \emph{transcendental}
functions. The jet enclosure below is that construction with the circular regions
supplied by the carrier's own tail bound, and it should be read as a descendant, not as
a new localization principle. The complementary classical device, power sums of the
zeros in a region by contour integration, is Delves and Lyness~\cite{DelvesLyness1967}.

\emph{The Hamiltonian reading.} That $\zeta$ is the partition function of a quantum
statistical system whose Hamiltonian has spectrum $\{\log n\}$ is Bost and
Connes~\cite{BostConnes1995}; the spectral side of the same frame --- an operator whose
eigenvalues are the ordinates --- is the Hilbert--P\'olya programme in the forms of Berry
and Keating~\cite{BerryKeating1999} and Connes~\cite{Connes1999}. In
\S\ref{sec:hamiltonian} the warp's generator of \S\ref{sec:contribution} is made a
self-adjoint operator on $\ell^2$; $\zeta$ is its trace, which is the Bost--Connes
identity for this operator, and the object the trace is taken of is the \emph{bank} ---
with its unitary orbit on the line, its $\ell^2$-boundary at the line, and every
$L$-series as one of its charts. No operator with the ordinates as spectrum is
constructed or claimed.

%======================================================================
\section{The geometry: the double-ended helix}\label{sec:geometry}

The object is three-dimensional, and this section fixes it before any series is
written. Nothing here is a picture attached to a one-dimensional argument: the
one-dimensional statements later in the paper are readouts \emph{of} this, taken in a
chart, and it is worth having the stage in hand first so that it is clear which
statements are about the object and which are about the chart.

\begin{definition}[the helix]\label{def:helix}
For a pitch $p>0$ and a radial rate $r>0$, the helix is the space curve
\[
  \gamma(k)\;=\;\bigl(r k\cos 2\pi k,\; r k\sin 2\pi k,\; p k\bigr),
  \qquad k\ge 0 ,
\]
Archimedean in the turn parameter $k$: radius $rk$, pitch $p$. The geometry is
\emph{double-ended} --- the helix together with its conjugate anti-helix, the mirror
image through the origin --- and both ends run $0\to+\infty$.
\end{definition}
\leanref{Geometry.helix, Geometry.helixVel, Geometry.helix_hasDerivAt,
Geometry.helix_cyl_radius}

\noindent
The cylindrical radius of $\gamma(k)$ is $\lvert rk\rvert$ and the speed is
\[
  \lvert\dot\gamma(k)\rvert^{2}\;=\;p^{2}+r^{2}+(2\pi r k)^{2},
\]
so the curve accelerates as it climbs: successive turns are longer, and that is the
whole source of the logarithm later.
\leanref{Geometry.helix_speed_sq, Geometry.speed}

\begin{definition}[the climber, and the height]\label{def:climber}
The \emph{climber} is $k(y)=e^{y}/p$. It registers a one-dimensional ordinate $y$
against a point of the geometry, at cylindrical height
\[
  z\;=\;p\,k(y)\;=\;e^{y},\qquad R(y)\;=\;r\,k(y)\;=\;(r/p)\,e^{y}.
\]
\end{definition}
\leanref{Geometry.kClimb, Geometry.height_eq_p_mul_kClimb, Geometry.radius,
Geometry.radius_eq, Geometry.gammaY, Geometry.gammaY_eq}

\noindent
So the object is Archimedean in its own turn parameter and exponential in the analytic
ordinate; both descriptions are correct, in different charts, and the climber is the
change of variable between them. \emph{The height is the geometric quantity; the
ordinate is its readout} (\S\ref{sec:descent}). We will say that a height and an
ordinate \emph{coincide}, never that they are equal --- they are quantities in
different charts, related by a registration.

\begin{theorem}[closed-form arclength]\label{thm:arclength}
For $r>0$,
\[
  S(k)\;=\;\int_0^{k}\lvert\dot\gamma\rvert
  \;=\;\frac{k}{2}\sqrt{p^{2}+r^{2}+4\pi^{2}r^{2}k^{2}}
      +\frac{p^{2}+r^{2}}{4\pi r}\,
       \operatorname{arsinh}\frac{2\pi rk}{\sqrt{p^{2}+r^{2}}} .
\]
For $r=0$ it degenerates to $S(k)=pk$.
\end{theorem}
\leanref{Geometry.arclength, Geometry.arclengthClosed, Geometry.arclength_closed_form,
Geometry.arclength_r_zero}

\noindent
The arclength is what the next section measures with, so it matters that it is in
closed form rather than an integral to be estimated: the ruler is exact.

%======================================================================
\section{The carrier: a numberline, harmonically scaled}\label{sec:carriergeom}

The geometry is the stage; it carries no marks. A \emph{carrier} is a numberline laid
along it --- a ruler --- and the carrier is where scaling acts. The distinction is not
decorative and the paper keeps it throughout:

\begin{center}
\begin{tabular}{@{}lll@{}}
\toprule
 & scaling by & acts on\\
\midrule
carrier scale & a harmonic \emph{constant} $\Delta$ & the carrier\\
warp & a \emph{function}, unit modulus & the fiber\\
\bottomrule
\end{tabular}
\end{center}

\noindent
Harmonic scaling does not act on the geometry. It multiplies the carrier's tick
spacing by a harmonic constant $\Delta$, and the carrier's sites then sit at equal
\emph{arclength} $x_n=\Delta n$ along the curve. The index is the arclength in units of
$\Delta$,
\[
  N(y)\;=\;\frac{S\bigl(k(y)\bigr)}{\Delta},
  \qquad\text{and at }\Delta=\pi/3,\qquad N(y)=\frac{3}{\pi}\,S\bigl(k(y)\bigr).
\]
\leanref{Geometry.Delta, Geometry.Nindex, Geometry.Nindex_eq}

\noindent
One geometry can host several carriers at once --- \emph{rails} --- differently
configured, running alongside one another in the same space without displacing each
other. Only one rail is used in this paper.

\begin{theorem}[the carrier scale forces the logarithm]\label{thm:gauge}
Call a weight $w$ on the carrier \emph{gauge} if rescaling the carrier by $\Delta>0$
shifts the weight of every site by one constant, $w(\Delta x)=w(x)+c(\Delta)$ for all
$\Delta,x>0$ --- that is, if the carrier scale acts on the fiber by a factor common to
all sites. Then:
\begin{enumerate}
\item[(i)] $\log$ is gauge, with $c(\Delta)=\log\Delta$;
\item[(ii)] a gauge weight continuous on $(0,\infty)$ is $a\log x+b$ for constants $a,b$;
\item[(iii)] the weights of the carrier-scaled and the unscaled generator differ by the
gauge constant, $\log(\Delta x_n)=\log x_n+\log\Delta$, which is
$H_\Delta=H+\log\Delta\cdot I$ (Theorem~\ref{thm:logweight}(iii)).
\end{enumerate}
\end{theorem}
\leanref{CarrierReadout.log_isGaugeWeight, CarrierReadout.gaugeWeight_eq_log,
CarrierReadout.logWeight_gauge_shift}

\begin{proof}
(i) is $\log(\Delta x)=\log x+\log\Delta$. (ii): $\varphi(u)=w(e^{u})-w(1)$ satisfies
$\varphi(u+v)=\varphi(u)+\varphi(v)$, by the gauge law at $\Delta=e^{u}$, $x=e^{v}$ and
again at $x=1$; a continuous additive map $\mathbb R\to\mathbb R$ is linear, so
$\varphi(u)=au$ and $w(x)=a\log x+w(1)$. (iii) is (i) at the sites.
\end{proof}

\noindent
So the site weight of the fiber is not chosen. The design fact that the carrier scale is
a gauge --- Proposition~\ref{prop:zeros}, measured in \S\ref{sec:validation} --- \emph{is}
the statement that the weight is logarithmic, by (ii); the affine freedom $a\log x+b$ is
the unit of the ordinate ($a$) and the carrier scale ($b=\log\Delta$).

\begin{proposition}[the geometric head]\label{prop:geohead}
For $r>0$ and $k\ge0$, $S(k)\ge\pi rk^{2}$; hence at $\Delta=\pi/3$ the carrier index
satisfies, for $p,r>0$,
\[
  N(y)\;\ge\;\frac{3r}{p^{2}}\,e^{2y}:
\]
the number of sites the carrier has laid down by the ordinate $y$ grows exponentially in
$y$, against the linear $2y/\pi$ that the lane saddle of \S\ref{sec:extractor} requires.
\end{proposition}
\leanref{CarrierReadout.arclength_ge_pi_r_sq, CarrierReadout.Nindex_ge_exp}

\begin{proof}
In the closed form of Theorem~\ref{thm:arclength} the root is at least $2\pi rk$ and the
$\operatorname{arsinh}$ term is non-negative, so $S(k)\ge\tfrac k2\cdot2\pi rk$; then
$k(y)=e^{y}/p$.
\end{proof}

That $\Delta$ is a gauge and not a parameter of the answer is the content of
Proposition~\ref{prop:zeros} below, and it is measured in \S\ref{sec:validation}: the
extracted ordinates are identical across a factor of $57$ in carrier scale.

%======================================================================
\section{The fiber: the $L$-function riding the carrier}\label{sec:fiber}

The carrier is a ruler; the \emph{fiber} is what rides it. The $L$-function's
representation lives on the carrier, and the phasors live on the fiber. As the fiber
grows from the origin outward, the conductor assigns a phasor to each carrier integer
consumed --- positive, negative, or neutral channels, supplied by the character. For
$\zeta$ the channels are the even and odd sites and the assignment is the
non-principal character mod $2$; that is where the two \emph{lanes} come from, and
they are not a device introduced for $\zeta$ but what a conductor of $2$ supplies.

Three properties of the fiber are used later and are worth stating as the geometry
gives them, not as the series gives them.

\emph{Phasors are born at zero magnitude and grow.} A phasor enters continuously at
magnitude $0$ at its own site and grows until finished. There is no site at the origin
and no phasor before the first, and in particular there is no convergence gate: the
fiber is defined at every height. (The abscissa of convergence is a fact about a
projected series, \S\ref{sec:descent}, and it is not a constraint on the fiber.)

\emph{Phasors spin in the logarithm because the carrier turns.} The site at arclength
$x_n$ carries the phasor $x_n^{-s}$; at $s=\tfrac12+it$ this has modulus $x_n^{-1/2}$
and phase $-t\log x_n$. The logarithm is not put in by hand, and \S\ref{sec:carriergeom}
has already said why once: Theorem~\ref{thm:gauge} forces it. The climber of
Definition~\ref{def:climber} says it a second time, as a registration. On the axial rail
--- $r=0$, where the carrier's arclength is its height (Theorem~\ref{thm:arclength}) ---
the climber reaches arclength $x$ at the ordinate $\log x$; so it reaches the site
$x_n=\Delta n$ exactly at the ordinate $\log x_n$, which is the site's weight, and at
$\Delta=\pi/3$ the carrier index at that ordinate is $N(\log x_n)=n$. The phase
$-t\log x_n$ is therefore the product of two registrations by the same climber: the
ordinate $t=\log z$ at which the fiber is read, and the ordinate $\log x_n$ at which the
site was reached. (On a rail of radial rate $r>0$ the same site is reached at the
ordinate $\log(p\,S^{-1}(x_n))$; the weight $\log x_n$ is the axial reading, and the
site's \emph{angle} --- its cell --- is read on the turning rail,
Proposition~\ref{prop:cellangle}.)
\leanref{CarrierReadout.climber_reaches_arclength,
CarrierReadout.climber_reaches_site_at_weight, CarrierReadout.Nindex_at_site,
CarrierReadout.mellinSpin_eq_registrations, cpow_critical_line,
eta_term_critical_line, norm_cpow_critical_line}

\emph{The warp.} A \emph{warp} is a scaling by a function: dynamic, unit modulus,
readout-preserving. It is the geometric knob --- it moves the fiber through
configuration space, and turning it changes \emph{the height at which the cancellation
completes}. The warp used here aligns the fiber's vanishing with the harmonic cell
edges, and the vanishing books no residue. There is no such thing as an inverse warp:
the reciprocal of a warp is again a warp.

\begin{theorem}[antipodal cancellation, logarithm-free]\label{thm:antipodal}
In the $\mu_6$ cell the antipode of a cell is its negative:
$\operatorname{cell}(n)+\operatorname{cell}(n+3)=0$.
\end{theorem}
\leanref{Readout.cell_antipodal_cancel}

\noindent
This is the mechanism in its native form: exact cancellation of the positive against
the negative side, forced by the $\pi/3$ helix angle alone, with no logarithm
anywhere. The vanishing of the fiber is a \emph{focal cancellation} --- the channels
point at a common height and cancel there --- and it is an event at a height, not a
root of a formula.

\begin{proposition}[the cell is the helix's angle]\label{prop:cellangle}
$\operatorname{cell}(n)=e^{\,i\cdot2\pi\cdot n/6}$: the $\mu_6$ cell of the site $n$ is the
angle of the helix at turn parameter $k=n/6$ --- six sites per turn --- and the helix at
that parameter is
\[
  \gamma(n/6)\;=\;\Bigl(r\tfrac n6\operatorname{Re}\operatorname{cell}(n),\;
                       r\tfrac n6\operatorname{Im}\operatorname{cell}(n),\;p\tfrac n6\Bigr).
\]
The antipode $n+3$ of Theorem~\ref{thm:antipodal} is half a turn up.
\end{proposition}
\leanref{CarrierReadout.cell_eq_helix_angle, CarrierReadout.helix_at_sixth_turn}

\begin{definition}[fiber phasor, its $\pm$ block, and a cancelling partner]\label{def:fiberphasor}
A \emph{fiber phasor} is a unimodular complex number $w$, $\lvert w\rvert=1$: the state
of one site of the fiber at a height. Its \emph{$\pm$ block} is the conjugate-pair
matrix $\operatorname{diag}(w,\bar w)$ --- the phasor read from the two ends of the
geometry, the end-swap being conjugation (Lemma~\ref{lem:crit}) --- and its two entries
are the phasor's \emph{eigenheights}; they \emph{phase-align} when $w=\bar w$. A
\emph{cancelling partner} of $w$ is an $a\in\mathbb C$ with $w+a=0$.
\end{definition}
\leanref{DeBranges.conjPairBlock}

\begin{theorem}[the fiber always leaves a residue]\label{thm:residue}
A unimodular phasor has a unique cancelling partner, and the $\pm$ block of a
unimodular fiber phasor has determinant $1$; so the two sides cannot both vanish, and
one cannot simultaneously phase-align two eigenheights and cancel exactly.
\end{theorem}
\leanref{Readout.fiber_residue_no_double_cancellation,
Readout.cancelling_partner_unique, Readout.fiber_pair_block_det_one}

\begin{definition}[the $\pm$ channels and the closure function]\label{def:channels}
For a real-valued character $\chi$ and an ordinate $t$ let
$a_n(t)=\tfrac{\pi}{3}\,n^{-1/2}e^{-it\log n}$ be the amplitude carried by the site $n$
--- the same for both channels, the character entering only through the split --- and
put
\[
  P^{(N)}_\chi(t)=\sum_{\substack{n\le N\\ \chi(n)=+1}}a_n(t),\qquad
  M^{(N)}_\chi(t)=\sum_{\substack{n\le N\\ \chi(n)=-1}}a_n(t),
\]
the \emph{positive} and \emph{negative channels} truncated at $N$. The \emph{closure
function} at height $Z>0$ is the $L$-value at the corresponding ordinate,
$\Phi_\chi(Z)=L(\chi,\tfrac12+i\log Z)$.
\end{definition}
\leanref{Focal.aTerm, Focal.Pchan, Focal.Mchan, Geometric.Phi}

\begin{theorem}[focal cancellation is the vanishing]\label{thm:focal}
For a non-principal real-valued $\chi$ and $Z>0$, the positive-minus-negative channel
closes at height $Z$ exactly when the $L$-value vanishes at the corresponding
ordinate:
\[
  P_\chi^{(N)}(\log Z)-M_\chi^{(N)}(\log Z)\;\xrightarrow[N\to\infty]{}\;0
  \quad\Longleftrightarrow\quad
  \Phi_\chi(Z)=0 .
\]
\end{theorem}
\leanref{Readout.focal_pm_cancellation_iff_L_zero}

\noindent
This is the bridge in its native form, and it is an \emph{iff}: the geometric event ---
the $\pm$ channels closing at a height --- and the analytic event --- the $L$-function
vanishing --- are the same occurrence seen in two charts. In the chart,
Definition~\ref{def:channels} makes $P^{(N)}_\chi-M^{(N)}_\chi$ the partial sum
$\tfrac{\pi}{3}\sum_{n\le N}\chi(n)n^{-1/2-i\log Z}$ of the $L$-series at
$\tfrac12+i\log Z$, and the theorem's chart reading is the conditional convergence of that
series on the line. The object's reading is the one displayed: the conductor sorts the
sites into two channels, and the fiber vanishes at a height exactly when they close
there.

\begin{remark}[$\zeta$ is not covered by Theorem~\ref{thm:focal}, and what covers it
instead]\label{rem:zetabridge}
The hypothesis $\chi\ne1$ excludes the trivial character, so
Theorem~\ref{thm:focal} does \emph{not} apply to $\zeta$. This is not an oversight of
the statement: the $\pm$ split there is by character \emph{value}, and for $\zeta$ that
split degenerates --- there is no non-trivial sign to sort by. Saying otherwise would be
the characteristic error of this subject, and we do not.

For $\zeta$ the corresponding decomposition is by \emph{angle} rather than by value:
the bank's phasors, with von Mangoldt weights and no character hypothesis, are sorted
into the six $\mu_6$ sectors by $t\log n$, and antipodal sectors are compared. The
decomposition is an exact partition, the antipodal pairing is a fixed-point-free
involution, and the imbalance vanishes at a truncation exactly when all three antipodal
pairs agree --- all machine-checked.
\leanref{ZetaSectorBank.sector_partition, ZetaSectorBank.bank_pair_decomposition,
ZetaSectorBank.antipode_involutive, ZetaSectorBank.imbalance_eq_zero_iff,
ZetaSectorBank.norm_zetaTerm}

The balance law built on top of that bookkeeping --- that the sectors balance in the
limit at every on-line zero --- was carried in an earlier version of this work as a
named hypothesis. It is not open. It is false, and the reason is elementary.
\end{remark}

\begin{theorem}[the antipodal arc obstruction]\label{thm:arcobstruction}
For every real $t$ and every $N$,
\[
  \frac{\sqrt3}{2}\sum_{n\le N}\frac{\Lambda(n)}{\sqrt n}\;\le\;\operatorname{imb}(N,t),
  \qquad
  \operatorname{imb}(N,t)=\sum_{j=0}^{2}\bigl\lvert S_j(N,t)-S_{j+3}(N,t)\bigr\rvert .
\]
Consequently $\operatorname{imb}(N,t)\not\to0$ for \emph{every} real $t$: the balance law
fails at every on-line zero, and since $\zeta$ has on-line zeros
($\gamma_1=14.134\ldots$, certified in \S\ref{sec:bandverify}) the law is false. The
compiled form of that last step is an equivalence, the law holding if and only if $\zeta$
has no zero on the critical line at all.
\end{theorem}
\leanref{ZetaSectorBank.imbalance_ge_half_mass, ZetaSectorBank.not_sectorBalanceAt,
ZetaSectorBank.angleSectorBalanceLaw_iff_no_critical_zero,
ZetaSectorBank.re_rotate_ge, ZetaSectorBank.re_rotate_le, ZetaSectorBank.norm_sub_ge,
ZetaSectorBank.floorConst_le_imbalance}

\begin{proof}
A sector is an arc of angular width $2\pi/6=\pi/3$. Every phasor in sector $j$
therefore lies within $\pi/6$ of that sector's midpoint direction $u_j$, so its
component along $u_j$ is at least $\cos(\pi/6)=\sqrt3/2$ times its modulus. The
antipodal sector $j+3$ is the same arc rotated by $\pi$, so its phasors have midpoint
direction $-u_j$ and their components along $u_j$ are at most $-\sqrt3/2$ times their
moduli. Subtracting therefore \emph{adds} the two masses:
\[
  \lvert S_j-S_{j+3}\rvert\;\ge\;\bigl\langle S_j-S_{j+3},\,u_j\bigr\rangle
  \;\ge\;\frac{\sqrt3}{2}\bigl(M_j+M_{j+3}\bigr),
\]
where $M_j$ is the mass in sector $j$; and the three antipodal pairs sweep the six
sectors exactly once, giving the displayed bound. For the consequence, no asymptotics
are needed: the $n=2$ term alone gives
$\operatorname{imb}(N,t)\ge\tfrac{\sqrt3}{2}\cdot\tfrac{\log2}{\sqrt2}>0.42$ for every
$N\ge2$, and a sequence bounded below by a positive constant does not tend to $0$.
\end{proof}

\noindent
Measured, for confirmation rather than for evidence: at $\gamma_1=14.1347\ldots$ and at
the non-zero $t=20$, the ratio $\operatorname{imb}/\text{mass}$ sits at $0.955$ for
every $N$ from $10^{3}$ to $4\times10^{6}$ --- above the proved floor $\sqrt3/2=0.866$,
and \emph{identical at the zero and the non-zero to three decimals}. So the statistic
does not merely fail to vanish; it does not distinguish zeros from non-zeros at all.
Two independent implementations agree --- one in \lean{numpy}
with a sieve, one in \texttt{mpmath} at $40$ digits with trial division and no
\lean{numpy} at all (\lean{checks/balance_law_refuted.py},
\lean{checks/balance_law_independent.py}); they match to six digits.

\begin{remark}
This is reported because it removes the one place where this work carried an unproved
hypothesis, and it removes it in the direction that is easy to get wrong. The correct
conclusion is not that the $\mu_6$ geometry is idle --- the log-free antipodal
cancellation of Theorem~\ref{thm:antipodal} is exact and is the mechanism --- but that
\emph{this particular statistic}, sorting the von Mangoldt bank into six arcs and
differencing antipodes, cannot detect a zero, because differencing an arc against its
own antipode is a sum and not a cancellation. Nothing elsewhere in the paper used it.
\end{remark}

\noindent
So the licence for the chart work on $\zeta$ is not Theorem~\ref{thm:focal} and not the
balance law. It is an exact algebraic identity, and that is the stronger footing:
the fiber \emph{factors} on the line as
$F_\Delta(\tfrac12+it)=w_\Delta(t)\,\Lambda(\tfrac12+it)$ with $w_\Delta(t)\ne0$
(Theorem~\ref{thm:weld}), so the readout's vanishing and $\zeta$'s vanishing are the
same set --- not in a limit, not conditionally, but by a factorisation. Everything from
\S\ref{sec:carrier} onward rests on that identity. Theorem~\ref{thm:focal} and
Remark~\ref{rem:zetabridge} say what the same phenomenon looks like in its geometric
form, and how much of that form is currently proved.

%======================================================================
\section{The descent to the chart}\label{sec:descent}

The one-dimensional $L$-function is not the object; it is the object read in a chart,
and the chart is reached by two projections, each with its own loss.

\begin{center}
\begin{tabular}{@{}llll@{}}
\toprule
step & map & drops & survives\\
\midrule
$3\text{-D}\to 2\text{-D}$ & Cayley/M\"obius & the radius & angle, height\\
$2\text{-D}\to 1\text{-D}$ & $y=\log z$ & the angle & the height\\
\bottomrule
\end{tabular}
\end{center}

\begin{theorem}[the two steps]\label{thm:descent}
The Cayley map $w\mapsto(w-i)/(w+i)$ carries the real readout line to the unit circle,
$\lVert\operatorname{cayley}(x)\rVert=1$ for real $x$; and the height readout
$y=\log z$ inverts globally, $\log(e^{\gamma})=\gamma$, with no branch cut.
\end{theorem}
\leanref{Readout.cayley, Readout.cayley_real_norm_one, Readout.heightReadout,
Readout.readout_exp, Readout.height_eq_exp_of_readout}

\noindent
Two consequences are used repeatedly and are stated here so that they are not
rediscovered later as difficulties.

\emph{The classical strip is the Cayley image of the unit circle.} It is a feature of
the chart, not of the geometry, and phenomena that live in it are chart artifacts. The
geometry has no free abscissa to vary: there is a height axis, and the critical
abscissa is not posited there but is where the double-ended reflection has its fixed
set. In particular the abscissa of convergence $\operatorname{Re}s>1$ is where a
projected series stops converging; the fiber has no convergence gate, and this paper
never treats the strip's boundary as a constraint on the object.

\emph{The height survives; the ordinate is its registration.} The logarithm appears
exactly once in the whole model, in the $2\text{-D}\to1\text{-D}$ step, and that is the
only place it appears. A vanishing event happens at a height $z$; its readout is the
ordinate $\gamma=\log z$. The two coincide; they are not the same quantity.

\emph{The functional equation is the double end.} The geometry has two ends, mirror
images through the origin, and a readout cannot tell which end it is reading. That
end-swap descends through the Cayley step to $s\mapsto1-s$. The functional equation is
therefore not an identity to be derived in the chart but a symmetry the object simply
has; the chart's devices for reconstructing it from inside the projection are chart
devices, and the paper does not use them to explain carrier facts.

%======================================================================
\section{The chart readout: lanes and the weld gauge}\label{sec:carrier}

From here to \S\ref{sec:extractor} the work is in the chart, and every object is a
readout of \S\S\ref{sec:geometry}--\ref{sec:fiber}. The reason for descending is
Theorem~\ref{thm:focal}: the events are the same, and in the chart they are
computable and machine-checkable.

\begin{definition}[carrier fiber]\label{def:fiber}
For $\Delta>0$ set
\[
  F_\Delta(s)\;=\;\Delta^{-s}\bigl(1-2^{1-s}\bigr)\zeta(s),
\]
the sum over the sites $x_n=\Delta n$ of the phasors $x_n^{-s}$ with the odd and even
sites entering with opposite sign.
\end{definition}
\leanref{carrierFiber}

\begin{theorem}[the fiber is the function]\label{thm:fiberisfunction}
Let $\Delta>0$; write $B_\Delta(s)=(x_n^{-s})_{n\ge1}$ for the bank on the carrier of scale
$\Delta$, $x_n=\Delta n$, and $\ell$ for the lane-balanced chart, the phasors summed in
pairs in the order they enter.
\begin{enumerate}
\item[(i)] $B_\Delta(s)=e^{-sH_\Delta}\mathbb 1$: the phasor at the site $x_n$ is
$\exp(-s\,\lambda_n(H_\Delta))$ with $\lambda_n(H_\Delta)=\log x_n$ the eigenvalue of the
scaled generator (\S\ref{sec:hamiltonian}); and $B_\Delta(s)=\Delta^{-s}B_1(s)$, the
carrier scale a gauge on the bank.
\item[(ii)] For $\operatorname{Re}s>0$, $s\ne1$: $\ell(B_\Delta(s))=F_\Delta(s)$. The chart
of the warped bank on the scaled carrier \emph{is} the fiber ---
$\Delta^{-s}(1-2^{1-s})\zeta(s)$, the analytic function and not a truncation of it --- on
the whole channel on which the ordered sum converges.
\item[(iii)] In the strip, $\zeta(s)=0$ if and only if $B_\Delta(s)\in\ker\ell$, at every
carrier scale.
\item[(iv)] The same identity holds on the three-dimensional phasors: with $\chi_2$ the
conductor-$2$ lane character, the planar projections of the $3$-D phasors
$\bigl(\chi_2(n)\,n^{-\sigma}e^{-iy\log n},\;(1-\lvert\chi_2(n)\rvert)\,n^{-\sigma}\bigr)$,
summed in pairs in the order they enter and scaled by $\Delta^{-s}$, equal
$F_\Delta(\sigma+iy)$ for $\sigma>0$, $\sigma+iy\ne1$. (The mass-axis coordinate vanishes:
for $\zeta$ the neutral channel is empty.)
\item[(v)] The one-dimensional function has no content beyond the chart: any entire
function that agrees with $\ell(B_1(s))$ on the channel equals $(1-2^{1-s})\zeta(s)$ at
every $s\ne1$.
\end{enumerate}
\end{theorem}
\leanref{CarrierReadout.fiber_eq_function, CarrierReadout.fiber3D_eq_function,
CarrierReadout.function_determined_by_fiber, CarrierReadout.plane_phasor3D_eq,
CarrierReadout.laneChart3D_eq, CarrierReadout.bank_eq_rpow_mul_spin, DVP.etaLim_eq_zeta}

\begin{proof}[Proof sketch]
(i): $x_n^{-s}=e^{-s\log x_n}$, and $(\Delta n)^{-s}=\Delta^{-s}n^{-s}$
(\lean{CarrierReadout.bank}$\Delta$\lean{_eq_exp},
\lean{CarrierReadout.bank}$\Delta$\lean{_eq_scale}). (ii): pairing consecutive phasors
gives the ordered sum $\sum_k\bigl((2k+1)^{-s}-(2k+2)^{-s}\bigr)$, which converges for
$\operatorname{Re}s>0$ and equals $(1-2^{1-s})\zeta(s)$ there --- the continuation of the
lane-balanced series across the abscissa of the unbalanced one, proved in the sources from
$\operatorname{Re}s>1$ by the identity theorem (\lean{DVP.etaLim_eq_zeta}); then the gauge
of (i) (\lean{CarrierReadout.laneChart_bank}$\Delta$). (iii): the lane factor does not
vanish in the strip (\lean{CarrierReadout.zeta_zero_iff_bank}$\Delta$\lean{_mem_ker}).
(iv): the planar projection of the $3$-D phasor at site $n$ is
$\chi_2(n)\,n^{-\sigma}e^{-iy\log n}=\chi_2(n)\,n^{-s}$. (v): both sides are analytic on
$\mathbb C\setminus\{1\}$, which is connected, and agree on the open channel.
\end{proof}

\noindent
Before this theorem the closed form of Definition~\ref{def:fiber} was, in the sources, the
\emph{definition} of \lean{carrierFiber}; the description ``the sum over the sites of the
phasors, the odd and even sites entering with opposite sign'' is now what is proved, on
the three-dimensional phasors and at every carrier scale. The boundary is stated exactly:
the ordered sum converges on $\operatorname{Re}s>0$ and nowhere to the left of it, and
there the function is determined by the chart, which is (v).

\noindent
The lane-balanced sum converges for $\operatorname{Re}s>0$ --- that is the whole point
of taking the lanes with opposite sign, and it is why the construction can be written
\emph{on} the critical line rather than continued to it. The unbalanced Dirichlet
series would need $\operatorname{Re}s>1$; the alternating one does not, and the
displayed closed form is an identity of convergent objects at $s=\tfrac12+it$. The
convergence is conditional, which is exactly the fact the head law of
\S\ref{sec:extractor} has to pay for.

\begin{definition}[weld factor]\label{def:weld}
\[
  w_\Delta(t)\;=\;\frac{\Delta^{-(\frac12+it)}\bigl(1-2^{\frac12-it}\bigr)}{\GR(\tfrac12+it)} ,
  \qquad \GR(s)=\pi^{-s/2}\Gamma(s/2).
\]
\end{definition}
\leanref{weldFactor}

\begin{theorem}[the weld gauge]\label{thm:weld}
For every $\Delta>0$ and every $t\in\R$,
\[
  e^{-i\arg w_\Delta(t)}\,F_\Delta(\tfrac12+it)\;=\;\lvert w_\Delta(t)\rvert\,
  \Lambda(\tfrac12+it)\;\in\;\R ,
\]
and $\lvert w_\Delta(t)\rvert=\Delta^{-1/2}\lvert 1-2^{\frac12-it}\rvert/\lvert\GR(\tfrac12+it)\rvert>0$.
\end{theorem}
\leanref{weld_gauge_exact, weld_gauge_im_zero, weldFactor_norm,
carrierFiber_eq_weldFactor_mul, lane_factor_ne_zero, weldFactor_ne_zero}

\begin{proof}[Proof sketch]
The factorisation $F_\Delta(\tfrac12+it)=w_\Delta(t)\Lambda(\tfrac12+it)$ is
$\eta=(1-2^{1-s})\zeta$ together with $\Lambda=\GR\zeta$, the $\Delta^{-s}$ cancelling
between fiber and gauge. Rotating a nonzero complex number by its own argument lands it
on its modulus, so the left side is $\lvert w_\Delta(t)\rvert\,\Lambda(\tfrac12+it)$;
and $\Lambda(\tfrac12+it)\in\R$ by the functional equation together with Schwarz
reflection (\S\ref{sec:formalization}). Positivity of $\lvert w_\Delta(t)\rvert$ is the
observation that $1-2^{\frac12-it}$ never vanishes, since
$\lvert2^{\frac12-it}\rvert=\sqrt2\ne1$; and that $\GR(\tfrac12+it)$ is finite and
nonzero, since $\Gamma$ has no zeros and $\tfrac14+\tfrac{it}{2}$ is never a pole.
\end{proof}

\begin{remark}\label{rem:theta}
No asymptotics enter. The Riemann--Siegel theta function never appears; its role is
played by $\arg\GR(\tfrac12+it)$, which it equals. Indeed
$\GR(s)=\pi^{-s/2}\Gamma(s/2)$ gives
\[
  \arg\GR(\tfrac12+it)\;=\;\arg\Gamma(\tfrac14+\tfrac{it}{2})-\tfrac{t}{2}\log\pi
  \;=\;\vartheta(t),
\]
which is the definition of $\vartheta$. The equality is one of \emph{arguments}, hence
modulo $2\pi$: $\arg$ is the principal value and $\vartheta$ the continuous branch, and
the two differ by the winding $2\pi k(t)$. Nothing tracks that branch, because
$w_\Delta$ uses the argument only through $e^{-i\arg\GR}$, on which $k$ acts trivially.
This is the whole reason the construction needs no asymptotic expansion of $\vartheta$:
the classical device computes $\vartheta$ to subtract it, and here it is divided out
exactly by the local factor it came from.

The script \lean{extractor/checks/theta_identity.py} checks this numerically at
$60$-digit working precision, at thirteen sample heights spanning
$t\in[0,10^{4}]$, including $t=2\pi$ and the first ordinate: worst
branch-corrected discrepancy $2.5\times10^{-57}$, i.e.\ agreement to working
precision at every sample. The check was pre-registered to refute the remark at any
sample exceeding $10^{-30}$.
\end{remark}

\begin{definition}[the readout]\label{def:readout}
$u_\Delta(t)=\lvert w_\Delta(t)\rvert\,\operatorname{Re}\Lambda(\tfrac12+it)$.
\end{definition}
\leanref{weldReadout, xiRe}

\begin{definition}[event]\label{def:event}
An \emph{event} is a real $t$ at which the readout vanishes --- equivalently, by
Proposition~\ref{prop:zeros} below, a zero of $\zeta$ on the critical line. Its
\emph{multiplicity} is the order of vanishing there, and an event of multiplicity $1$
is called \emph{simple}. Counts of events are always counts with multiplicity unless
said otherwise. The word is used rather than ``zero'' as a standing reminder that an
event is an \emph{on-line} zero: the readout sees exactly these, and whether they
exhaust the zeros of $\zeta$ in a band is the question \S\ref{sec:pairing} keeps open
and \S\ref{sec:notclaimed}(iv) restates.
\end{definition}

\begin{proposition}[the readout is Hardy's $Z$ in the carrier's gauge]\label{prop:isZ}
For every $\Delta>0$ and every real $t$,
\[
  u_\Delta(t)\;=\;\lvert w_\Delta(t)\rvert\,\bigl\lvert\GR(\tfrac12+it)\bigr\rvert\;Z(t),
  \qquad Z(t)=e^{i\vartheta(t)}\zeta(\tfrac12+it),
\]
and both factors in front are strictly positive. Consequently $u_\Delta$ and $Z$ have the
same zeros, the same orders and the same sign at every $t$.
\end{proposition}

\begin{proof}
$\Lambda(\tfrac12+it)=\GR(\tfrac12+it)\zeta(\tfrac12+it)$ and
$\vartheta=\arg\GR(\tfrac12+it)$ (Remark~\ref{rem:theta}), so
$\Lambda(\tfrac12+it)=\lvert\GR(\tfrac12+it)\rvert\,e^{i\vartheta(t)}\zeta(\tfrac12+it)
=\lvert\GR(\tfrac12+it)\rvert\,Z(t)$. Multiply by $\lvert w_\Delta(t)\rvert$ and apply
Definition~\ref{def:readout}. Positivity of the two factors is
Theorem~\ref{thm:weld}.
\end{proof}

\begin{remark}[what this costs and what it buys]\label{rem:isZ}
It should be said plainly, because it governs how the rest of the paper is to be read.
\emph{The carrier readout is Hardy's $Z$, up to a strictly positive factor.} It follows
that \S\S\ref{sec:ordinate}--\ref{sec:capstone} --- reality on the line, the sign change
across an odd-order zero, uniqueness and simplicity by Rolle and the intermediate value
theorem --- are the classical $Z$/$\Xi$ isolation argument. We do not claim them as new
mathematics, and a reader who already knows that argument will recognise every step.

What the identification buys is the other direction, and it is used twice. First, the
certification is \emph{evaluator-agnostic}: every theorem here consumes $u_\Delta$ only
through its sign, so the sign may be taken from any evaluator of $Z$ --- which is what
closes (O4), and it is why a two-ended head can be substituted for the lane bank without
touching a proof. Second, it is what makes the geometry's contribution locatable rather
than decorative: since the isolation argument is classical, the question ``what does the
carrier actually do here?'' has a definite answer, and \S\ref{sec:loadbearing} gives it.
\end{remark}

\begin{proposition}[events]\label{prop:zeros}
$u_\Delta(t)=0 \iff \zeta(\tfrac12+it)=0$, for every $\Delta>0$; and the sign of
$u_\Delta$ does not depend on $\Delta$.
\end{proposition}
\leanref{weldReadout_eq_zero_iff, weldReadout_carrier_ratio, weldReadout_sign_carrier_free}

\begin{proof}[Proof sketch]
By Theorem~\ref{thm:weld} the readout is $\lvert w_\Delta(t)\rvert\Lambda(\tfrac12+it)$
with $\lvert w_\Delta(t)\rvert>0$, so it vanishes exactly where $\Lambda(\tfrac12+it)$
does; and $\Lambda=\GR\zeta$ with $\GR(\tfrac12+it)$ finite and nonzero, so that is
exactly where $\zeta(\tfrac12+it)$ does. Changing $\Delta$ multiplies the readout by
the positive ratio $\lvert w_\Delta(t)\rvert/\lvert w_{\Delta'}(t)\rvert
=(\Delta'/\Delta)^{1/2}$, which cannot change a sign.
\end{proof}

%======================================================================
\section{The ordinate variable and the entire readout}\label{sec:ordinate}

Everything so far has been written on the line $s=\tfrac12+it$, where the readout is
a function of the real variable $t$. For the counting arguments we need an
\emph{entire} function of a complex variable whose restriction to the real axis is
that readout, so that order, multiplicity and the argument principle are available.
The substitution that does this is $s=\tfrac12+iz$: the functional equation's
involution $s\mapsto1-s$ becomes $z\mapsto-z$, and complex conjugation on the line
becomes conjugation in $z$. This is Riemann's own change of variable~\cite{Riemann1859}, not
something the carrier supplies, and we record it as his.

Write $\Lambda_0=\lean{completedRiemannZeta}_0$ for the entire completed zeta, so that
$\Lambda(s)=\Lambda_0(s)-\frac1s-\frac1{1-s}$.

\begin{definition}[the entire readout]\label{def:xi}
\[
  \Xi(z)\;=\;\bigl(\tfrac12+iz\bigr)\bigl(\tfrac12-iz\bigr)\,
             \Lambda_0\bigl(\tfrac12+iz\bigr)\;-\;1 .
\]
\end{definition}

\leanref{xiEntire, xiEntire_differentiable}

\begin{remark}[normalisation]\label{rem:normalisation}
This $\Xi$ is not Riemann's. With $\xi(s)=\tfrac12 s(s-1)\Lambda(s)$ one has
$s(1-s)\Lambda(s)=-2\xi(s)$, so
\[
  \Xi(z)\;=\;-2\,\xi(\tfrac12+iz),
\]
an overall factor of $-2$ against the classical normalisation. Nothing below depends on
the constant --- every statement is about zeros, orders and sign \emph{changes} --- but
a reader comparing formulas against Riemann's or Titchmarsh's should carry it.
\end{remark}

\noindent
Entirety is manifest: $\Lambda_0$ is entire and the remaining factors are polynomials.
No pole is cancelled by hand and no removable singularity is argued away --- the
$-1$ is exactly what the two poles contribute, as the next statement records.

\begin{proposition}[pole clearing]\label{prop:xiclear}
For $z$ with $\tfrac12+iz\notin\{0,1\}$,
$\Xi(z)=(\tfrac12+iz)(\tfrac12-iz)\,\Lambda(\tfrac12+iz)$; that is, $\Xi$ is
$s(1-s)\Lambda(s)$ in the ordinate variable. On the real axis,
\[
  \Xi(t)\;=\;\bigl(\tfrac14+t^{2}\bigr)\,\Lambda\bigl(\tfrac12+it\bigr),
  \qquad t\in\R ,
\]
and $\tfrac14+t^{2}>0$, so $\Xi\restriction\R$ is a strictly positive multiple of the
readout: same sign, same zeros, same orders.
\end{proposition}
\leanref{xiEntire_eq_completed, xiEntire_ofReal, quarter_add_sq_pos}

\begin{proof}[Proof sketch]
Write $A=\tfrac12+iz$ and $B=\tfrac12-iz$, so that $A+B=1$. With
$\Lambda=\Lambda_0-\tfrac1A-\tfrac1B$ at $s=A$, the clearing is the identity
\[
  AB\Bigl(L-\tfrac1A-\tfrac1B\Bigr)\;=\;ABL-1 ,
\]
valid for any $L$ whenever $A,B\ne0$ and $A+B=1$: expand the left side to
$ABL-B-A$ and substitute $A+B=1$. The two factors are never inspected --- only their
sum is used. On the real axis $AB=(\tfrac12+it)(\tfrac12-it)=\tfrac14+t^{2}$, which is
positive, giving the displayed form.
\end{proof}

\begin{proposition}[reality]\label{prop:xireal}
$\Xi(t)\in\R$ for every real $t$, and $\Xi(t)=0\iff\zeta(\tfrac12+it)=0$.
\end{proposition}
\leanref{xiEntire_im_ofReal, xiEntire_ofReal_eq_zero_iff, xiEntire_re_eq_zero_iff}

\begin{proof}[Proof sketch]
Immediate from Proposition~\ref{prop:xiclear}: $\Xi(t)$ is the real number
$\tfrac14+t^{2}$ times $\Lambda(\tfrac12+it)$, which is real, and the multiplier is
positive, so $\Xi(t)=0$ exactly when $\Lambda(\tfrac12+it)=0$, i.e.\ exactly when
$\zeta(\tfrac12+it)=0$.
\end{proof}

\noindent
So $\Xi$ is a real-on-$\R$ entire function with the on-line zeros of $\zeta$ as its
real zeros, carrying the same multiplicities. Every count below is a count of
zeros of $\Xi$.

\begin{theorem}[multiplicity sign law]\label{thm:signflip}
Let $G$ be analytic at a real point $t_0$, real on the real axis, not identically
zero near $t_0$, of order $m$ there. Then for all sufficiently small $h>0$,
\[
  (-1)^m\,\operatorname{Re}G(t_0-h)\,\operatorname{Re}G(t_0+h)\;>\;0 .
\]
Consequently the carrier readout's sign is multiplied by $(-1)^m$ across an event of
multiplicity $m$, at every carrier scale.
\end{theorem}
\leanref{sign_flip_of_order, weldReadout_sign_flip}

\begin{proof}[Proof sketch]
Factor $G(z)=(z-t_0)^m g(z)$ with $g$ analytic and $g(t_0)\ne0$. Along the real axis
$(t-t_0)^m$ is a nonzero real for $t\ne t_0$, so $g$ inherits reality there by a
punctured limit, and $g$ keeps the sign of $g(t_0)$ on a small enough interval. Then
\[
  \operatorname{Re}G(t_0-h)\,\operatorname{Re}G(t_0+h)
  \;=\;(-h)^m h^m\,g(t_0-h)g(t_0+h)
  \;=\;(-1)^m h^{2m}\,g(t_0-h)g(t_0+h),
\]
and the last product is positive for small $h$. Multiplying by $(-1)^m$ gives the
claim. The carrier statement follows because the readout is a positive multiple of
$\Xi\restriction\R$ (Remark~\ref{rem:whichderiv}).
\end{proof}

\begin{remark}
The statement is with multiplicity. This is what removes any simplicity hypothesis
from the chain; it is also the honest price of a sign bit, which cannot see an event
of even order.
\end{remark}

%======================================================================
\section{Off-line pairing and the strip count}\label{sec:pairing}

\begin{lemma}[the reflection is conjugation]\label{lem:crit}
In the ordinate variable the involution $s\mapsto 1-\bar s$, whose fixed set is the
critical line, is $z\mapsto\bar z$.
\end{lemma}
\leanref{crit_is_conj, ordinate_eq, offCriticalLine_iff}

\begin{proof}[Proof sketch]
With $s=\tfrac12+iz$ one has $\bar s=\tfrac12+\overline{iz}=\tfrac12-i\bar z$, hence
\[
  1-\bar s\;=\;1-\tfrac12+i\bar z\;=\;\tfrac12+i\bar z ,
\]
which is the point of ordinate $\bar z$. So $s\mapsto1-\bar s$ reads as
$z\mapsto\bar z$, whose fixed set is $\operatorname{Im}z=0$, i.e.\
$\operatorname{Re}s=\tfrac12$.
\end{proof}

\begin{theorem}[off-line pairing]\label{thm:pairing}
A zero of $\Xi$ off the real axis --- a zero of $\zeta$ off the critical line, which is
not an event in the sense of Definition~\ref{def:event} --- has a distinct partner at the
same ordinate with the same multiplicity.
\end{theorem}
\leanref{xiEntire_conj, xiEntire_order_conj, offAxis_pairs, offCriticalLine_partner}

\begin{proof}[Proof sketch]
The key identity is $\Xi(\bar z)=\overline{\Xi(z)}$, and it consumes \emph{both} the
conjugation symmetry $\overline{\Lambda_0(\bar s)}=\Lambda_0(s)$ and the functional
equation $\Lambda_0(1-s)=\Lambda_0(s)$; neither alone suffices, because
$z\mapsto\bar z$ is the composite $s\mapsto1-\bar s$ of the two. Given it, $\Xi(z)=0$
forces $\Xi(\bar z)=0$, and $\bar z\ne z$ precisely when $z$ is off the axis, which by
Lemma~\ref{lem:crit} is off the critical line. Orders agree by transporting the local
factorisation $\Xi(w)=(w-z)^m h(w)$ through conjugation.
\end{proof}

\begin{theorem}[the strip count's parity]\label{thm:stripparity}
For a bounded height band with no event at either endpoint, the parity of the total
multiplicity of \emph{all} events in the band --- on the critical line or off it ---
is the readout's sign change.
\end{theorem}
\leanref{bounded_strip_zeros_finite, even_sum_of_involution,
offline_strip_multiplicity_even, full_count_parity_eq_weld_parity}

\begin{proof}[Proof sketch]
The band is a compact rectangle in the ordinate variable and $\Xi$ is entire and not
identically zero, so it holds finitely many events. Off the axis, conjugation is a
fixed-point-free involution preserving the order, so the off-line zeros split into
two-element orbits, each contributing $2m$ to the total --- an even number, whatever
$m$ is. The total multiplicity is therefore congruent mod $2$ to the on-line total,
and the on-line half is Theorem~\ref{thm:count}.
\end{proof}

\noindent
Theorem~\ref{thm:count} is stated and proved in \S\ref{sec:count}, and its proof uses
nothing from this section, so the forward reference carries no circularity. That is
mechanically checkable and was checked: none of the twelve identifiers defined in
\lean{CarrierOffLinePairing.lean}, the module carrying this section, occurs anywhere in
\lean{CarrierIntervalCount.lean}, the module carrying \S\ref{sec:count}.

\begin{proposition}[the boundary]\label{prop:boundary}
Every event in the band is interior: $0<\operatorname{Re}s<1$.
\end{proposition}
\leanref{xiEntire_ne_zero_of_im_eq_neg_half, xiEntire_ne_zero_of_im_eq_half,
stripBand_zeros_interior, stripBand_zeros_in_open_strip}

\begin{proof}[Proof sketch]
In the ordinate variable $\operatorname{Re}s=\tfrac12-\operatorname{Im}z$, so the two
edges $\operatorname{Re}s\in\{0,1\}$ are $\operatorname{Im}z=\pm\tfrac12$. On the edge
$\operatorname{Re}s=1$ non-vanishing is $\zeta(1+it)\ne0$; the edge
$\operatorname{Re}s=0$ follows from it by the functional equation. The one point where
$\Lambda$ has a pole rather than a value is handled by the entire form: at $z=-i/2$,
i.e.\ $s=1$, the factor $\tfrac12-iz$ vanishes and
$\Xi(-i/2)=0\cdot\Lambda_0(1)-1=-1\ne0$.
\end{proof}

%======================================================================
\section{The interval count}\label{sec:count}

\begin{theorem}[interval count parity]\label{thm:count}
For $a\le b$ with no event at either endpoint,
\[
  \ord(a,b)\ \text{is even}\iff u_\Delta(a)\,u_\Delta(b)>0 ,
\]
where $\ord(a,b)$ is the total multiplicity of the on-line events in $(a,b)$.
\end{theorem}
\leanref{zeros_finite, sign_const_of_no_zeros, zeroCount_split, interval_sign_parity,
zeroCount_even_iff, weldReadout_interval_parity}

\begin{proof}[Proof sketch]
Strong induction on the number of events in $(a,b)$, which is finite. With none, the
readout is continuous and non-vanishing on $[a,b]$, hence of constant sign, so the
product is positive and the count is $0$. Otherwise isolate one event $t_0$, split
$(a,b)$ at points either side of it at which the readout does not vanish, apply
Theorem~\ref{thm:signflip} across $t_0$ to pick up a factor $(-1)^{m}$, and apply the
inductive hypothesis to the two shorter intervals; the counts add and the signs
multiply.
\end{proof}

\begin{theorem}[scope]\label{thm:scope}
The on-line count equals the full strip count for a band exactly when no event in
that band lies off the critical line.
\end{theorem}
\leanref{stripCount_eq, stripCount_eq_zeroCount_iff}

\noindent
This is stated so that the index certified by the on-line chain is not confused with
the classical index. The machinery nowhere assumes the two agree; where the
distinction matters it is flagged, in \S\ref{sec:notclaimed}(iv).

%======================================================================
\section{A count of one without a count}\label{sec:countone}

\begin{theorem}\label{thm:countone}
Let $a<b$. If $\Xi'$ does not vanish on $[a,b]$ and $u_\Delta(a)u_\Delta(b)<0$, then
$(a,b)$ contains exactly one event and it is simple.
\end{theorem}
\leanref{zeroCount_eq_one_of_deriv_ne_zero, subsingleton_zeros_of_deriv_ne_zero,
exists_zero_of_sign_change, order_eq_one_of_deriv_ne_zero}

\begin{proof}[Proof sketch]
Three steps. \emph{At least one:} the readout is continuous and changes sign, so the
intermediate value theorem supplies an event. \emph{At most one:} if there were two,
Rolle applied to $\Xi\restriction\R$ between them would give a point of $[a,b]$ with
$\Xi'=0$, contrary to hypothesis. \emph{Simple:} at the event $\gamma$ we have
$\Xi(\gamma)=0$ and $\Xi'(\gamma)\ne0$, which is order exactly one.
\end{proof}

\begin{remark}[which function's derivative, and why it matters]\label{rem:whichderiv}
The sign hypothesis is the carrier readout's and the derivative hypothesis is
$\Xi$'s, and the two are not interchangeable. By
Proposition~\ref{prop:xiclear} and Theorem~\ref{thm:weld},
\[
  u_\Delta(t)\;=\;\rho_\Delta(t)\,\Xi(t),\qquad
  \rho_\Delta(t)\;=\;\frac{\lvert w_\Delta(t)\rvert}{\tfrac14+t^{2}}\;>\;0 ,
\]
so $u_\Delta$ and $\Xi$ have the same zeros and the same sign, and the sign
hypotheses transfer exactly, at every carrier scale --- that is the content of
\lean{sign_product_iff}. Their derivatives, however, satisfy
$u_\Delta'=\rho_\Delta'\,\Xi+\rho_\Delta\,\Xi'$, and away from an event the first
term need not vanish; so $u_\Delta'\ne0$ on $[a,b]$ and $\Xi'\ne0$ on $[a,b]$ are
genuinely different hypotheses, and it is the latter that is formalised.

At an event the distinction disappears: $\Xi(\gamma)=0$ forces
$u_\Delta'(\gamma)=\rho_\Delta(\gamma)\,\Xi'(\gamma)$ with $\rho_\Delta(\gamma)>0$, so
$u_\Delta'(\gamma)\ne0\iff\Xi'(\gamma)\ne0$ and ``simple'' means the same thing read
either way. What is \emph{not} formalised here is the version of
Theorem~\ref{thm:countone} whose derivative hypothesis is $u_\Delta'\ne0$; it is true,
by the same Rolle argument applied to $u_\Delta$, but it is not what the Lean
statements consume, and we state the hypothesis as it is checked.
\end{remark}

\begin{remark}[why not a density bound]\label{rem:density}
The apparent alternative --- ``a short interval holds at most two events'' --- is not
a weaker statement. Recall the Riemann--von Mangoldt decomposition of the counting
function $N(T)$ of zeros with ordinate in $(0,T]$,
\[
  N(T)\;=\;\frac{\vartheta(T)}{\pi}+1+S(T),\qquad
  S(T)\;=\;\frac1\pi\arg\zeta(\tfrac12+iT),
\]
the argument taken by continuous variation; $S$ is the entire difficulty, and it is
what every classical count has to control. Since
$N(b)-N(a)=\bigl(\vartheta(b)-\vartheta(a)\bigr)/\pi+S(b)-S(a)$ and the main term is
$w\log(t/2\pi)/2\pi+O(w^{2})$ for a bracket of width $w=b-a$, bounding the count by
two amounts to
\[
  S(b)-S(a)\;\le\;2-\frac{w\log(t/2\pi)}{2\pi}+O(w^{2}).
\]
So the ``density bound'' is a bound on the variation of $S$ over a short interval ---
the same quantity, not an easier one. And approaching it from the other side does not
help either: Jensen's formula on a disc of radius comparable to the spacing gives
$O(\log t)$ events, not two. The hypothesis of Theorem~\ref{thm:countone} is pointwise
and computable instead, which is the whole reason it is the one used.
\end{remark}

%======================================================================
\section{The chain and the capstone}\label{sec:capstone}

\begin{theorem}[the index from a certified chain]\label{thm:chain}
Let $t_0<t_1<\dots<t_n$ be such that across each cell the readout changes sign and
$\Xi'$ does not vanish. Then $\ord(t_0,t_n)=n$; and for $1\le i\le n$ the counts
to either side of the $i$-th cell are $i-1$ and $i$, and that cell holds a unique
event, which is simple.
\end{theorem}
\leanref{zeroCount_chain, ith_event_of_chain}

\begin{proof}[Proof sketch]
Theorem~\ref{thm:countone} applied to each cell gives exactly one event there, simple,
so each cell contributes $1$; the counts add over the ladder by
\lean{zeroCount_split}, giving $\ord(t_0,t_n)=n$ and $\ord(t_0,t_{i-1})=i-1$,
$\ord(t_0,t_i)=i$.
\end{proof}

\begin{theorem}[the capstone]\label{thm:capstone}
Under the hypotheses of Theorem~\ref{thm:chain}, for any $t_{i-1}\le c\le d\le t_i$
with $u_\Delta(c)u_\Delta(d)<0$ there is $\gamma\in(c,d)$ with
$\zeta(\tfrac12+i\gamma)=0$, and $\gamma$ is the unique event of the $i$-th cell.
\end{theorem}
\leanref{ith_event_enclosed, sign_product_iff, ith_event_enclosed_carrier}

\begin{proof}[Proof sketch]
By Theorem~\ref{thm:chain} the $i$-th cell $(t_{i-1},t_i)$ holds a unique event
$\gamma$. The sub-bracket $[c,d]$ lies inside that cell and the readout changes sign
across it, so by the intermediate value theorem it contains an event; uniqueness in
the cell forces that event to be $\gamma$. Proposition~\ref{prop:zeros} converts
``event'' to $\zeta(\tfrac12+i\gamma)=0$.
\end{proof}

\noindent
Every hypothesis is a finite computation: $n+3$ sign evaluations of the readout, and
non-vanishing of $\Xi'$ on each cell (Remark~\ref{rem:whichderiv}).

%======================================================================
\section{The jet enclosure}\label{sec:jets}

Theorem~\ref{thm:capstone} encloses the $i$-th event in a bracket $(c,d)$ obtained
from sign evaluations alone. That bracket is as narrow as the bisection that produced
it and no narrower. This section sharpens it to an interval of width $2\varepsilon$
using the logarithmic jets of $\Xi$ at a nearby base point.

Fix a base point $y_0$ and set $d=y_0-\gamma$ for the event $\gamma$ in question. The
jets are the power sums
\[
  A_m\;=\;\sum_{j}\,(y_0-\gamma_j)^{-m},
\]
and the classical observation is that $A_m/A_{m+1}\to d$ when $\gamma$ is the ordinate
nearest $y_0$. What we need is not the limit but an \emph{enclosure} valid at finite
$m$, with every error that entered the computation carried explicitly.

\begin{theorem}[the ratio enclosure]\label{thm:ratio}
Let $d\in\C^{\times}$, $m\in\N$, and let $P,Q\in\C$ and $\alpha,\beta\ge0$ satisfy
\[
  \bigl\lvert P-d^{-m}\bigr\rvert\le\alpha,\qquad
  \bigl\lvert Q-d^{-(m+1)}\bigr\rvert\le\beta,\qquad
  \beta<\lvert d\rvert^{-(m+1)} .
\]
Then
\[
  \Bigl\lvert \frac{P}{Q}-d \Bigr\rvert
  \;\le\;\frac{\alpha+\lvert d\rvert\,\beta}{\lvert d\rvert^{-(m+1)}-\beta}.
\]
If moreover $d$ is real and the right-hand side is $\le\varepsilon$, then
$\lvert\operatorname{Re}(P/Q)-d\rvert\le\varepsilon$.
\end{theorem}
\leanref{ratio_enclosure, re_enclosure}

\begin{proof}[Proof sketch]
The numerator identity is
\[
  P-dQ\;=\;\bigl(P-d^{-m}\bigr)-d\bigl(Q-d^{-(m+1)}\bigr),
\]
since $d\cdot d^{-(m+1)}=d^{-m}$; so
$\lvert P-dQ\rvert\le\alpha+\lvert d\rvert\beta$. For the denominator, the reverse
triangle inequality gives
$\lvert Q\rvert\ge\lvert d\rvert^{-(m+1)}-\beta>0$, which in particular makes $Q\ne0$.
Dividing, $\lvert P/Q-d\rvert=\lvert P-dQ\rvert/\lvert Q\rvert$ is at most the stated
quotient. Taking real parts can only shrink the error: for $d$ real,
$\lvert\operatorname{Re}(P/Q)-d\rvert=\lvert\operatorname{Re}(P/Q-d)\rvert
\le\lvert P/Q-d\rvert$.
\end{proof}

\noindent
The hypothesis $\beta<\lvert d\rvert^{-(m+1)}$ is exactly what keeps $Q$ away from
zero, and it is checkable. The two budgets $\alpha,\beta$ are the \emph{only} route by
which any approximation enters: the finite-bank truncation error and the jet tail both
land there, and the theorem is indifferent to which contributed what. This is the
ancestor of the argument in Gargantini and Henrici \cite{GargantiniHenrici1971}, where
circular arithmetic on the same power-sum ratios yields inclusion discs for polynomial
roots; the contribution here is that the two budgets are separated and the statement is
machine-checked.

The counting routine enters the next statement as an interface, and the interface has
exactly one requirement. It is worth isolating, because it is the whole of what the
theorem asks of a count.

\begin{definition}[counting routine]\label{def:routine}
A \emph{counting routine} is a function $R:\R\to\N$ that is \emph{additive along the
axis between non-events}: whenever $x\le y$ and $\Xi(x),\Xi(y)\ne0$,
\[
  R(y)\;=\;R(x)+\ord(x,y).
\]
\end{definition}
\leanref{CountRoutine}

\noindent
Nothing else is required --- not that $R$ counts from zero, not that it agrees with
$N(t)$, not that it is computed in any particular way. Any procedure whose increments
are the true event counts qualifies, and the theorem below is stated relative to
whichever one is supplied. This is what \S\ref{sec:notclaimed}(iv) means when it says
the certified index is the index \emph{of the count that was given}.

\begin{theorem}[the product theorem]\label{thm:product}
Let $R$ be a counting routine with $R(a)+1=i$ and $R(b)=i$, let $\Xi(a),\Xi(b)\ne0$
and $a\le b$. Then $\Xi$ has exactly one zero $\gamma$ in $(a,b)$; and if, for that
$\gamma$, the data $y_0,P,Q,m,\alpha,\beta,\varepsilon$ satisfy the hypotheses of
Theorem~\ref{thm:ratio} with $d=y_0-\gamma$, then
\[
  \gamma\;\in\;\bigl[\,y_0-\operatorname{Re}(P/Q)-\varepsilon,\;\;
                      y_0-\operatorname{Re}(P/Q)+\varepsilon\,\bigr].
\]
\end{theorem}
\leanref{gamma_enclosure, count_eq_one_of_endpoints, unique_event_of_count_eq_one,
one_le_order_of_zero}

\begin{proof}[Proof sketch]
Additivity (Definition~\ref{def:routine}) applied to $a\le b$ gives
$R(b)=R(a)+\ord(a,b)$, and the two supplied counts then force $\ord(a,b)=1$. Now
$\ord(a,b)$ is by definition the \emph{sum} of the multiplicities of the events in
$(a,b)$, and every event has multiplicity at least $1$; a sum of positive integers
equal to $1$ has exactly one summand, equal to $1$. So there is exactly one zero
$\gamma$ of $\Xi$ in $(a,b)$, and it is simple --- note that this comes from the count
alone, with no derivative hypothesis and no appeal to
Theorem~\ref{thm:countone}. Applying Theorem~\ref{thm:ratio} to that $\gamma$ with
$d=y_0-\gamma$ bounds $\lvert\operatorname{Re}(P/Q)-(y_0-\gamma)\rvert$ by
$\varepsilon$, which rearranges to the stated interval.
\end{proof}

\noindent
Two features of this statement carry the weight. The counting routine enters only as
an interface: any procedure satisfying Definition~\ref{def:routine} may be substituted,
and the theorem does not ask how it is computed. And the event's multiplicity is
\emph{certified by the count}, not assumed --- the passage from $R(b)-R(a)=1$ to ``one
simple zero'' is the sum-of-orders argument in the proof above, so no simplicity
hypothesis is imported here or anywhere else in the paper.

It is worth being precise about what this buys, because it is easy to overstate. The
enclosure of Theorem~\ref{thm:capstone} obtains its count from the chain, whose
hypotheses are sign evaluations and non-vanishing of $\Xi'$; the enclosure here obtains
it from $R$, and inherits whatever $R$ certifies. The two are independent routes to the
same integer, and the paper proves both rather than choosing.

\subsection*{Measured conditioning}
The enclosure is a theorem for any $\alpha,\beta$; how \emph{small} $\varepsilon$ gets
is a question about the jets, and it is governed by
\[
  \kappa\;=\;\frac{\lvert y_0-\gamma_1\rvert}{\lvert y_0-\gamma_2\rvert},
\]
the distance ratio from the base point to the nearest and second-nearest ordinates.
The error of $A_m/A_{m+1}$ should contract by $\kappa$ per order. Measured against the
true ordinates at $50$-digit precision
(\lean{extractor/checks/jet_sweep.py}), over ten base points with $\kappa$ ranging from
$0.0143$ to $0.3426$, the mean contraction over $m=5,\dots,15$ agreed with $\kappa$ to
between $1.2\%$ and $2.7\%$ --- confirmed at all ten, against a pre-registered
threshold of $5\%$.

The law degrades exactly where it should. At a base point placed between two ordinates
at nearly equal distance ($\kappa=0.882$) the measured contraction was $1.29$: the
error does not decay at all, and it oscillates, because the two nearest contributions
enter with opposite signs and there is no dominant term to extract. This is not a
defect of Theorem~\ref{thm:ratio}, which remains true; it is the statement that a
badly placed base point yields a large $\varepsilon$. A large $\varepsilon$ is a wide
interval, never a wrong one.

\begin{remark}
This entire section is \emph{off} the certification path. The enclosure of
Theorem~\ref{thm:capstone} --- which ordinate, and that it is unique --- is complete
without any jet, any tail bound, or any conditioning hypothesis. What the jets buy is
narrowness, and if they are badly conditioned the only consequence is that the
interval is wide.
\end{remark}

%======================================================================
\section{The extractor}\label{sec:extractor}

\subsection{The head law}
The lane split is not free: it manufactures a stationary point, and that stationary
point sets the cost. Writing the $n$-th term of the lane-balanced bank at height $T$,
its phase is
\[
  \varphi(n)\;=\;\pi n-T\log(\Delta n),
\]
the $\pi n$ supplied by the alternating lane sign and the $T\log(\Delta n)$ by the
phasor. Then $\varphi'(n)=\pi-T/n$, which vanishes at
\[
  n_\ast \;=\; T/\pi .
\]
Below $n_\ast$ the terms oscillate and cancel; at $n_\ast$ they align. A head window
that is tapered to zero at its edge therefore has to carry $n_\ast$ strictly inside its
flat core, or the truncation cuts the bank exactly where it is not cancelling. With a
window whose flat core is the first half, this forces
\[
  X \;>\; 2T/\pi .
\]
\emph{The head is linear in the ordinate, not $O(\sqrt T)$.} This is the single
structural cost of the construction, and it is the price of the lane split: the
Riemann--Siegel head is $O(\sqrt T)$ precisely because its phase has no such
stationary point in range.

This is a \emph{design} rule: it says where to put the window, not what accuracy a
given head achieves, and the two are worth separating. The script
\lean{extractor/checks/head_threshold.py} measures the smallest head $X_{\min}$
reaching a fixed relative accuracy, as a multiple of $2T/\pi$:

\begin{center}\footnotesize
\begin{tabular}{@{}rrrr@{}}
\toprule
$T$ & $\mathrm{tol}=10^{-6}$ & $10^{-9}$ & $10^{-12}$\\
\midrule
$500$   & $1.12$ & $1.37$ & $1.69$\\
$1500$  & $1.00$ & $1.10$ & ---\\
$5000$  & $0.95$ & $0.99$ & ---\\
$15000$ & $0.96$ & $0.98$ & ---\\
\bottomrule
\end{tabular}
\end{center}

\noindent
So the accuracy threshold tracks $2T/\pi$ to within a few percent at large $T$, but it
is \emph{not} approached from one side: at a loose tolerance and large height the rule
is mildly conservative (ratio $0.95$), and at a tight tolerance and small height it is
not enough on its own (ratio $1.69$). The dashes are not omissions. They record that a
\lean{float64} bank \emph{cannot} reach $10^{-12}$ relative accuracy at those heights at
any head length: the phase $T\log(\Delta n)$ carries absolute error of order
$T\log(X)\cdot2.2\times10^{-16}$, which floors the bank's relative accuracy near
$10^{-13}$ at $T\sim10^{3}$ and near $3\times10^{-10}$ at $T=1.5\times10^{4}$. The
\lean{--tol} parameter therefore sets the head \emph{length} via the law below; it does
not by itself deliver the corresponding accuracy, and past the floor a longer head buys
nothing. This is the same arithmetic limit that (O1) proposes to remove.

Past the threshold the convergence is fast and the right variable is the excess.
Writing $s=\sqrt{\pi X/2-T}$, the relative error of the truncated bank falls like
$\exp(-cs)$. Measured at $60$-digit working precision against
$\Delta^{-s}(1-2^{1-s})\zeta(s)$
(\lean{extractor/checks/head_law_decay.py}, \lean{extractor/checks/head_law_moments.py}),
over heights $T=143,500,1500$, carrier scales $\Delta\in\{\pi/3,1\}$ and moment orders
$r=0,\dots,4$:
\[
  -\ln(\mathrm{err})/s\;\in\;[1.400,\,1.485]\quad\text{for } s\gtrsim10,
\]
across thirteen orders of magnitude of error. The dependence on the moment order is
slight --- the five orders $r=0,\dots,4$ span $0.085$ at fixed $s$ --- and the
dependence on $\Delta$ is nil to every digit printed. For $s$ below about $10$, that is
for a head only just past the threshold, the effective rate is \emph{larger} (up to
$2.7$ at $s\approx5.6$), so the exponential law is approached from above and using a
constant from the asymptotic band is the conservative choice. The implementation sets
$X=2(T+K)/\pi$ with $K=(\ln(1/\mathrm{tol})/1.43)^2$; the constant $1.43$ sits inside
the measured band, and at $\mathrm{tol}=10^{-16}$ the difference between $1.43$ and the
band's lower end $1.40$ is thirteen terms of head.

\begin{remark}[why this had to be measured in high precision]
A double-precision run cannot see this law at all. The bank summed in \lean{float64}
saturates near $10^{-13}$ relative, so its error stops falling once the head passes the
threshold and the measured $-\ln(\mathrm{err})/s$ then decays like $1/s$ --- an artifact
of the arithmetic, not a property of the bank. The figures above were therefore taken
with the bank summed in \texttt{mpmath}; the reported thirteen orders of magnitude are
what $60$ digits of working precision will support, and no claim is made beyond them.
\end{remark}

\subsection{Algorithm}
The extractor takes an index $i$ and returns an interval.

\begin{description}
\item[Stage 1 (anchor).] Invert the archimedean winding $\vartheta$ for a coarse
  ordinate: solve $\vartheta(T)/\pi+1=i$ by Newton on a smooth monotone function. Cost
  is polylogarithmic and the result is within a few mean spacings.
\item[Stage 2 (bracket).] Bisect on the predicate $\mathrm{Count}(t)\ge i$, which is
  monotone in $t$. This is the one place the count enters, and it enters only as an
  \emph{interface}: any routine satisfying Definition~\ref{def:routine} may be
  substituted, and the certification in Stage~5 is stated relative to whatever it
  returns (\S\ref{sec:capstone}, \S\ref{sec:jets}). Bisection rather than a fixed-step
  march is not a detail --- see \S\ref{sec:validation} on Lehmer's pair.
\item[Stages 3--4 (moments and jets).] Write $v$ for the vector of site phasors
  $\bigl(x_n^{-s}\bigr)_n$ over the head, $\ell$ for the lane-balanced summation
  functional that forms the fiber from it (the alternating sum of
  Definition~\ref{def:fiber}), and $H$ for the log-weight operator of
  \S\ref{sec:contribution} (Theorem~\ref{thm:logweight}), so that $\ell(v)=F_\Delta(s)$
  truncated to the head. One
  pass over the head evaluates the bank and, simultaneously, the log-weighted moments
  $M_r=\ell(H^r v)$, since the weights $\log x_n$ are already formed. A Bell (cumulant) recursion converts the moments into
  the log-jets $A_m=\sum_j(y_0-\gamma_j)^{-m}$, and the ratio $A_m/A_{m+1}$ locates the
  nearest event with error $O(\kappa^m)$, $\kappa$ the ratio of distances to the
  nearest and next-nearest events. The whole jet vector costs one pass, not $m$.
\item[Stage 5 (certificate).] Narrow to a sign bracket around the jet estimate and
  return it. The returned interval is certified \emph{given} the count, which is
  exactly the shape of the product theorem, Theorem~\ref{thm:product} --- \emph{not}
  of the capstone. The distinction matters and is set out in (O1): the implementation
  takes the count route, which carries no derivative hypothesis at all.
\end{description}

\noindent A noise floor is estimated from the head tolerance and the bank magnitude; a
bracket narrower than the floor is reported as \emph{not certified} rather than
returned silently. That guard was added after an early version produced false
certificates by narrowing past its own arithmetic.

\subsection{Validation}\label{sec:validation}
The harness checks three things at every index: that the count interface actually
reached $\mathrm{Count}(a)=i-1$ and $\mathrm{Count}(b)=i$; that Stage~5 returned an
interval rather than hitting its noise floor; and that the interval contains the
reference ordinate. The reference is \texttt{mpmath}'s
\texttt{zetazero}, which is an independent implementation --- it is not derived from
anything in this paper.

\begin{remark}[two layers, and which one this table is]\label{rem:twolayers}
The word \emph{certified} is used in two senses in this paper and they should not be run
together. This section reports the \lean{float64} extractor, which \emph{locates}: it is
fast, it is validated against an independent reference at every index below, and its
internal guard is a measured noise floor, which is a heuristic. The certification proper
is a separate layer, \S\ref{sec:obligations}(O1): an Euler--Maclaurin evaluator with an
explicit remainder, which proves the signs, the non-vanishing of $\Xi'$ and the count on
the intervals this layer produces. The table below is evidence that the locator works;
it is not the certificate, and nothing in it is claimed to be.
\end{remark}

\begin{center}
\begin{tabular}{lrr}
\toprule
sample & indices & failures\\
\midrule
exhaustive $1\le i\le200$ & 200 & 0\\
random $201\le i\le 2\times10^{4}$ & 90 & 0\\
spot checks to $i=10^{6}$ & 5 & 0\\
Lehmer's pair $i=6709,6710$ & 2 & 0\\
\midrule
total & 297 & 0\\
\bottomrule
\end{tabular}
\end{center}

\noindent\emph{Lehmer's pair~\cite{Lehmer1956} is the interesting one.} The two ordinates are
\[
  \gamma_{6709}=7005.062866174921,\qquad \gamma_{6710}=7005.100564672656:
\]
a separation of $0.0377$ against a local
mean spacing of $0.8955$, i.e.\ $0.042$ of a mean spacing. A fixed-step bracket search
with any step comparable to the mean spacing walks straight over both; bisection on the
monotone count predicate does not, because it never needs to resolve the pair --- it
only needs the count to cross. Both are returned certified, and each contains its
reference ordinate. This was a genuine failure of an earlier version (a fixed-step
march at $0.30$ spacings returned $\gamma_{6711}$ for $i=6710$), and it is the reason
Stage~2 is stated as bisection on a predicate rather than as a search.

\emph{Carrier-scale invariance.} The whole construction is supposed to be
$\Delta$-free, and this is testable rather than assumed: running $i=1$ and $i=100$ on
$\Delta\in\{\pi/6,\pi/3,1,2,30\}$ --- a factor of $180/\pi\approx57$ between the
smallest and largest carrier scale --- returns the same ordinates to every digit
printed. That is Proposition~\ref{prop:zeros} as a measurement.

%======================================================================
\section{Eigenvalues and the rank drop: reading the abscissa on the
line}\label{sec:pencil}

Everything so far has used the readout's \emph{sign}, which is one bit per evaluation
and is blind to the abscissa by construction. The moment vector of Stage~3 carries
more than that, and this section extracts it. The claim is that the abscissa of a zero
--- whether it lies on the critical line at all --- is recoverable from data computed
\emph{only on the line}, as the imaginary part of an eigenvalue.

\subsection{The jets are a power sum over the fiber's zeros}
Write $G(y)=F_\Delta(\tfrac12+iy)$ for the fiber read in the ordinate variable: a
meromorphic function of $y$ whose zeros are those of $\zeta$ (non-trivial and trivial)
and of the lane factor, and whose one pole --- $\zeta$'s pole at $s=1$ --- sits at
$y=-\tfrac{i}{2}$. Near a real base point $y_0$, for $m\ge2$,
\begin{equation}\label{eq:powersum}
  A_m\;=\;\frac{(-1)^{m-1}}{(m-1)!}\,\bigl(\log G\bigr)^{(m)}(y_0)
  \;=\;\sum_j \mu_j\,d_j^{\,m},
  \qquad d_j=\frac{1}{y_0-y_j},
\end{equation}
the sum over the zeros \emph{and the pole} of $G$, with $\mu_j$ the multiplicity of a
zero and $\mu=-1$ at the pole. The series converges absolutely for $m\ge2$: the
non-trivial zeros number $O(\log T)$ per unit ordinate, so
$\sum_j\lvert y_0-y_j\rvert^{-m}<\infty$, and the trivial zeros, the lane nodes and the
pole sit at distance $\gtrsim y_0$ from a base point on the line. This is the $m$-th
logarithmic derivative of a Hadamard
product of genus one: for $m\ge2$ the exponential factor and the convergence factors
differentiate away and what remains is exactly the power sum. At order one it is
\emph{not} exact --- an earlier version of this section said it was, and that was an
erratum: $(\log G)'(y_0)$ carries, besides a symmetrically summed
$\sum_j\mu_j/(y_0-y_j)$, the constant $-i\log\Delta$ from the carrier scale and the
constants of the Hadamard factor and of $\Gamma_{\mathbb R}$. That the carrier scale
shows itself in the jets \emph{only} at order one, and there as the shift $\log\Delta$,
is Theorem~\ref{thm:logweight}(iii) read on the jets. For $m\ge2$ the jets are an exact
\emph{power sum} in the reciprocal distances, and \eqref{eq:powersum} is what makes the
rest of this section possible: a power sum determines its nodes. The extractor uses
$m\ge2$ only (its default order is $9$), and from a base point on the line the pole and
the trivial zeros are nodes at distance $\gtrsim y_0$, contributing $O(y_0^{-m})$.

The $A_m$ cost one pass. The bank evaluates $M_r=\ell(H^rv)$ for $r=0,\dots,2K+1$
simultaneously (Stage~3), and a Bell recursion converts $M_r\mapsto A_m$; no additional
evaluation of $\zeta$, and no differentiation of anything, is involved.

\subsection{The node dictionary}
Which $y_j$ occur is the point. $G$ is $\Delta^{-s}(1-2^{1-s})\zeta(s)$ at
$s=\tfrac12+iy$, so its zeros are those of $\zeta$ together with those of the lane
factor. Under $s=\tfrac12+iy$ a zero of $\zeta$ at $\rho=\beta+i\gamma$ sits at
\begin{equation}\label{eq:dict}
  y\;=\;\gamma-i\bigl(\beta-\tfrac12\bigr),
  \qquad\text{so}\qquad
  \operatorname{Im}y=0\iff\beta=\tfrac12 .
\end{equation}
\emph{The imaginary part of the node is the signed distance of the zero from the
critical line.} A node on the real axis is an on-line zero; a node off it reports its
own abscissa as $\beta=\tfrac12-\operatorname{Im}y$. Nothing in the computation left the
line to learn this.

\subsection{The pencil}
Recovering nodes from a power sum is classical: it is Prony's problem, taken here in
the numerically stable matrix-pencil form of Hua and Sarkar~\cite{HuaSarkar1990}.
Build the Hankel data matrix
$Y=(A_{1+i+j})_{0\le i\le K,\,0\le j\le K}$, take its singular value decomposition,
keep the $K_r$ dominant right singular directions $V$, and form the shifted pair
$(V_{1:},V_{:1})$; the eigenvalues of $V_{1:}^{+}V_{:1}$ are the $d_j$, and
$y_j=y_0-1/d_j$.

\subsection{What the rank counts}
The rank of $Y$ is the number of \emph{distinct} nodes. This is worth stating carefully
because it is easy to misread, and because it is the instrument's one structural
limitation. Measured directly on synthetic power sums with known nodes
(\lean{extractor/checks/rank_drop.py}):

\begin{center}
\begin{tabular}{lll}
\toprule
configuration & numerical rank & $\sigma_2/\sigma_1$\\
\midrule
$1,2,3,4$ distinct nodes & $1,2,3,4$ & --- \\
one node of multiplicity $2$ & $1$ & $1.3\times10^{-18}$\\
one node of multiplicity $1$ & $1$ & $1.3\times10^{-18}$\\
\bottomrule
\end{tabular}
\end{center}

\noindent
The last two rows are the limitation: a double node and a simple node give
\emph{bit-identical} singular-value profiles. Multiplicity enters \eqref{eq:powersum}
as the coefficient $\mu_j$, not as an extra exponential, so the rank cannot see it.
\textbf{The rank detects coincidence, not multiplicity.} A rank of $1$ where the count
says $2$ means two events have merged to within the resolution; it does not certify a
double zero, and we do not use it to.

The approach to a merger is quantitative. With two nodes at separation $\varepsilon$,
the measured $\sigma_2/\sigma_1$ was $4.58\times10^{-3},\ 4.24\times10^{-5},\
4.20\times10^{-7},\ 4.20\times10^{-9}$ at $\varepsilon=10^{-1},\dots,10^{-4}$: a clean
\[
  \sigma_2/\sigma_1\;\approx\;4.2\,\varepsilon^{2},
\]
so in double precision, where the floor is near $10^{-16}$, two events remain distinct
down to $\varepsilon\approx5\times10^{-9}$ and are indistinguishable below it. That
number is the resolution of the instrument, and it is a property of the arithmetic
rather than of the geometry.

\subsection{The calibration standard the fiber carries}\label{sec:lanecontrol}
An instrument that claims to measure the abscissa must be checked against zeros whose
abscissa is known and \emph{is not} $\tfrac12$. The fiber supplies them. The lane
factor $1-2^{1-s}$ vanishes exactly when $1-s\in\frac{2\pi i}{\log2}\Z$, that is on the
line $\operatorname{Re}s=1$, and under \eqref{eq:dict} these sit at
\begin{equation}\label{eq:lane}
  y\;=\;-\frac{2\pi k}{\log 2}-\frac{i}{2},\qquad k\in\Z:
  \qquad \operatorname{Im}y=-\tfrac12\ \text{exactly},\quad
  \text{spacing } \frac{2\pi}{\log2}=9.064720 .
\end{equation}
These are genuine off-line nodes, in closed form, carried by the object itself --- not
planted, not synthetic, and not adjustable. A working pencil must place them at
$\operatorname{Im}y=-\tfrac12$, equivalently report $\beta=1$ for them.

Measured (\lean{extractor/checks/pencil_lane_control.py}), at an anchor $y_0=17.5$
placed between the first two ordinates so that no node dominates, with $K=8$:

\begin{center}
\begin{tabular}{lll}
\toprule
returned node & identification & error\\
\midrule
$18.129441-0.500000\,i$ & lane, $k=-2$ of \eqref{eq:lane} & $2.0\times10^{-14}$\\
$14.134707+0.000026\,i$ & $\gamma_1$, on-line & $1.8\times10^{-5}$\\
$21.022108+0.000122\,i$ & $\gamma_2$, on-line & $6.8\times10^{-5}$\\
\bottomrule
\end{tabular}
\end{center}

\noindent
From one moment vector computed on the critical line, the pencil returns $\beta=1$ for
the lane zeros to fourteen digits and $\beta=\tfrac12$ for $\zeta$'s zeros, and it
separates the two populations. At $y_0=12$ the $k=-1$ lane node comes back as
$9.064720-0.500005\,i$, against the predicted $9.064720-\tfrac{i}{2}$.

This control is what makes the abscissa reading falsifiable rather than decorative, and
it has already done work: it refuted two successive implementations of the pencil --- a
wrong data-matrix shape, and a spurious conjugation that returned $\overline{y}$ in
place of $y$ --- both of which had produced clean-looking certificates while inverting
precisely the off-line information the pencil exists to read.

\subsection{Lehmer's pair, resolved}
The hardest close pair in the low range is $\gamma_{6709},\gamma_{6710}$, separated by
$0.0377$, or $0.042$ of the local mean spacing. At an anchor between them
(\lean{extractor/checks/lehmer_rank.py}), with $K=4$, the pencil returns rank $2$ with
$\sigma_2/\sigma_1=1.6\times10^{-10}$ --- far above the floor --- and the two nodes
\[
  7005.100564673+2.1\times10^{-12}i,\qquad
  7005.062866175-2.7\times10^{-12}i,
\]
matching the reference ordinates to $9.1\times10^{-13}$ and $3.2\times10^{-11}$, with
both abscissae $\beta=\tfrac12$ to $3\times10^{-12}$. So the pair is resolved as two
distinct events and both are certified on-line by the instrument, at a separation forty
times finer than would trouble it.

\subsection{The singular values, certified}\label{sec:certrank}
The rank reading above was a measurement. It need not be, and making it a certificate is
worth doing because the rank is the one part of this instrument that carries a
\emph{count}, and counts are what the rest of the paper certifies.

The chain is entirely within the machinery of \S\ref{sec:obligations}, with one new
element. Certified values of $\Xi$ give certified \emph{Taylor coefficients} at an anchor
$c$ by the trapezoid rule on a circle $\lvert z-c\rvert=R$: the $n$-point rule returns the
aliased coefficient $\sum_{i\ge0}a_{k+in}R^{k+in}$, and bounding $\lvert\Xi\rvert\le M'$ on
a larger disc of radius $R'>R$ gives $\lvert a_j\rvert\le M'/R'^{\,j}$ by Cauchy, hence
\[
  \Bigl\lvert a_kR^k-\widehat a_k\Bigr\rvert
  \;\le\;\frac{M'\,(R/R')^{k+n}}{1-(R/R')^{n}} ,
\]
an explicit geometric bound. The point of routing through $\Xi$ rather than $\Xi'/\Xi$ is
that $\Xi$ is \emph{entire}: no lower bound on $\lvert\Xi\rvert$ is needed anywhere, and
the region bound (O5) supplies $M'$. Complex interval arithmetic then carries the
coefficients through $f_k=\Xi^{(k)}/\Xi$ and the cumulant recursion to certified jets
$A_m$; Weyl's inequality,
$\lvert\sigma_i(H)-\sigma_i(H_{\text{true}})\rvert\le\lVert E\rVert_2\le\lVert E\rVert_F$,
turns those into certified singular values.

Measured (\lean{extractor/checks/certified_rank.py}), at three anchors with $K=3$: the
error norm $\lVert E\rVert_F$ came out between $3.7\times10^{-29}$ and
$5.3\times10^{-27}$, every singular value is enclosed to that width, and \emph{every
consecutive pair of enclosures separates} --- so the ordering of the spectrum is proved,
not observed. For the anchor $c=17{.}5$: $\sigma_1=6.773\times10^{-1}$,
$\sigma_2=8.088\times10^{-2}$, $\sigma_3=1.970\times10^{-3}$,
$\sigma_4=1.992\times10^{-5}$, each with radius $2.6\times10^{-28}$.

\subsection{Propose and verify: the node locations, certified}\label{sec:proposeverify}
Theorem~\ref{thm:rankcert} certifies the rank for the contour-restricted moments. It says
nothing about \emph{where} the nodes are, and that is what the pencil supplies: the node
locations, abscissa included. Those need not be trusted either --- they can be
\emph{verified}:

\begin{enumerate}
\item[(1)] the contour certifies the count $K$ in the band, with no pencil involved;
\item[(2)] the pencil, run on certified jets, \emph{proposes} $K$ nodes with abscissae;
\item[(3)] the certified machinery of \S\ref{sec:obligations}(O1) verifies each proposal
--- a certified sign change with $\Xi'\ne0$ on a small interval about it;
\item[(4)] if the number verified equals $K$, every node is accounted for.
\end{enumerate}

\noindent
The tail never enters: step (1) counts and step (3) locates, and neither consults the
pencil's own error. The script \lean{checks/pencil_verify.py} measures this at the
anchors the extractor actually uses, with $K=2$ and a $\pm10^{-4}$ verification window:

\begin{center}\footnotesize
\begin{tabular}{@{}rrrrl@{}}
\toprule
$i$ & proposed node & $\operatorname{Im}y$ & node error & verified\\
\midrule
$1$ & $14.1347251374$ & $2.8\times10^{-36}$ & $4.3\times10^{-9}$ & yes\\
$10$ & $49.7738359156$ & $-2.1\times10^{-29}$ & $3.4\times10^{-6}$ & yes\\
$100$ & $236.524229631$ & $4.1\times10^{-26}$ & $3.4\times10^{-8}$ & yes\\
$1000$ & $1419.42248108$ & $-6.5\times10^{-24}$ & $1.3\times10^{-7}$ & yes\\
\bottomrule
\end{tabular}
\end{center}

\noindent
All four verify, and $\operatorname{Im}y$ --- the abscissa offset --- comes back at
$10^{-24}$ or smaller in every case.

\begin{remark}[the conditioning is not incidental]
A first attempt ran this at $c=17{.}5$, midway between $\gamma_1$ and $\gamma_2$, and
\emph{all} proposals failed verification: the pencil returned $14.068$, $20.872$,
$27.666$ against true ordinates $14.135$, $21.022$, $25.011$. That is the
$\kappa\approx0.88$ case of \S\ref{sec:jets} --- no nearest node dominates --- and the
failure is the conditioning law doing exactly what it says. The extractor never uses such
an anchor: Stage~2 bisects first, which is what makes Stage~4 well conditioned. The
propose-and-verify pipeline inherits that dependence, and a verification window must be
scaled to the conditioning, not fixed.
\end{remark}

\subsection{The rank, certified}\label{sec:rankcertified}
An earlier version of this section left the \emph{rank} uncertified. The obstruction was
that $A_m=\sum_j d_j^{\,m}$ runs over \emph{all} zeros, so the exact Hankel matrix is not
of finite rank and ``rank $K$'' is only a statement about a threshold; certifying it
looked to require an explicit Riemann--von Mangoldt constant to bound the tail
$\sum_{\lvert c-\gamma_j\rvert>D}d_j^{\,m}$.

No such constant is needed. The tail is an artifact of summing over all zeros, and
restricting to a contour removes it by construction.

\begin{theorem}[the rank]\label{thm:rankcert}
Let $C$ be a contour carrying no zero of $\Xi$, enclosing exactly $K$ zeros of $\Xi$, all
simple and at distinct points. Put
\[
  s_m\;=\;\frac{1}{2\pi i}\oint_C z^{m}\,\frac{\Xi'(z)}{\Xi(z)}\,dz ,
  \qquad m=0,1,2,\dots
\]
Then $s_m=\sum_{j=1}^{K}\gamma_j^{\,m}$, and the Hankel matrix
$H=(s_{i+j})_{0\le i,j\le K}$ has rank exactly $K$.
\end{theorem}

\begin{proof}
$z^m$ is entire, so by the residue theorem the only contributions are the zeros inside
$C$; each being simple, $s_m=\sum_j\gamma_j^{\,m}$. Let $V$ be the $(K+1)\times K$ matrix
$V_{i,k}=\gamma_k^{\,i}$. Then $s_{i+j}=\sum_k\gamma_k^{\,i}\gamma_k^{\,j}$, that is
$H=VV^{\mathsf T}$. $V$ is a Vandermonde in $K$ \emph{distinct} nodes with $K+1$ rows,
hence of full column rank $K$. Over $\mathbb C$ that is not yet enough on its own ---
$VV^{\mathsf T}$ is not $VV^{*}$, and $AA^{\mathsf T}$ can have smaller rank than $A$
--- but $V^{\mathsf T}$ has full row rank $K$, so it maps $\mathbb C^{K+1}$ \emph{onto}
$\mathbb C^{K}$, on which $V$ is injective; hence
$\operatorname{rank}H=\operatorname{rank}(VV^{\mathsf T})=K$.
\end{proof}

\noindent
Each hypothesis is discharged by machinery already built: no zero on $C$ by a certified
positive lower bound for $\lvert\Xi\rvert$ there --- sample, then fill the gaps with
$\lvert\Xi(z)\rvert\ge\lvert\Xi(z_k)\rvert-\tfrac{L}{2}M/\rho$, the same device as the
sound winding criterion; exactly $K$ zeros inside by the certified argument principle;
simple and distinct by \S\ref{sec:bandverify}.

Measured (\lean{extractor/checks/rank_certified.py}):

\begin{center}\footnotesize
\begin{tabular}{@{}rrrlll@{}}
\toprule
centre & $R$ & $K$ & $s_0$ & $s_1$ & Hankel rank\\
\midrule
$14.1347$ & $0.30$ & $1$ & $1.0$ & $14.1347251417$ & $1$\\
$17.6$ & $3.60$ & $2$ & $2.0$ & $35.1567647805$ & $2$\\
$19.5$ & $6.00$ & $3$ & $3.0$ & $60.1676223607$ & $3$\\
\bottomrule
\end{tabular}
\end{center}

\noindent
$s_0$ returns the count as an integer, $s_1$ matches the sum of the reference ordinates to
twelve digits, and the trailing singular value falls to between $10^{-18}$ and $10^{-20}$
--- the arithmetic floor --- in every case. On the first contour, which passes within
$0.3$ of $\gamma_1$, the lower bound for $\lvert\Xi\rvert$ needed $n=8192$ samples to come
out positive ($3.6\times10^{-4}$); a contour near a zero costs sampling, which is the only
price this route pays.

\begin{remark}[what changed, and what did not]
The full-jet Hankel is still not of finite rank, and its numerical rank is still a
threshold statement --- that was never wrong. What was wrong was the inference that
certifying \emph{a} rank therefore required a zero-density input. The rank question is
well posed for the contour-restricted moments and ill posed for the unrestricted ones, and
posing it correctly removes the need for the constant rather than supplying it.
\end{remark}

\subsection{The inertia: what the sign of the Hankel spectrum counts}\label{sec:inertia}
Theorem~\ref{thm:rankcert} restricts the moments to a contour, $s_m=\sum_j y_j^m$ over the
$K$ nodes enclosed, and reads the rank. The same matrix carries a second integer.

\begin{lemma}[inertia]\label{lem:inertia}
Let the enclosed nodes consist of $r$ real ones $y_1,\dots,y_r$ and $q$ conjugate pairs
$z_j,\bar z_j$, all distinct, and let $H_n=(s_{i+j})_{0\le i,j<n}$ with $n\ge r+2q$. Then
the inertia of $H_n$ is $(r+q,\;q,\;n-r-2q)$: the number of negative eigenvalues is
exactly the number of off-line pairs enclosed, and $H_n\succeq0$ if and only if $q=0$.
\end{lemma}

\begin{proof}
With $v(y)=(1,y,\dots,y^{n-1})^{\mathsf T}$,
$H_n=\sum_i v(y_i)v(y_i)^{\mathsf T}+\sum_j\bigl[v(z_j)v(z_j)^{\mathsf T}
+v(\bar z_j)v(\bar z_j)^{\mathsf T}\bigr]$. Writing $v(z_j)=A_j+iB_j$ with $A_j,B_j$
real, the bracket is $2A_jA_j^{\mathsf T}-2B_jB_j^{\mathsf T}$ --- the minus sign is the
whole mechanism --- so $H_n=UJU^{\mathsf T}$ with
$U=[\,v(y_1)\cdots v(y_r)\;A_1\,B_1\cdots A_q\,B_q\,]$ and
$J=\operatorname{diag}(1,\dots,1,\,2,-2,\dots,2,-2)$. $U$ has full column rank $r+2q$ (a
Vandermonde in distinct nodes followed by an invertible change of columns), and
Sylvester's law of inertia gives the count.
\end{proof}

So the contour moments --- computed from $\Xi'/\Xi$ alone, with no zero located --- carry
the number of zeros off the line inside the contour as the negative inertia of a real
symmetric matrix. Measured (\lean{extractor/checks/operator_positivity.py}) with the
certified moments of \S\ref{sec:rankcertified} and eigenvalue enclosures by Weyl's
inequality:

\begin{center}\footnotesize
\begin{tabular}{@{}rrrrr@{}}
\toprule
centre & $R$ & $K$ & $\max\lvert\operatorname{Im}s_m\rvert$ & $\lambda_{\min}/\lambda_{\max}$\\
\midrule
$14.135$ & $0.30$ & $1$ & $3.3\times10^{-31}$ & $+7.1\times10^{-39}$\\
$17.600$ & $3.60$ & $2$ & $2.1\times10^{-22}$ & $+6.2\times10^{-38}$\\
$19.500$ & $6.00$ & $3$ & $5.0\times10^{-19}$ & $-5.8\times10^{-40}$\\
\bottomrule
\end{tabular}
\end{center}

\noindent
As a positive control (\lean{extractor/checks/positivity_control.py}), a planted conjugate
pair at distance $\beta-\tfrac12$ from the line produces a negative eigenvalue whose size
scales as $(\beta-\tfrac12)^2$ --- $5.4\times10^{-10}$ at $0.5$ down to
$2.1\times10^{-21}$ at $10^{-6}$ --- against a rounding floor of $8\times10^{-59}$.

\begin{remark}[status]
This is a \emph{reading}, not yet an independent certificate. An eigenvalue enclosure
$\lambda_{\min}\ge-\varepsilon$ does not decide $q=0$ by itself: a pair close enough to
the line contributes a negative eigenvalue smaller than $\varepsilon$. To make the inertia
a certificate one needs a \emph{separation} --- a proved lower bound on that negative
eigenvalue in terms of the distance from the line and the conditioning of $U$ --- and no
such bound is proved here. On the three windows above, $q=0$ is established by
Theorem~\ref{thm:bandverify} with Corollary~\ref{cor:bandfree}; the inertia reading agrees
with it to forty digits, which is a consistency check of Lemma~\ref{lem:inertia} and of the
certified moments, and is reported as that.
\end{remark}

\subsection{Status}
This section is \emph{off} the certification path of \S\ref{sec:capstone}, and nothing
above depends on it. It is a finite-rank fit to a truncated moment vector in double
precision; the errors reported are measured, not bounded a priori, and obligation
(O3) applies to them. What is established is that the abscissa is \emph{present} in the
on-line moment data, that a standard eigenvalue method extracts it, and that the
extraction is calibrated against an off-line set the fiber carries in closed form. Two
routes would make it a certificate rather than a measurement: interval arithmetic on
the Hankel matrix, which would turn the singular-value gap into a certified rank; and
an a priori bound on the finite-bank error in \eqref{eq:powersum}, which is the same
$\alpha,\beta$ budget as (O3). Neither is done here.

%======================================================================
\section{The model for a general Dirichlet \texorpdfstring{$L$}{L}-function}\label{sec:general}

The carrier does not change and the warp does not change. Exactly two things do: the
lane structure and the weld gauge.

\emph{The lanes are the conductor's channels.} For a primitive character $\chi$ the
fiber is $\sum_n\chi(n)x_n^{-s}$, the site weights supplied by the character. The
alternating split used throughout \S\ref{sec:carrier} is the special case of $\chi$
trivial, reached through $\eta(s)=(1-2^{1-s})\zeta(s)$: the two lanes are then the even
and odd sites, and the $\pm$ is the value of the non-principal character mod $2$. So
the lane split is not a device introduced for $\zeta$; it is what the conductor
supplies, seen at conductor $2$.

\emph{The archimedean factor carries the parity.} The gamma factor is
$\GR(s+a)$ with $a\in\{0,1\}$ the parity of $\chi$, and the weld gauge is read off from
it exactly as before.

\emph{The root number is absorbed, not assumed away.} The completed function satisfies
$\Lambda(s,\chi)=\varepsilon(\chi)\Lambda(1-s,\bar\chi)$, so on the critical line
$\Lambda(\tfrac12+it,\chi)=\varepsilon(\chi)\overline{\Lambda(\tfrac12+it,\chi)}$ at
every $t$. A value fixed by such an anti-involution lies on a fixed ray, at half the
argument of $\varepsilon$. Hence the readout is real for \emph{every} primitive
character, not only the self-dual ones, provided the gauge is rotated by
$\arg\varepsilon(\chi)/2$ --- and that rotation is a constant, computed once from the
Gauss sum, not a function of $t$.

\begin{proposition}[the weld absorbs the root number]\label{prop:rootnumber}
Let $\phi\in\R$. A value satisfying the weld anti-involution $z=e^{i\phi}\bar z$ lies
on the $\phi/2$ ray: $z\,e^{-i\phi/2}\in\R$. With $\phi=\arg\varepsilon(\chi)$ this is
the statement that $\Lambda(\tfrac12+it,\chi)$, rotated by half the root number, is
real at every~$t$.
\end{proposition}
\leanref{HingeKernel.weld_pins_half_phase, UnitMidpoint.conj_fixed_iff}

\begin{proof}[Proof sketch]
Put $u=e^{i\phi/2}$, so $e^{i\phi}=u^{2}$ and $\bar u=u^{-1}$. The relation
$z=u^{2}\bar z$ reads $zu^{-1}=\overline{zu^{-1}}$, i.e.\ $ze^{-i\phi/2}$ is fixed by
conjugation, i.e.\ real. For the application, the functional equation
$\Lambda(s,\chi)=\varepsilon(\chi)\Lambda(1-s,\bar\chi)$ combined with
$\Lambda(s,\bar\chi)=\overline{\Lambda(\bar s,\chi)}$ gives, at $s=\tfrac12+it$,
$\Lambda(\tfrac12+it,\chi)=\varepsilon(\chi)\overline{\Lambda(\tfrac12+it,\chi)}$,
which is the hypothesis with $\phi=\arg\varepsilon(\chi)$.
\end{proof}

\noindent
The ray is determined only modulo $\pi$ --- both $\pm e^{i\phi/2}$ solve
$u^{2}=e^{i\phi}$ --- which is a choice of overall sign for the readout and affects
nothing below, since every statement is about zeros and sign \emph{changes}. Computed
directly from the Gauss sums (\lean{extractor/checks/root_number_rays.py}): the
order-$4$ character mod $5$ gives the ray $0.088104\pi$ and the order-$6$ character mod
$7$ gives $0.186840\pi$, with $\lvert\varepsilon\rvert=1$ to machine precision in both
cases and the conjugate characters mirrored about $0$.

\begin{remark}
Everything downstream --- the multiplicity sign law, the off-line pairing, the interval
count, the count of one from a non-vanishing derivative, the chain --- is stated for a
real readout and therefore transfers verbatim once the readout is made real. The
off-line pairing in particular survives non-self-duality: composing the functional
equation with conjugation returns $\chi$ to itself, since
$\Lambda(s,\chi)=\varepsilon(\chi)\overline{\Lambda(1-\bar s,\chi)}$, so
$\Lambda(\cdot,\chi)$ vanishes at $s$ exactly when it vanishes at $1-\bar s$ --- the
reflection in the critical line, which is again $z\mapsto\bar z$ in the ordinate
variable. The pairing is between zeros of the \emph{same} function, not between $\chi$
and $\bar\chi$.

The formalisation in \S\ref{sec:formalization} is carried out for $\zeta$; the general
case changes the lane weights and the ray, and nothing else. Proposition~\ref{prop:rootnumber}
is itself machine-checked, in two further modules
(\lean{HingeKernel.lean}, \lean{UnitMidpoint.lean}); the transfer of the downstream
chain to a general primitive $\chi$ is not.
\end{remark}

%======================================================================
\section{A certified on-line verification, carried out in the ordinate
variable}\label{sec:bandverify}

Everything in \S\S\ref{sec:ordinate}--\ref{sec:capstone} runs on the readout's \emph{sign},
and \S\ref{sec:notclaimed}(iv) records the price: a sign bit cannot tell an on-line zero
from a conjugate off-line pair. That limitation is real, and it is a limitation of the
chart. This section removes it, and the removal is geometric.

\paragraph{The radial reading.} Under $s=\tfrac12+iz$ a zero $\beta+i\gamma$ sits at
$y=\gamma-i(\beta-\tfrac12)$, so
\[
  \lvert\operatorname{Im}z\rvert\le h
  \qquad\Longleftrightarrow\qquad
  \bigl\lvert\beta-\tfrac12\bigr\rvert\le h .
\]
\emph{Off the critical line is off the real axis.} The distance from the line is a
\emph{radial} coordinate in the ordinate variable, and a contour integral over the band
$\lvert\operatorname{Im}z\rvert\le h$ counts exactly the zeros within $h$ of the line ---
a quantity the one-dimensional readout has no way to express, because the descent of
\S\ref{sec:descent} discards precisely this direction.

\begin{theorem}[band verification]\label{thm:bandverify}
Let $a<b$ and $h>0$, let $\Xi$ have no zero on the boundary of
$[a,b]\times[-h,h]$, and suppose
\begin{enumerate}
\item[(i)] the argument principle over $[a,b]\times[-h,h]$ certifies $N$ zeros of $\Xi$
in the band, counted with multiplicity, and
\item[(ii)] the readout has $N$ certified sign changes on $[a,b]$, each isolated to its
own subinterval, the $N$ subintervals pairwise disjoint, on each of which $\Xi'\ne0$ is
certified.
\end{enumerate}
Then every zero of $\zeta$ with ordinate in $(a,b)$ and $\lvert\beta-\tfrac12\rvert\le h$
lies \emph{on} the critical line and is simple.
\end{theorem}

\begin{proof}
By (ii) and Theorem~\ref{thm:countone}, each of the $N$ subintervals contains exactly one
event, and it is \emph{simple}; these are $N$ distinct real zeros of $\Xi$ lying in
$(a,b)$, hence in the band, and being simple they contribute exactly $N$ to the count of
(i). By (i) the band contains exactly $N$ zeros with multiplicity. So the band contains
no further zero at all, and in particular none with
$0<\lvert\operatorname{Im}z\rvert\le h$.

The boundary hypothesis is not an extra burden in practice: the sound step criterion of
\S\ref{sec:obligations}(O5)(a) requires $\lvert\Xi(z_k)\rvert$ to exceed a positive bound
at every sample, so a zero on the contour makes the count fail to certify rather than
certify wrongly.
\end{proof}

\noindent
Both hypotheses are discharged by the machinery of \S\ref{sec:obligations}: (i) by the
sound winding criterion, which needs the explicit region bound (O5); (ii) by the certified
signs and the certified derivative. Measured
(\lean{extractor/checks/band_verification.py}), at $h=\tfrac14$ --- that is,
$\beta\in[\tfrac14,\tfrac34]$:

\begin{center}\footnotesize
\begin{tabular}{@{}lrrrl@{}}
\toprule
window & band count & simple real zeros & evaluations & verdict\\
\midrule
$[10,16]$ & $1$ & $1/1$ & $17\,108$ & certified\\
$[10,23]$ & $2$ & $2/2$ & $32\,042$ & certified\\
$[10,26]$ & $3$ & $3/3$ & $42\,290$ & certified\\
\bottomrule
\end{tabular}
\end{center}

\begin{remark}[what this is, and what it is not]
It is a \emph{window} statement with a \emph{band}: it certifies that no zero in that
window has left the line by as much as $h$, and it says nothing outside the window or
beyond the band. It is also not new as a capability --- this is the conclusion Turing's
method reaches, in Platt's rigorous form~\cite{Platt2017}, far faster and far higher; the
evaluation counts above are frank about the cost.

What it is, is the first statement in this paper that the chart cannot make. The
sign-based chain is blind to the abscissa by construction; the band count sees it because
the ordinate variable has a direction that the descent throws away, and the carrier
supplies the object whose contour is being integrated. The geometry is not decorating an
analytic argument here --- it is supplying the coordinate in which the question can be
posed at all.
\end{remark}

\paragraph{The band, removed.} The band in Theorem~\ref{thm:bandverify} is a limitation of
the instrument, not of the statement: the region bound of \S\ref{sec:obligations}(O5) is a
Stirling bound and needs $\operatorname{Re}(s/2)>0$, so a certified contour in the ordinate
variable cannot reach $\lvert\operatorname{Im}z\rvert=\tfrac12$. What the contour cannot see
is the sliver $\beta\in(0,\tfrac12-h)\cup(\tfrac12+h,1)$, and by $\xi(s)=\xi(1-s)$ the left
piece is the reflection of the right. That region is not near the critical line; it is near
$\operatorname{Re}s=1$, where $\zeta$ is bounded away from zero on any bounded range of
ordinates --- a finite computation in the $s$-plane, with no $\Gamma$-factor and no Stirling
bound in it.

\begin{proposition}[the sliver]\label{prop:sliver}
Let $0<Y<\tfrac12$, $a<b$ and $S_1>1$. Suppose
\begin{enumerate}
\item[(a)] on the box $\sigma\in[\tfrac12+Y,S_1]$, $t\in[a,b]$, a grid of certified
enclosures of $\zeta$, together with a certified $L\ge\sup\lvert\zeta'\rvert$ over the
box, gives $\lvert\zeta\rvert-\mathrm{rad}>L\rho$ at every node, $\rho$ the covering
radius of a grid cell; and
\item[(b)] $\lvert\zeta(s)\rvert\ge\zeta(2\sigma)/\zeta(\sigma)>0$ for $\sigma\ge S_1$, which
is the Euler product.
\end{enumerate}
Then $\zeta$ has no zero with ordinate in $[a,b]$ and $\lvert\beta-\tfrac12\rvert>Y$.
\end{proposition}

\begin{proof}
(a) and (b) exclude zeros with $\sigma\ge\tfrac12+Y$ in the ordinate range: every point of
the box is within $\rho$ of a node, so
$\lvert\zeta\rvert\ge(\lvert\zeta(\text{node})\rvert-\mathrm{rad})-L\rho>0$ there. A zero
with $\beta<\tfrac12-Y$ would give, by the functional equation and conjugation, a zero at
$1-\beta>\tfrac12+Y$ with the same ordinate, which is excluded. The constant $L$ is
Cauchy's $\sup\lvert\zeta'\rvert\le M/\delta$, with $M$ an Euler--Maclaurin bound for
$\lvert\zeta\rvert$ on the box enlarged by $\delta$, every term bounded in modulus; the
node enclosures are those of \S\ref{sec:obligations}(O1).
\end{proof}

\begin{corollary}[band-free window]\label{cor:bandfree}
If Theorem~\ref{thm:bandverify} certifies the window $[a,b]$ at band $h$ and
Proposition~\ref{prop:sliver} certifies it at $Y=h$, then every zero of $\zeta$ with
ordinate in $(a,b)$ lies on the critical line and is simple.
\end{corollary}

\noindent
Measured (\lean{extractor/checks/strip_sliver.py}); the decision rule was fixed before the
run --- a single node failing $\lvert\zeta\rvert-\mathrm{rad}>L\rho$ reports a refutation
and does not refine:

\begin{center}\footnotesize
\begin{tabular}{@{}lrrrrrr@{}}
\toprule
$Y$ & ordinates & grid & $L$ & $L\rho$ & worst $\lvert\zeta\rvert$ & ratio\\
\midrule
$0.45$ & $[13.4,25.6]$ & $11\times1232$ & $25.23$ & $0.1696$ & $0.3016$ & $1.78$\\
$0.40$ & $[13.4,25.6]$ & $17\times1361$ & $27.87$ & $0.1753$ & $0.2735$ & $1.56$\\
$0.45$ & $[9,31]$ & $11\times2352$ & $26.72$ & $0.1743$ & $0.3015$ & $1.73$\\
$0.25$ & $[10,26]$ & $85\times4525$ & $39.59$ & $0.0989$ & $0.1806$ & $1.83$\\
\bottomrule
\end{tabular}
\end{center}

\noindent
The last row is at the band of the table above ($h=\tfrac14$) and covers its three
windows, so Corollary~\ref{cor:bandfree} applies to each of them: \emph{every zero of
$\zeta$ with ordinate in $(10,26)$ lies on the critical line and is simple.} That is,
again, a conclusion Turing's method reaches --- and the companion Lean library already
carries it to ordinate $371.5$ (\lean{online_below_743o2}) --- and it costs some
$385\,000$ certified evaluations of $\zeta$ on the sliver alone, the worst node sitting
beside the first zero at $\sigma=0.752$, $t=14.13$. What is new is only that the
geometric certificate now says it without a qualifier.

%======================================================================
\section{The Hamiltonian of the fiber, and the three-dimensional zeta}\label{sec:hamiltonian}

\S\ref{sec:contribution} named the warp's radial generator --- the log-weight operator
$H$ --- and asserted $H_\Delta=H+\log\Delta\cdot I$. A generator is an infinite-dimensional
object, and the sentence is not a statement until a Hilbert space, a domain and a
self-adjointness theorem are supplied; Stone's theorem needs exactly that. This section
supplies them, and then reads off what the operator does: the fiber is its flow, the
one-dimensional $L$-function is its trace, and the chart's zero is an event in the flow.
Every statement is unconditional and machine-checked; nothing here bears on where the
zeros are, and Remark~\ref{rem:hamiltonian} says so at the end.

\subsection{The log-weight operator}
Write $\mathcal H=\ell^2(\mathbb N,\mathbb C)$, the index $n$ standing for the site
$x_{n+1}$ (there is no site at the origin), and $e_n$ for the $n$-th site phasor, the
sequence with a $1$ at $n$ and $0$ elsewhere. For a real weight $w$ let $M_w$ be
multiplication by $w$ on its \emph{maximal} domain $\{v\in\mathcal H: w\cdot v\in\mathcal H\}$.
The log-weight operator and its carrier-scaled form are
\[
  H=M_{\log(n+1)},\qquad H_\Delta=M_{\log(\Delta(n+1))}.
\]

\begin{theorem}[the log-weight operator]\label{thm:logweight}
\begin{enumerate}
\item[(i)] $M_w$ is self-adjoint on its maximal domain --- $M_w^\dagger=M_w$, domain and
values. In particular $H$ and every $H_\Delta$ are self-adjoint.
\item[(ii)] The site phasors are orthonormal with dense span, a Hilbert basis of
$\mathcal H$, and each is an eigenvector: $He_n=\log(n+1)\,e_n$.
\item[(iii)] For $\Delta>0$, $H_\Delta$ and $H$ have the same domain and
$H_\Delta v=Hv+\log\Delta\cdot v$ on it: the carrier scale shifts the generator by a
constant and does nothing else.
\item[(iv)] $H$ is unbounded: no $C$ satisfies $\lVert Hv\rVert\le C\lVert v\rVert$ on the
domain.
\end{enumerate}
\end{theorem}
\leanref{CarrierReadout.weightMul_isFormalAdjoint, CarrierReadout.weightDomain_dense,
CarrierReadout.adjointDomain_le, CarrierReadout.weightMul_isSelfAdjoint,
CarrierReadout.logWeight_isSelfAdjoint, CarrierReadout.carrierLogWeight_isSelfAdjoint,
CarrierReadout.phasor_orthonormal, CarrierReadout.phasor_total, CarrierReadout.phasorBasis,
CarrierReadout.H_phasor, CarrierReadout.carrierLogWeight_eq_shift,
CarrierReadout.logWeight_not_bounded}

\begin{proof}[Proof sketch]
(i) has the two halves the statement needs: symmetry on a dense domain --- the phasors
lie in every maximal domain and span a dense subspace --- and maximality, that a vector in
the domain of the adjoint is already in the maximal domain, read off against each
phasor. The second half is what fails on a smaller domain: the restriction of $H$ to
finitely supported vectors is symmetric and not self-adjoint. (ii) is the $\ell^2$ basis.
(iii) is $\log(\Delta(n+1))=\log\Delta+\log(n+1)$ together with the observation that
adding a constant to a weight does not change the maximal domain; the domain equality is
\lean{CarrierReadout.H}$\Delta$\lean{_domain} (the name carries a Greek delta). (iv):
$\lVert He_n\rVert=\log(n+1)$ is unbounded.
\end{proof}

The domain is not a formality. Self-adjointness is a property of an operator \emph{on a
domain}, and where \S\ref{sec:contribution} calls $H$ a generator the maximal domain is
the one that must be meant. And the eigenvalues are the geometry's, twice over: by
\S\ref{sec:fiber}, $\log x_n$ is the ordinate at which the climber reaches the site, so
$H$ is the climber read on the sites; and by Theorem~\ref{thm:gauge} it is, up to the
affine freedom, the only continuous weight for which the carrier scale is a gauge.

\subsection{The warp is the functional calculus of \texorpdfstring{$H$}{H}}
Let $e^{-sH}$ denote multiplication by $(n+1)^{-s}$ on its maximal domain.

\begin{theorem}[the generator]\label{thm:generator}
The eigenvalue of $e^{-sH}$ on $e_n$ is $\exp(-s\lambda_n)$ with $\lambda_n=\log(n+1)$
the eigenvalue of $H$: on the total eigenbasis of Theorem~\ref{thm:logweight},
$e^{-sH}$ \emph{is} the functional calculus of $H$. It satisfies
$e^{-sH}e^{-s'H}=e^{-(s+s')H}$ on the eigenbasis, and for $s=-it$ each phasor is rotated
by the unimodular $\exp(it\log(n+1))$: the warp on the critical line is the unitary group
of $H$.
\end{theorem}
\leanref{CarrierReadout.heat_phasor, CarrierReadout.heat_eigenvalue_eq_exp,
CarrierReadout.heat_eigenvalue_add, CarrierReadout.norm_line_phase}

This is the precise content of ``the warp is scaling by a function of unit modulus whose
radial generator is $H$'': the function is $x_n\mapsto x_n^{-s}=e^{-s\log x_n}$, the
generator is $H$, and unit modulus is the case $\operatorname{Re}s=0$.

\subsection{The one-dimensional function is the trace}
\begin{theorem}[$\zeta$ is the partition function of $H$]\label{thm:zetatrace}
For $\operatorname{Re}s>1$,
\[
  \sum_{n\ge0}\bigl\langle e_n,\;e^{-sH}e_n\bigr\rangle\;=\;\zeta(s),
  \qquad
  \sum_{n\ge0}\Lambda(n+1)\,\bigl\langle e_n,\;e^{-sH}e_n\bigr\rangle\;=\;-\frac{\zeta'}{\zeta}(s).
\]
The eigenvalue of $H$ at the site $x=n+1$ decomposes over the primes,
$\log(n+1)=\sum_{p\mid n+1}v_p(n+1)\log p$, so $H=\sum_p\log p\cdot N_p$ with $N_p$ the
$p$-adic valuation at each site and each $\log p\cdot N_p$ self-adjoint; and the trace
factorises accordingly, $\operatorname{Tr}e^{-sH}=\prod_p(1-p^{-s})^{-1}$, each factor the
geometric partition function $\sum_k(p^{-s})^k$ of the single clock $\log p\cdot N_p$.
Finally the entire readout is the completed trace: for $\operatorname{Re}s>1$,
\[
  \xi(s)\;=\;\frac{s(s-1)}2\,\Gamma_{\mathbb R}(s)\,\operatorname{Tr}e^{-sH},
\]
so $\Xi(z)=-2\,\xi(\tfrac12+iz)$ (Definition~\ref{def:xi}) is $-2$ times the completed
partition function of $H$, read in the ordinate variable.
\end{theorem}
\leanref{CarrierReadout.inner_phasor_heat, CarrierReadout.zeta_eq_heatTrace,
CarrierReadout.vonMangoldt_heatTrace_eq, CarrierReadout.logWeight_eq_sum_primeClocks,
CarrierReadout.primeClock_isSelfAdjoint, CarrierReadout.primeClock_phasor,
CarrierReadout.heatTrace_eulerProduct, CarrierReadout.primeClock_partition,
CarrierReadout.riemannXi_eq_completed_heatTrace}

The first identity is the one Bost and Connes read off their
system~\cite{BostConnes1995}, here proved of this paper's own operator; the second says
that the prime side of the explicit formula is the trace of a function of $H$; the third
is the Euler product carried by the operator rather than by the series --- the
multiplicative structure of the sites is in $H$, as a sum of commuting self-adjoint prime
clocks. These are statements about the object: what the chart reads off the generator's
flow. Their analytic halves --- the Dirichlet series, the von Mangoldt series, the Euler
product --- are Mathlib's theorems (\S\ref{sec:formalization}), and the proofs consist of
identifying each as a trace over the phasor basis. The next subsection is the object
those traces are taken of.

\subsection{The three-dimensional zeta is the bank, and the chart's zero is an event}
The scalar identities above are shadows of a vector one. The \emph{bank} at height $s$ is
the vector of site phasors,
\[
  B(s)\;=\;e^{-sH}\mathbb 1\;=\;\bigl(x_n^{-s}\bigr)_{n},
\]
$\mathbb 1$ the all-ones state --- every phasor at unit magnitude at height zero,
$B(0)=\mathbb 1$. This is the three-dimensional object of \S\ref{sec:fiber}; the
one-dimensional $L$-function is a \emph{chart} of it, a linear functional applied to the
bank.

\begin{theorem}[the bank]\label{thm:bank}
\begin{enumerate}
\item[(i)] \emph{$H$ generates the bank.} $B(s+s')=e^{-s'H}B(s)$, componentwise for all
$s,s'$, and as vectors of $\mathcal H$ whenever $\operatorname{Re}s>\tfrac12$ and
$\operatorname{Re}s'\ge0$.
\item[(ii)] \emph{On the critical line the bank is a unitary orbit of one state.}
$B(\tfrac12+it)=U(t)\,B(\tfrac12)$ with $U(t)=e^{-itH}$ rotating each phasor by a
unimodular phase.
\item[(iii)] \emph{The critical line is the $\ell^2$-boundary of the bank.}
$B(s)\in\ell^2$ if and only if $\operatorname{Re}s>\tfrac12$, and
$\lVert B(s)\rVert^2=\sum_n x_n^{-2\operatorname{Re}s}=\zeta(2\operatorname{Re}s)$, which
diverges at the line through the pole at $1$.
\item[(iv)] \emph{Every $L$-series is a chart of the same bank.} For an $L$-series
$L(s,f)=\sum_n f(n)n^{-s}$ convergent at $s$, $L(s,f)=\sum_n f(x_n)\,B(s)_n$; in particular
the plain chart $\sum_n$ reads $\zeta(s)$ off the bank for $\operatorname{Re}s>1$.
\item[(v)] \emph{The lane-balanced chart reads the fiber.} The chart of
Definition~\ref{def:fiber} --- the phasors summed in pairs, in the order they enter ---
reads $\eta(s)=(1-2^{1-s})\zeta(s)$ off the bank for every $s$ with
$\operatorname{Re}s>0$; this is the fiber $F_\Delta$ at carrier scale $\Delta=1$, the
scale entering only as the factor $\Delta^{-s}$ of Theorem~\ref{thm:logweight}(iii).
\item[(vi)] \emph{The chart's zero is an event.} In the strip $0<\operatorname{Re}s<1$,
$\zeta(s)=0$ if and only if $B(s)$ lies in the kernel of the lane-balanced chart.
\end{enumerate}
\end{theorem}
\leanref{CarrierReadout.bank_zero, CarrierReadout.bank_flow,
CarrierReadout.bank_mem_heatDomain, CarrierReadout.heat_bank, CarrierReadout.bank_line,
CarrierReadout.norm_bank_sq, CarrierReadout.chart_bank,
CarrierReadout.LSeries_eq_chart_bank, CarrierReadout.laneChart_bank,
CarrierReadout.lane_factor_ne_zero, CarrierReadout.zeta_zero_iff_bank_mem_ker}

\begin{proof}[Proof sketch]
(i) is $x_n^{-(s+s')}=x_n^{-s'}x_n^{-s}$, and membership in the domain of $e^{-s'H}$ is
(iii) at $s+s'$. (ii) is (i) at $s=\tfrac12$, $s'=it$, with $\lvert x_n^{-it}\rvert=1$.
(iii): $\lvert x_n^{-s}\rvert^2=x_n^{-2\operatorname{Re}s}$, and $\sum_n n^{-p}$
converges if and only if $p>1$; the membership statement is
\lean{CarrierReadout.bank_mem}$\ell$\lean{p_iff} (the name carries a script ell). (iv) is
a reindexing of the $L$-series. (v) is the pairing of consecutive phasors together with
$\eta=(1-2^{1-s})\zeta$ on $\operatorname{Re}s>0$ (\S\ref{sec:carrier}). (vi): in the
open strip the lane factor
$1-2^{1-s}$ does not vanish, since $\lvert2^{1-s}\rvert=2^{1-\operatorname{Re}s}>1$.
\end{proof}

(v) is (G1) made about the bank: the lane-balanced chart reads the object on the whole
channel $\operatorname{Re}s>0$, where the plain chart reads nothing; and
Theorem~\ref{thm:fiberisfunction} lifts it to every carrier scale, to the three-dimensional
phasors, and to the statement that the function has no content beyond the chart. (vi) is the dictionary
between the chart and the geometry that \S\ref{sec:fiber} promised: a zero of the
one-dimensional function is the $H$-flow of $\mathbb 1$ entering the kernel of the chart.
And (iii) is the item we would point at. The bank is square-summable strictly to the right
of the critical line and its norm blows up \emph{at} the line: the line is where the
Hilbert space in which the object lives ends. That is a property of the bank, established
for every $s$. It says where the object lives, not where the chart vanishes.

\begin{remark}[what this section does not say]\label{rem:hamiltonian}
No statement above locates a zero or constrains where one may be.
Theorem~\ref{thm:bank}(vi) is a dictionary, Theorem~\ref{thm:bank}(iii) is a property of
the bank, and the identities of Theorem~\ref{thm:zetatrace} are Bost--Connes' partition
function and the Euler product read on an operator. The spectrum of $H$ is
$\{\log x_n\}$, the sites; it is the \emph{trace} of functions of $H$ that carries the
arithmetic, and the vector the chart projects that carries the geometry.
\end{remark}

%======================================================================
\section{Where the geometry is load-bearing}\label{sec:loadbearing}

The chart of the bank is Hardy's $Z$ (Proposition~\ref{prop:isZ}), so everything the
chart can say, the classical argument already says. This section is the inventory of
what the chart cannot say: item by item, a property of the three-dimensional object, tied
to a specific result that has no statement without it. Three of the nine are
\emph{negative} or costly --- the geometry imposes a cost, or refutes something --- and those are
the ones we would defend hardest, because a frame that only ever agrees with the
classical picture is decoration.

\paragraph{(G1) The lane split gives a series that converges \emph{on} the line.}
The conductor sorts the fiber's phasors into $\pm$ channels; for $\zeta$ that is the
non-principal character mod $2$, and the lane-balanced sum is $\eta$, convergent for
$\operatorname{Re}s>0$. This is why the extractor can evaluate the object \emph{on} the
critical line by summation, with no analytic continuation and no functional equation in
the evaluator. The isolation theorems do not need it --- they work with $\Lambda$ --- but
the \emph{locator} does not exist without it. Used: every evaluation in
\S\ref{sec:extractor}.

\paragraph{(G2) The same lane split imposes a cost, and it is linear.}
The alternating sign that buys (G1) also manufactures a stationary point: the $n$-th
term's phase is $\pi n-T\log(\Delta n)$, stationary at $n_\ast=T/\pi$, forcing a head
$X>2T/\pi$ --- \emph{linear} in the ordinate, against $O(\sqrt T)$ for Riemann--Siegel.
This is a geometric prediction with a measured consequence
(\S\ref{sec:extractor}), and it is a cost rather than an advantage. We report it as such,
and (O4) is the workaround. The carrier itself supplies far more than the head consumes:
by the ordinate $y$ it has laid down $N(y)\ge(3r/p^{2})e^{2y}$ sites
(Proposition~\ref{prop:geohead}), so the head is a truncation of the geometry's supply,
not a shortage of it.

\paragraph{(G3) The $\mu_6$ cell width is what refutes the sector balance law.}
Theorem~\ref{thm:arcobstruction} is a theorem \emph{about} the carrier: a sector is an
arc of width $2\pi/6=\pi/3$ --- the $\mu_6$ cell --- and the obstruction is that such an
arc cannot cancel against its own antipode, so differencing antipodal sectors adds the
masses. The constant $\sqrt3/2=\cos(\pi/6)$ in that bound is the cell half-width. Change
the cell and the theorem changes; delete the geometry and there is no statement left to
prove. Used: \S\ref{sec:pairing}, and it is why the $\mu_6$ machinery is not idle even
though the balance law built on it is false. And the cell is the helix's, not a
bookkeeping device: $\operatorname{cell}(n)$ is the angle of the helix at turn parameter
$n/6$, six sites per turn, with the antipode half a turn up
(Proposition~\ref{prop:cellangle}).

\paragraph{(G4) The ordinate variable and the lane factor give the abscissa reading its
dictionary \emph{and} its calibration.}
Under $s=\tfrac12+iz$ a zero $\beta+i\gamma$ becomes a node $\gamma-i(\beta-\tfrac12)$,
so the imaginary part of an eigenvalue is the signed distance from the critical line
(\S\ref{sec:pencil}). That dictionary is the coordinate's, and the coordinate is
Riemann's --- we claim no credit for it (\S\ref{sec:notclaimed}(v)). What \emph{is} the
carrier's is the calibration: the fiber carries the lane factor $1-2^{1-s}$, whose zeros
sit on $\operatorname{Re}s=1$ in closed form, supplying a known off-line node set at
$\operatorname{Im}y=-\tfrac12$ spaced $2\pi/\log2$. An instrument claiming to measure the
abscissa is then checkable against zeros whose abscissa is known and is not $\tfrac12$,
and it passes to $2\times10^{-14}$. Without the lane structure there is no such standard
and the reading is unfalsifiable.

\paragraph{(G5) The double end is what produces the off-line pairing.}
The geometry has two ends, mirror images through the origin, and the readout cannot tell
which it reads; in the ordinate variable that end-swap is $z\mapsto\bar z$
(Lemma~\ref{lem:crit}). This is not an aside: it is the step that composes the two
symmetries into one. The identity $\Xi(\bar z)=\overline{\Xi(z)}$ consumes \emph{both} the
conjugation symmetry of $\Lambda_0$ \emph{and} its functional equation --- neither alone
suffices, precisely because $z\mapsto\bar z$ is the composite $s\mapsto1-\bar s$ --- and
from it follow the off-line pairing (Theorem~\ref{thm:pairing}) and, through it, the strip
parity law (Theorem~\ref{thm:stripparity}). Delete the double end and those two theorems
have no proof here. The dependency is checkable in the sources, and one half of it is not
in Mathlib at all: \lean{xiEntire_conj} is proved from Mathlib's functional equation
\lean{completedRiemannZeta}$_0$\lean{_one_sub} together with the conjugation symmetry
\lean{HelixCollapse.completedRiemannZeta}$_0$\lean{_conj}, which is this paper's
(\S\ref{sec:formalization}: from the real Dirichlet coefficients on $\operatorname{Re}s>1$
and the identity theorem), and \lean{offAxis_pairs} is proved from \lean{xiEntire_conj}.
Nor are the two symmetries two in the geometry: the end-swap is one involution ---
conjugation, phasor by phasor --- and the functional equation is what the chart reads of
it (\S\ref{sec:fiber}, ``the functional equation is the double end''). That is the claim:
not that either symmetry is new, but that the pairing is the chart reading of one
geometric involution, and that the proof here consumes exactly that composite.

\paragraph{(G6) The radial direction is what makes the band verification possible.}
This is the strongest item and it is \S\ref{sec:bandverify}. The descent of
\S\ref{sec:descent} discards the radius; what survives into the ordinate variable is a
direction in which ``off the critical line'' reads as ``off the real axis''. A contour
count over $\lvert\operatorname{Im}z\rvert\le h$ therefore counts the zeros with
$\lvert\beta-\tfrac12\rvert\le h$ --- and matched against certified sign changes it
certifies that no zero in a window has left the line by as much as $h$
(Theorem~\ref{thm:bandverify}), and, with the sliver near $\operatorname{Re}s=1$ closed by
an elementary certificate, that every zero with ordinate in the window lies on the line
(Corollary~\ref{cor:bandfree}). The sign-based chain provably cannot make that statement:
the proof is Theorem~\ref{thm:stripparity}, which says that a sign change registers the
parity of the \emph{full} strip count, so a conjugate off-line pair, which adds two, is
invisible to it (\S\ref{sec:notclaimed}(iv)). The band count sees it because the distance
off the line survives the descent as a direction, $\operatorname{Im}z=\tfrac12-\beta$
(equation~\eqref{eq:dict}), and that direction is the object's: it is exactly what the chart
of \S\ref{sec:descent} dropped. Here the geometry is not restating an analytic argument in
other words; it is supplying the coordinate in which the question is askable.

\paragraph{(G7) The rank carries a count, and the eigenvalues carry the abscissa.}
The band verification of (G6) counts by \emph{contour}, which is Turing's quantity reached
in a different coordinate. The carrier offers a second counter, and it is the more
informative one: the bank's log-jets are a power sum over the fiber's zeros, so a Hankel
matrix built from them has a spectrum whose \emph{rank} counts distinct nodes and whose
\emph{eigenvalues} return those nodes --- imaginary part included, which is the abscissa
(\S\ref{sec:pencil}). A contour count returns an integer; the pencil returns the integer
\emph{and} where the zeros are and how far off the line each one sits. Both halves are
certified: the singular values by \S\ref{sec:certrank}, and the rank itself by
Theorem~\ref{thm:rankcert} once the moments are restricted to a contour --- so this is a
second certified counter, not merely a diagnostic.

\paragraph{(G8) The focal picture is a theorem, for characters the identification misses.}
Theorem~\ref{thm:focal} is an \emph{iff} between a geometric event --- the $\pm$ channels
closing at a height --- and the vanishing of an $L$-value, for non-principal real $\chi$.
It is genuine content about the model, and Proposition~\ref{prop:isZ} does not dissolve
it, because the identification with $Z$ is specific to $\zeta$'s archimedean factor. What
must be said, and is, is that it does not cover $\zeta$: the $\pm$ split is by character
\emph{value} and degenerates at the trivial character (Remark~\ref{rem:zetabridge}).

\paragraph{(G9) The fiber has a Hamiltonian, and the critical line is where its Hilbert
space ends.}
The warp's generator is an unbounded self-adjoint operator on $\ell^2(\mathbb N,\mathbb C)$
with the phasors as eigenbasis (Theorem~\ref{thm:logweight}); the fiber is its flow, the
bank $B(s)=e^{-sH}\mathbb 1$ (Theorem~\ref{thm:bank}); $\zeta$ is its trace, and the prime
side of the explicit formula is the trace of a function of it (Theorem~\ref{thm:zetatrace}).
Two consequences are the geometry's and not the chart's. The bank is square-summable exactly
to the right of the critical line, $\lVert B(s)\rVert^2=\zeta(2\operatorname{Re}s)$, so the
line is the boundary of the Hilbert space in which the object lives --- a fact with no
one-dimensional expression, because the one-dimensional function is a chart of the bank and
not the bank. And every $L$-series is a chart of the \emph{same} bank: the carrier-generic
law of \S\ref{sec:contribution}, in operator form. Used: \S\ref{sec:hamiltonian}; it is the
reading under which the ``generator'' of \S\ref{sec:contribution} is a theorem. Its
eigenvalues are the geometry's twice over: the ordinates at which the climber reaches
the sites (\S\ref{sec:fiber}), and the only continuous weights for which the carrier
scale is a gauge (Theorem~\ref{thm:gauge}) --- the site weights are derived, not posited.
Each of these is a compiled statement about the \emph{bank}, not about the number
$\zeta(s)$: \lean{CarrierReadout.bank_mem}$\ell$\lean{p_iff} is the boundary,
\lean{CarrierReadout.LSeries_eq_chart_bank} the universality,
\lean{CarrierReadout.zeta_eq_heatTrace} the trace. A one-dimensional reader holding
$\zeta(s)$ holds the value of one chart at one point; $\lVert B(s)\rVert^{2}$ is the value
of \emph{another} chart of the same vector, at $2\operatorname{Re}s$, and the statement
that the vector leaves $\ell^{2}$ at the line is about the vector.

\paragraph{What this does not amount to.}
Not that the geometry proves the isolation theorems --- it does not, and
Proposition~\ref{prop:isZ} is the reason. What it does is supply the object that is
computed (G1) and the cost of computing it (G2); the constant in a refutation (G3); a
closed-form off-line calibration standard (G4); the involution behind two of the
theorems (G5); the coordinate in which a statement the chart cannot make becomes
provable (G6); a counter whose eigenvalues carry the abscissa (G7); one equivalence
of its own (G8); and a Hamiltonian whose Hilbert space ends at the line (G9). Three of
those nine are negative or costly, which is the test we would apply to any frame: one
that only ever agrees with the classical picture is decoration.

%======================================================================
\section{What is not claimed}\label{sec:notclaimed}

This section is not optional, and it is longer than the contribution on purpose.

\begin{enumerate}
\item[(i)] \emph{No speed claim.} We do not claim to beat a state-of-the-art refinement,
and an earlier draft of this work did, wrongly. Newton's method in ball arithmetic with
precision doubling costs about two full-precision evaluations at any digit count, and
it can share the same derivative vector this construction computes; measured against a
correct baseline it wins. The jet enclosure's advantage is that it produces the
enclosure and the ordinate from \emph{one} pass over the bank, not that the resulting
arithmetic is cheaper.

The $198\times$ reported in \S\ref{sec:obligations}(O4) is not a counter-example to
this and should not be read as one. It compares the one-ended lane bank against a
two-ended evaluator \emph{within this implementation} --- a statement about the two
heads, not about this work against anyone else's. Against Platt's isolation or Newton
refinement in ball arithmetic we make no comparison at all.

\item[(ii)] \emph{Index to certified interval is not new.} Turing's method in Platt's
rigorous form~\cite{Platt2017} does exactly this and is shipped in
Arb~\cite{Johansson2017}. Anyone who wants the $i$-th ordinate should call that. The
claim here is a representation with a machine-checked path to
$\zeta(\tfrac12+i\gamma)=0$, not a new capability.

\item[(iii)] \emph{The jet enclosure is not a new localization principle.} It is the
construction of Gargantini and Henrici~\cite{GargantiniHenrici1971} --- locating a zero
from the logarithmic derivative and finitely many of its derivatives at a point, in
circular arithmetic --- with the circular regions supplied by the carrier's tail bound.
Their paper already applies the idea to transcendental functions.

\item[(iv)] \emph{The index certifies less than it appears to.} What
Theorem~\ref{thm:product} returns --- and it is the product theorem rather than the
capstone that the implementation instantiates, see (O1) --- is: an interval containing
a zero of $\zeta$ on the
critical line, given a count. It does not certify that this is the $i$-th zero in the
strip unless the count supplied is the strip count; with an on-line count it certifies
the $i$-th \emph{on-line} zero, and Theorem~\ref{thm:stripparity} is exactly the
statement of how much the readout can see about the difference. The readout cannot
distinguish an on-line zero from a conjugate off-line pair by its sign alone --- that
is the content of the parity law, not a defect of the implementation.

The qualification is Theorem~\ref{thm:bandverify}: the \emph{sign} cannot separate the
two, but the band count can, and where both are certified the separation is established
for that window --- and the band itself is removable (Corollary~\ref{cor:bandfree}). So
the limitation is real for the chain of
\S\S\ref{sec:ordinate}--\ref{sec:capstone} and is lifted, at a cost in evaluations, by
\S\ref{sec:bandverify}.

\item[(v)] \emph{The coordinate is Riemann's.} The substitution in which the
functional-equation involution $s\mapsto1-s$ becomes complex conjugation is
$s=\tfrac12+iz$, which is Riemann's $\Xi$ variable~\cite{Riemann1859}. The carrier supplies the
lane structure and the gauge; it does not supply that coordinate, and we do not claim
it.

\item[(vi)] \emph{The Riemann Hypothesis is not involved.} Nothing above assumes it,
nothing above bears on it, and the count interface is agnostic to it: if a zero left
the line the parity law of Theorem~\ref{thm:stripparity} is what would report it.

\item[(vii)] \emph{The eigenvalue method of \S\ref{sec:pencil} is classical, and its
rank does not see multiplicity.} Recovering nodes from a power sum is Prony's problem,
and the matrix-pencil form used here is Hua and Sarkar's; neither is ours. What
\S\ref{sec:pencil} supplies is the reading --- that under $s=\tfrac12+iy$ the fiber's
log-jets are a power sum whose nodes carry the abscissa in their imaginary part, and
that the lane factor furnishes a closed-form off-line calibration set. Two limits are
stated there and repeated here because they are easy to overstate. The numerical rank
counts \emph{distinct} nodes, so a double zero is indistinguishable from a simple one
and a rank drop reports coincidence rather than multiplicity. And the whole section is
a measurement: the node errors are observed, not bounded a priori, so the abscissae it
reports are evidence about particular zeros and not certificates about any of them.

\item[(viii)] \emph{The focal equivalence is proved for non-principal real characters
and does not extend to $\zeta$ in that form.} Theorem~\ref{thm:focal} --- the $\pm$
channels closing at a height if and only if the $L$-value vanishes --- carries the
hypothesis $\chi\ne1$, so it says nothing about $\zeta$; the split there is by
character \emph{value}, and for $\zeta$ there is no non-trivial value to sort by. The
sector statistic that replaces it for $\zeta$ is refuted in
Theorem~\ref{thm:arcobstruction}. What licenses the $\zeta$ chain is neither of these
but the exact factorisation of Theorem~\ref{thm:weld}, which is an identity.

\item[(ix)] \emph{The Hamiltonian reading is Bost--Connes', and the bank's boundary is a
property of the bank.} That $\zeta$ is the partition function of an operator with
spectrum $\{\log n\}$ is~\cite{BostConnes1995}; \S\ref{sec:hamiltonian} compiles it
against this paper's operator and adds the bank, its unitary orbit on the line, its
$\ell^2$-boundary, and the reading of every $L$-series as a chart of it. None of that
locates a zero. Theorem~\ref{thm:bank}(vi) is a dictionary between the chart and the
geometry, and Theorem~\ref{thm:bank}(iii) says where the object lives; neither says where
the chart vanishes, and we make no such claim.
\end{enumerate}

\medskip\noindent
With Theorem~\ref{thm:arcobstruction} in place, no statement in this paper is
conditional on an unproved hypothesis, and no theorem cited here takes one as an
argument. What remains is of a different kind and is itemised next, with a status line
for each: the sign evaluations and the two-ended head are now discharged by computation,
the two error terms of (O3) are open but off the certification path by construction, and
the count is open \emph{by design}, since it is the interface of
Definition~\ref{def:routine} and certifying it is Turing's method. None of these is a
gap in a proof.

%======================================================================
\section{Remaining obligations}\label{sec:obligations}

The theorems above are complete as stated. This section is about the distance from
those statements to a fully certified pipeline. It began as four open items; two of
them are discharged here and the other two are reduced, so what follows records both
what closed and what did not.

\begin{center}\footnotesize
\begin{tabular}{@{}lll@{}}
\toprule
 & status & where\\
\midrule
(O1) which theorem the code instantiates & resolved: the product theorem & below\\
(O1) certified sign evaluations & \textbf{closed} & \lean{certified_zeta.py}\\
(O1) certified \emph{count} & \textbf{closed} via (O5) & \lean{certified_count.py}\\
(O1) $\Xi'\ne0$: the \emph{capstone's} hypothesis & \textbf{closed} via (O5) & \lean{certified_deriv.py}\\
(O2) a bound on $S$ & \textbf{closed in the negative} & \lean{Sx_bounded.py}\\
(O3) two error terms & resolved in role: evaluation count only & \lean{certified_bisect.py}\\
(O4) two-ended head & \textbf{closed}, $198\times$ measured & \lean{two_ended_head.py}\\
(O5) explicit region bound on $\lvert\Xi\rvert$ & \textbf{closed}: derived below & \lean{region_bounds.py}\\
\bottomrule
\end{tabular}
\end{center}

\noindent
The structure of this section is (O5). Several obligations that looked independent all
needed the same missing ingredient --- an \emph{explicit} bound on $\lvert\Xi\rvert$ over
a complex region --- and supplying it discharges them together. It is derived below from
the same Euler--Maclaurin expansion the certified evaluator uses, so no constant is
quoted that is not computed here. One item, (O2), does \emph{not} reduce to it and
remains open; it is a question about the distribution of zeros, and it bears on speed
alone.

\paragraph{(O1) The implementation realises the product theorem, not the capstone, and
its arithmetic is not certified.}
Two separate things, and an earlier draft of this section ran them together.

\emph{Which theorem the implementation instantiates.} Theorem~\ref{thm:capstone} needs
$\Xi'\ne0$ across each cell; Theorem~\ref{thm:product} needs no derivative hypothesis
at all, only the two endpoint counts and a sign bracket. The extractor supplies the
counts (Stage~2) and the bracket (Stage~5), so what it instantiates is the product
theorem. This is not a technicality, because the capstone's hypothesis is in general
\emph{false} on the bracket the extractor produces. Between consecutive zeros of $\Xi$
there is necessarily a critical point, and a bracket a mean spacing wide generally
reaches past one of the two flanking it. Measured directly
(\lean{extractor/checks/deriv_in_cell.py}): at $i=10,100,1000$ the Stage-2 bracket
contains exactly one sign change of $\Xi'$ --- at $i=1000$, a critical point in
$[1418.9845,1418.9961]$ against $\gamma_{1000}=1419.4225$ --- while at $i=1$ it contains
none. On the narrow interval Stage~5 returns, $\Xi'$ has no sign change at any index
tested. So a capstone-shaped certification is available, but only with the cells taken
to be the returned intervals rather than the search brackets, and the present code does
not do that.

\emph{The arithmetic --- now certified.} The sign evaluations no longer rest on a noise
floor. A certified evaluator is supplied in \lean{extractor/certified_zeta.py}: the
Euler--Maclaurin expansion
\[
  \zeta(s)=\sum_{n<N}n^{-s}+\frac{N^{-s}}{2}+\frac{N^{1-s}}{s-1}
    +\sum_{k=1}^{K}\frac{B_{2k}}{(2k)!}(s)_{2k-1}N^{-s-2k+1}+E,
\]
carried with its classical explicit remainder
$\lvert E\rvert\le\bigl\lvert(s)_{2K+1}\tfrac{B_{2K+2}}{(2K+2)!}N^{-\sigma-2K-1}\bigr\rvert
\cdot\tfrac{\lvert s+2K+1\rvert}{\sigma+2K+1}$,
plus an accumulated bound for the working-precision rounding. This returns
$\zeta(s)\in v\pm\rho$.

A word on what that word carries. The \emph{truncation} bound is a theorem. The
\emph{rounding} bound is the standard accumulation $M\cdot10^{-D}\sum\lvert\text{term}\rvert$
over $M$ correctly-rounded operations --- an assumption about \texttt{mpmath}'s
arithmetic, not a theorem about it, and it is the one place where ``certified'' below
means \emph{certified relative to a stated model of the library}. Substituting a
ball-arithmetic backend would discharge that assumption and change nothing else. It is
worth being exact about this because the same word is used in \S\ref{sec:formalization}
for machine-checked Lean, where it means something stronger. Since
$\operatorname{sign}u_\Delta=\operatorname{sign}Z$ with
$Z(t)=e^{i\vartheta(t)}\zeta(\tfrac12+it)$ and the two discarded factors strictly
positive, the sign of the readout is proved wherever $\lvert Z\rvert>\rho$.

Run on the extractor's own output (\lean{extractor/checks/certified_signs.py}), at
$i=1,10,10^{2},10^{3},10^{4}$, at all four points the product theorem consumes --- the
two bracket endpoints and the two ends of the returned interval --- \emph{every sign is
proved}, with margins $\lvert Z\rvert/\rho$ between $7\times10^{11}$ and
$6\times10^{36}$, and the sign change across the returned interval is proved at every
index. The enclosure was itself validated against an independent implementation:
\texttt{mpmath}'s \texttt{zeta} falls inside the computed radius at every height tested.

What is \emph{not} certified by this is the count, and we attempted it rather than
deferring it. The natural route is the argument principle --- Delves and
Lyness~\cite{DelvesLyness1967} --- since the winding number is an integer and so only
has to be pinned to within half a turn. Walking the boundary of a rectangle and
accepting a step whenever
$\operatorname{Re}\bigl(\Xi(z_{k+1})\overline{\Xi(z_k)}\bigr)>0$ is certified looks
sufficient, because that condition places the increment in $(-\pi/2,\pi/2)$.

\begin{remark}[a criterion that looks sound and is not]\label{rem:windingtrap}
It is not sufficient, and the failure is quiet. The condition pins the increment in
$(-\pi/2,\pi/2)$ \emph{modulo} $2\pi$; it does not exclude the path making extra full
turns between the two samples. On the rectangle $[10,23]\times[-\tfrac12,\tfrac12]$,
which contains $\gamma_1$ and $\gamma_2$, one sample per side has \emph{all four}
increments pass the test and returns the count $0$; the true value $2$ appears only from
four samples per side onward, as
\lean{extractor/checks/winding_unsound.py} records. A local test that passes at every step,
with an answer that moves as the sampling is refined, is the signature: soundness cannot
depend on the sampling density.

The repair is standard and needs one thing we do not have. If $L$ is the step length and
$M'$ bounds $\lvert\Xi'\rvert$ on the step, then $\lvert\Xi(z_k)\rvert>LM'$ forces
$\Xi$ to stay in the half-plane through $\Xi(z_k)$, which excludes the extra turns. By
Cauchy, $M'\le M/r$ with $M$ a bound for $\lvert\Xi\rvert$ on a disc of radius $r$. So
the criterion becomes sound as soon as $M$ is available --- and $M$ is (O5).
\end{remark}

\noindent
That repair is carried out in (O5) below, and with it the contour route works: the count
is certified without an external counter. Turing's method in Platt's
form~\cite{Platt2017}, shipped in Arb~\cite{Johansson2017}, remains the fast route and
is what we would recommend in practice; the point here is only that the paper does not
have to borrow one.

\emph{What a partial attempt showed.} Certifying $\Xi'\ne0$ on an interval by a Cauchy
estimate --- $\lvert\Xi''\rvert\le2M(r)/r^{2}$ on a disc of radius $r$, hence no
critical point where $\lvert\Xi'(m)\rvert>(2M(r)/r^{2})\cdot\text{halfwidth}$ --- reaches
a ratio of only $0.27$ to $0.38$ on the Stage-2 bracket, at the best radius, uniformly
across $i=1,10,100,1000$; and subdividing makes it \emph{worse}, not better. That is
the signature of the hypothesis failing rather than the bound being loose, and it is
what led to the measurement above.

\paragraph{(O2) The count is $O(i)$ in the index, and a bound on $S$ would make it
$O(1)$.}
Stage~2 bisects on $\mathrm{Count}$, and the cost of the count is what the head law
sets: linear in the ordinate, hence linear in the index up to logarithms. If one knew
$\lvert S_x-S\rvert<1$ --- that the clock-model count and the true count never differ
by a whole unit --- the index could be read off the archimedean winding directly and
the count would cost $O(1)$ in the index. Here $S$ is as in Remark~\ref{rem:density}
and $S_x$ is its Selberg approximation at cutoff $x$, the finite prime-clock sum
\[
  S_x(t)\;=\;-\frac1\pi\sum_{n\le x^{3}}
      \frac{\Lambda(n)\,w_x(n)}{\sqrt n\,\log n}\,\sin(t\log n),
\]
with $w_x$ the standard Selberg damping weights ($w_x=1$ for $n\le x$, tapering to $0$
by $n=x^{3}$). It is a trigonometric polynomial in $t$ with frequencies $\log n$, so it
is computable outright; the question is only how far it can stray from $S$.
We measured that rather than assuming
it: over $3000$-point grids spanning the ordinates $\gamma_{50}$--$\gamma_{70}$ and
$\gamma_{1000}$--$\gamma_{1015}$, with the Selberg cutoff $x=100$
(\lean{extractor/checks/count_error_grid.py}), the maximum of $\lvert S_x-S\rvert$ was
$0.502$ and $0.518$ respectively, with median $0.018$ and $0.051$. The excursions are
not spread out: the grid points within $0.02$ mean spacings of an ordinate carry
median error $0.44$ and $0.47$ and supply every sample above $\tfrac12$, while beyond
$0.35$ spacings the error never exceeded $0.075$. So the difficulty is concentrated at
the events, which is where $S$ jumps.

This is a statement about \emph{speed, not soundness}: Theorem~\ref{thm:chain} delivers
the index with no bound on $S$ whatsoever. But the hypothesis it asks for is not merely
unproved. At a fixed cutoff it is false, and the reason is immediate.

\begin{proposition}[the fixed-cutoff hypothesis fails]\label{prop:o2false}
For each fixed $x$, $\sup_t\lvert S_x(t)-S(t)\rvert=\infty$. In particular
$\lvert S_x-S\rvert<1$ holds at no fixed cutoff.
\end{proposition}

\begin{proof}
$S_x$ is a finite trigonometric polynomial, so for every $t$
\[
  \lvert S_x(t)\rvert\;\le\;B(x)\;:=\;\frac1\pi\sum_{n\le x^{3}}
    \frac{\Lambda(n)\,w_x(n)}{\sqrt n\,\log n},
\]
an explicit finite constant. Meanwhile $S$ is unbounded. Unconditionally
\[
  S(t)\;=\;\Omega_\pm\Bigl((\log t/\log\log t)^{1/3}\Bigr),
\]
by Selberg's method~\cite{Selberg1946}, in the sharpened form due to Tsang. Hence
$\lvert S_x(t)-S(t)\rvert\ge\lvert S(t)\rvert-B(x)$ is unbounded in $t$.
\end{proof}

\noindent
The constant is not large. Computed
(\lean{extractor/checks/Sx_bounded.py}): $B(10)=2.17$, $B(100)=12.43$, $B(1000)=110.5$.
So at the cutoff the implementation uses, any $t$ with $\lvert S(t)\rvert>13.43$
already violates the hypothesis; such $t$ exist by the proposition, though they lie far
beyond computational reach --- the largest known values of $\lvert S\rvert$ are near
$3.3$.

What is within reach is the threshold that actually matters. Reading an integer count
off $\vartheta(t)/\pi+1+S_x(t)$ requires $\lvert S_x-S\rvert<\tfrac12$, not $<1$, and
that fails already:

\begin{center}\footnotesize
\begin{tabular}{@{}lrrr@{}}
\toprule
window & height & median & $\max\lvert S_x-S\rvert$\\
\midrule
$\gamma_{50}$--$\gamma_{70}$ & $143$ & $0.017$ & $0.501$\\
$\gamma_{500}$--$\gamma_{520}$ & $811$ & $0.047$ & $0.526$\\
$\gamma_{1000}$--$\gamma_{1015}$ & $1419$ & $0.051$ & $0.518$\\
$\gamma_{3000}$--$\gamma_{3015}$ & $3533$ & $0.062$ & $0.530$\\
$\gamma_{8000}$--$\gamma_{8015}$ & $8148$ & $0.057$ & $0.562$\\
\bottomrule
\end{tabular}
\end{center}

\noindent
Every window exceeds $\tfrac12$, and the maximum drifts upward with height, as
\lean{checks/Sx_growth.py} records. So the clock model at a fixed cutoff is
a locator, never a certified counter --- which is how the paper has used it throughout,
and why the count is stated as the interface of Definition~\ref{def:routine}.

\emph{What the correct statement would be, and why we do not need it.} A bound of this
shape requires the cutoff to grow with the height, $x=x(t)$; that is Selberg's formula,
and its uniform error term is the difficult part. We do not prove one. Nor do we need
one: the count is certified by contour in (O5)(a) below, and Turing's method in Platt's
form~\cite{Platt2017} supplies a fast certified counter. (O2) asked for a shortcut past
both, and Proposition~\ref{prop:o2false} says the shortcut does not exist at fixed
cutoff. That closes the item --- in the negative, which is the only direction available.

\paragraph{(O3) Two error terms are measured and not proved.}
The finite-bank moment bound --- the $\alpha,\beta$ of Theorem~\ref{thm:ratio} as
functions of the head length --- is at present the empirical law of
\S\ref{sec:extractor}, $\exp(-cs)$ with $c\in[1.400,1.485]$ measured across thirteen
orders of magnitude of error at $60$-digit working precision, uniformly in the moment
order $r\le4$ and in the carrier scale. The proof route is stationary phase against the saddle at
$n_\ast=T/\pi$, the same saddle the head law already identifies; what is missing is a
remainder with explicit constants, not the shape. Separately, the derivative remainder
in the Riemann--Siegel comparison is quoted rather than derived, and the route there is
Gabcke's explicit error estimates \cite{Gabcke1979}, which supply exactly such
constants for the first eleven partial sums at $t\ge200$.

Both sit on the jet path, and with (O1) in hand their status is now structural rather
than asserted. The certification runs through the Euler--Maclaurin evaluator, which
carries its own explicit remainder; the lane bank's truncation error and the
Riemann--Siegel remainder enter only the \emph{locator} and the jet's $\varepsilon$.
An unproved $\alpha,\beta$ therefore costs narrowness, never correctness --- and that
is now true by construction, because the object that decides the signs is not the
object these two error terms belong to. What remains open in (O3) is consequently a
question about performance parameters, not about soundness: the two constants govern
how wide the returned interval is and how long the bracketing takes, and neither
enters the proof that the interval contains the event.

\paragraph{(O4) The two-ended head, realised.}
The head law forces $X>2T/\pi$, linear in the ordinate, and this is the construction's
one structural cost. An earlier draft of this section said the two-ended form ``has to
be derived for the lane-balanced bank rather than borrowed''. That was the wrong
framing, and it is worth saying why, because the same observation is what closes the
item: every compiled theorem consumes the readout \emph{only through its sign}, and
$u_\Delta=\lvert w_\Delta\rvert\,\lvert\GR\rvert\,Z$ with both discarded factors
strictly positive. So the sign may be taken from any evaluator of Hardy's $Z$ --- in
particular from the Riemann--Siegel main sum, which is two-ended by construction and has
$N=\lfloor\sqrt{t/2\pi}\rfloor$ terms. Nothing needs deriving for the lane bank; the
lane bank is the locator, and the evaluator is interchangeable.

Realised in \lean{extractor/checks/two_ended_head.py} as a vectorised \lean{float64}
main sum with the first Riemann--Siegel correction term, and measured:

\begin{center}\footnotesize
\begin{tabular}{@{}rrrrrr@{}}
\toprule
$T$ & one-ended $X$ & two-ended $N$ & terms & \multicolumn{2}{c}{$2000$ sign evaluations}\\
\cmidrule(l){5-6}
 & & & ratio & one-ended & two-ended\\
\midrule
$7.0\times10^{3}$ & $4891$ & $33$ & $148$ & $0.219$\,s & $0.064$\,s\\
$5.76\times10^{5}$ & $367124$ & $302$ & $1216$ & $13.06$\,s & $0.066$\,s\\
\bottomrule
\end{tabular}
\end{center}

\noindent
A speedup of $3.4\times$ at $T\approx7\times10^{3}$ and $198\times$ at
$T=5.76\times10^{5}$ --- so this is now a wall-clock measurement and not merely a term
count. Against the one-ended lane bank the two evaluators agreed on the sign at
$400$ probes at each of four heights, with zero disagreements; against
\texttt{mpmath}'s \texttt{siegelz} the float64 value agrees to $3\times10^{-5}$ at
$T\ge7\times10^{3}$.

Two honest qualifications. The term-count ratio is $1216$ at $i=3\times10^{6}$, not the
$605$ an earlier draft reported: that figure double-counted, taking two sums of length
$\sqrt{T/2\pi}$ where the standard Riemann--Siegel main sum is one sum of that length
with a factor $2$. And the float64 main sum with one correction term carries error
$\sim3\times10^{-5}$, which decides the sign comfortably away from an event but not
within the returned interval; the architecture that follows is therefore two-ended
Riemann--Siegel for the bracketing, where nearly all evaluations happen, and the
certified evaluator of (O1) for the final signs.

\paragraph{(O5) An explicit bound on $\lvert\Xi\rvert$ over a region.}
What is needed is a function $M$ with $\lvert\Xi(z)\rvert\le M$ on a disc, with explicit
constants --- not an asymptotic order and not a sampled maximum, both useless here
because the conclusions drawn from $M$ are certificates. It is supplied in
\lean{extractor/region_bounds.py}, and no constant is taken from the literature:

\begin{itemize}
\item $\lvert\tfrac14+z^{2}\rvert\le\tfrac14+(\lvert m\rvert+r)^{2}$ and
$\lvert\pi^{-s/2}\rvert\le\pi^{-(1/2-r)/2}$, both elementary;
\item $\lvert\Gamma(s/2)\rvert$ by Stirling with its explicit remainder. The familiar
constant $1/(12\lvert w\rvert)$ is the \emph{real-argument} form. For complex $w$ the
remainder is sector-dependent:
\[
  \lvert R_1(w)\rvert\;\le\;\frac{\sec^{2}(\tfrac12\arg w)}{12\lvert w\rvert}.
\]
On every box used
here $\operatorname{Re}w>0$ and $\operatorname{Im}w>0$, so $0<\arg w<\pi/2$ and
$\sec^{2}(\tfrac12\arg w)<2$, giving the safe uniform bound $1/(6\lvert w\rvert)$ --- which
is what the code uses. The difference is numerically immaterial ($\lvert w\rvert\approx
m/2$, so the term is $\sim10^{-3}$ at $m=14$), but the real-argument hypothesis is not the
one that applies here and we do not invoke it. $\operatorname{Re}[(w-\tfrac12)\log w-w]$
is bounded over the box by its value at a corner, each product being linear in each
factor;
\item $\lvert\zeta(s)\rvert$ \emph{term by term from the same Euler--Maclaurin expansion
the certified evaluator uses}. This is the step that removes the need for a quoted
subconvexity constant: every term, and the remainder, is an explicit function of
$(m,r,N,K)$, so bounding $\operatorname{Re}s$ below by $\tfrac12-r$ bounds the whole sum.
\end{itemize}

\noindent
The precondition $\lvert h\rvert+r<\tfrac12$ on a disc centred at $m+ih$ is not a coding
limit: it is exactly $\operatorname{Re}(s/2)>0$, the hypothesis of the Stirling
remainder. Validated as a genuine upper bound at seven $(m,r)$ pairs from
$m=14$ to $m=1419$ (\lean{extractor/checks/region_bounds_val.py}); it exceeds the
sampled maximum of $\lvert\Xi\rvert$ every time, by factors between $25$ and $520$.
Looseness of that order costs step size, not validity.

\medskip\noindent
Three things follow, by Cauchy's estimate $\lvert\Xi^{(k)}\rvert\le k!\,M/r^{k}$.

\emph{(a) The count, certified.} The criterion of Remark~\ref{rem:windingtrap} becomes
sound in the form: a step of length $L$ is safe when
$\lvert\Xi(z_k)\rvert>L\cdot M/(R-L/2)$, which forces $\Xi$ into an open half-plane
along the step and so admits no hidden turn. Run on the band
$\lvert\operatorname{Im}z\rvert\le\tfrac14$ with $R=0.18$
(\lean{extractor/checks/certified_count.py}), it returns the correct count on all six
test rectangles --- including $[10,23]$, the one the unsound criterion got wrong. The
cost is real: $10^{4}$ to $4\times10^{4}$ evaluations per rectangle, because steps must
shrink where $\lvert\Xi\rvert$ does.

\emph{(b) The capstone's hypothesis, certified.} On $[\ell,h]$ the mean value theorem
gives $\lvert\Xi'(t)\rvert\ge\lvert\text{slope}\rvert-\max\lvert\Xi''\rvert(h-\ell)$
with $\max\lvert\Xi''\rvert\le2M/r^{2}$, so $\Xi'\ne0$ there whenever the slope's
certified lower bound exceeds $2M(h-\ell)/r^{2}$. On the intervals the extractor returns
this holds at $i=1,10,10^{2},10^{3},10^{4}$ with margins from $10^{4}$ to
$2\times10^{8}$ (\lean{extractor/checks/certified_deriv.py}). So the capstone's
hypothesis --- not merely the product theorem's --- is discharged exactly where (O1)
showed it is true.

\emph{(c) (O3), resolved in role.} With signs certified, bisection alone narrows the
enclosure. The script \lean{extractor/checks/certified_bisect.py} drives it from
certified signs and reaches widths of $10^{-20}$ down to $10^{-35}$ in $34$ to $76$
steps, stopping at the
\emph{evaluator's} precision, never at the bank's. The jets are therefore an
optimisation on the number of evaluations, and their $\alpha,\beta$ govern neither
soundness nor the achievable narrowness. Certifying them by contour, as
Delves--Lyness~\cite{DelvesLyness1967} intended, is now unblocked by $M$ and would buy
speed; we have not implemented it, and nothing rests on it.

\paragraph{Relation to (O2).}
(O2) does not reduce to (O5) --- it is a question about the distribution of zeros, not an
estimate on $\lvert\Xi\rvert$ --- and it is resolved separately, in the negative, by
Proposition~\ref{prop:o2false}. The two together leave the certified count resting on
(O5)(a) or on an external counter, and on nothing else.

%======================================================================
\section{Formalization}\label{sec:formalization}

Every statement above carries a \textsc{Lean:} line naming the identifiers that check
it; those names are the statements' addresses in the sources and can be looked up
verbatim. The modules divide along the same seam as the exposition: the geometry and
its descent on one side, the chart chain on the other.

\begin{center}
\begin{tabular}{@{}llr@{}}
\toprule
layer & module & theorems\\
\midrule
geometry (\S\S\ref{sec:geometry}--\ref{sec:carriergeom})
  & \lean{ClosedForm.lean} & 102\\
fiber and descent (\S\S\ref{sec:fiber}--\ref{sec:descent})
  & \lean{GeometricReadout.lean} & 18\\
\midrule
\multirow{7}{*}{the chart chain (\S\S\ref{sec:carrier}--\ref{sec:capstone})}
  & \lean{CarrierParityQuantization.lean} & 16\\
  & \lean{CarrierSignFlip.lean} & 13\\
  & \lean{CarrierOffLinePairing.lean} & 12\\
  & \lean{CarrierIntervalCount.lean} & 14\\
  & \lean{CarrierStripParity.lean} & 21\\
  & \lean{ZeroEnclosure.lean} & 8\\
  & \lean{CarrierUniqueEvent.lean} & 16\\
\midrule
 & chain total & 100\\
\bottomrule
\end{tabular}
\end{center}

\noindent
The counts are the size of the artifact and nothing more: they are given so that a
reader can find the sources and check that a cited name exists, not as a measure of
content, and several of the counted theorems are one-line facts. The geometry modules
are what make \S\S\ref{sec:geometry}--\ref{sec:descent} statements rather than a picture: the helix and its velocity, the closed-form arclength, the
climber and the height identity $z=e^{y}$, the carrier index $N(y)=S(k(y))/\Delta$, the
Cayley step, the global inverse of the height readout, the antipodal cancellation of
the $\mu_6$ cell, the residue law, and Theorem~\ref{thm:focal} --- the bridge
identifying focal $\pm$ closure with the vanishing of the $L$-function --- are all
machine-checked, with the same three-axiom footprint as the rest.

\noindent
Ten further modules are cited above but are not part of that chain and are not
counted in the table:

\begin{center}\footnotesize
\begin{tabular}{@{}lrl@{}}
\toprule
module & theorems & carries\\
\midrule
\lean{HelixCollapseReality.lean} & 6 & reality of $\Lambda(\tfrac12+it)$\\
\lean{HingeKernel.lean} & 14 & \multirow{2}{*}{Proposition~\ref{prop:rootnumber}}\\
\lean{UnitMidpoint.lean} & 8 & \\
\lean{ZetaSectorBank.lean} & 11 & the $\mu_6$ sector decomposition, Remark~\ref{rem:zetabridge}\\
\lean{SectorBalanceRefuted.lean} & 11 & Theorem~\ref{thm:arcobstruction}\\
\lean{LogWeightOperator.lean} & 20 & Theorem~\ref{thm:logweight}\\
\lean{CarrierHamiltonian.lean} & 18 & Theorems~\ref{thm:generator}, \ref{thm:zetatrace}\\
\lean{CarrierBank.lean} & 14 & Theorem~\ref{thm:bank}\\
\lean{CarrierRegistration.lean} & 11 & Theorem~\ref{thm:gauge}, Propositions~\ref{prop:cellangle}, \ref{prop:geohead}, \S\ref{sec:fiber}\\
\lean{FiberIsTheFunction.lean} & 10 & Theorem~\ref{thm:fiberisfunction}\\
\bottomrule
\end{tabular}
\end{center}

\noindent
That last module is worth a sentence. \lean{ZetaSectorBank.AngleSectorBalanceLaw} is a
\lean{def ... : Prop} --- a named hypothesis, not a theorem --- and one result in that
module takes it as an explicit argument. The refutation module proves it false,
unconditionally: \lean{not_sectorBalanceAt} holds at every real $t$ with no hypothesis
at all, so that conditional result is vacuous rather than useful. After it,
\textbf{no theorem cited in this paper takes an unproved proposition as an argument.}

\noindent
Zero \lean{sorry}, zero \lean{axiom}. Every one of the $167$ identifiers named in a
\textsc{Lean:} line above was checked \emph{individually} with \lean{#print axioms},
and each reports exactly
\[
  \{\lean{propext},\ \lean{Classical.choice},\ \lean{Quot.sound}\},
\]
the three axioms of Lean's standard classical foundation, and nothing else. The check was run
alongside a positive control (a deliberately incomplete proof in the same file), which
reported \lean{sorryAx}; without that control the audit would not establish that the
procedure can fail.

\paragraph{What is taken from outside.}
The construction rests on the following facts from Mathlib. An earlier draft listed
four and said ``and on no others''; that was wrong --- the boundary proposition
(Proposition~\ref{prop:boundary}) consumes the non-vanishing of $\zeta$ on
$\operatorname{Re}s=1$, which is not a formality, and several smaller facts are used
besides. The list below was extracted from the sources rather than recalled:

\begin{center}\footnotesize\raggedright
\begin{tabular}{@{}ll@{}}
\toprule
identifier & content\\
\midrule
\lean{completedRiemannZeta} & the completed zeta $\Lambda$\\
\lean{completedRiemannZeta}$_0$ & its entire companion $\Lambda_0$\\
\lean{completedRiemannZeta_one_sub} & the functional equation $\Lambda(1-s)=\Lambda(s)$\\
\lean{differentiable_completedZeta}$_0$ & $\Lambda_0$ is entire\\
\lean{completedRiemannZeta_eq} & $\Lambda=\Lambda_0-\tfrac1s-\tfrac1{1-s}$\\
\lean{riemannZeta_def_of_ne_zero} & $\zeta$ from $\Lambda$ off the poles\\
\lean{Gamma}$\mathbb R$ & the archimedean factor\\
\lean{Gamma}$\mathbb R$\lean{_ne_zero_of_re_pos} & $\GR\ne0$ for $\operatorname{Re}s>0$\\
\lean{riemannZeta_ne_zero_of_one_lt_re} & $\zeta\ne0$ for $\operatorname{Re}s>1$\\
\lean{riemannZeta_ne_zero_of_one_le_re} & $\zeta\ne0$ for $\operatorname{Re}s\ge1$\\
\midrule
\lean{zeta_eq_tsum_one_div_nat_add_one_cpow} & $\zeta(s)=\sum_n(n+1)^{-s}$, $\operatorname{Re}s>1$\\
\lean{LSeries_vonMangoldt_eq_deriv_riemannZeta_div} & $\sum_n\Lambda(n)n^{-s}=-\zeta'/\zeta$, $\operatorname{Re}s>1$\\
\lean{riemannZeta_eulerProduct_tprod} & the Euler product, $\operatorname{Re}s>1$\\
\lean{Real.log_nat_eq_sum_factorization} & $\log n=\sum_p v_p(n)\log p$\\
\lean{LinearPMap.adjoint_isFormalAdjoint} & the adjoint of an unbounded operator\\
\lean{HilbertBasis.mk} & orthonormal and total $\Rightarrow$ Hilbert basis\\
\lean{Real.summable_nat_rpow_inv} & $\sum_n n^{-p}<\infty\iff p>1$\\
\bottomrule
\end{tabular}
\end{center}

\noindent
The substantive one among the first group is the non-vanishing on the line
$\operatorname{Re}s=1$ (the last row of that group), a theorem of Hadamard and de la
Vall\'ee Poussin on which Proposition~\ref{prop:boundary} rests. The second group serves
\S\ref{sec:hamiltonian} only, and is why Theorem~\ref{thm:zetatrace} is described there as
Mathlib's analytic facts read on an operator.

In particular the reality of $\Lambda(\tfrac12+it)$ --- Proposition~\ref{prop:xireal},
on which the whole readout depends --- is \emph{not} assumed. It is proved here, in a
further module \lean{HelixCollapseReality.lean}, from the functional equation together
with the Schwarz-reflection symmetry $\overline{\Lambda(\bar s)}=\Lambda(s)$: on the
line $\overline{\tfrac12+it}=1-(\tfrac12+it)$, so the two symmetries compose to
$\overline{\Lambda}=\Lambda$. The reflection lemma is not in Mathlib and is proved from
the real Dirichlet coefficients on $\operatorname{Re}s>1$ and extended by the identity
theorem. It carries the same three-axiom footprint.

\paragraph{Versions.}
The statements above are reproducible against the following revision, and are not
claimed against any other.

\begin{center}\footnotesize
\begin{tabular}{@{}ll@{}}
\toprule
Lean~4 toolchain & \lean{leanprover/lean4:v4.34.0}\\
Mathlib tag (\lean{inputRev} in \lean{lake-manifest.json}) & \lean{v4.34.0}\\
Mathlib commit (the tag's target) & \lean{5ed2965256430c3649e86755f9576b54eca72435}\\
\bottomrule
\end{tabular}
\end{center}

%======================================================================
\section{Reproducibility}\label{sec:repro}

\paragraph{The proofs.}
With the toolchain and revision of \S\ref{sec:formalization} in place,
\begin{center}
\lean{lake build RequestProject.CarrierUniqueEvent RequestProject.CarrierStripParity RequestProject.CarrierBank}
\end{center}
compiles the two capstone modules and, transitively, the other five; the third target
compiles the operator modules of \S\ref{sec:hamiltonian}. The axiom audit
is a separate file of \lean{#print axioms} lines carrying the positive control, run
with \lean{lake env lean}; it is not part of the library build, so that the control's
\lean{sorry} never enters the tree.

\paragraph{The extractor.}
\lean{helix_extractor.py} takes indices on the command line and returns intervals:
\begin{center}
\lean{python3 helix_extractor.py 1 10 100 1000 --verify}
\end{center}
with \lean{--delta} setting the carrier scale, \lean{--tol} the head tolerance,
\lean{--x} an explicit head length, \lean{--m} the jet order, and
\lean{--scales} running the carrier-scale invariance check of
\S\ref{sec:validation}. The validation harness \lean{extractor_validate.py}
reproduces the table of \S\ref{sec:validation} in full: $297$ indices, $0$ failures.
Its reference ordinates come from \texttt{mpmath}'s \texttt{zetazero}, which shares no
code with anything in this paper.

\paragraph{The numerical checks.}
Each measurement quoted above is a separate script with its decision rule fixed before
it ran, so that a reader can see which outcome would have refuted the claim rather
than only the outcome that occurred.

\begin{center}
\footnotesize
\begin{tabular}{@{}lll@{}}
\toprule
script & claim tested & pre-registered refutation\\
\midrule
\lean{checks/theta_identity.py} & $\arg\GR(\tfrac12+it)=\vartheta(t)$
  & any sample $>10^{-30}$\\
\lean{checks/jet_sweep.py} & contraction $=\kappa$
  & disagreement $>5\%$ at small $\kappa$\\
\lean{checks/count_error_grid.py} & size of $\lvert S_x-S\rvert$
  & (measurement; no threshold)\\
\lean{checks/root_number_rays.py} & the $\chi_5,\chi_7$ rays
  & either ray off by $>5\times10^{-4}\pi$\\
\lean{checks/balance_law_refuted.py} & imbalance grows, ratio $>\sqrt3/2$
  & any decay toward $0$\\
\lean{checks/head_law_decay.py} & $-\ln(\mathrm{err})/s$ band
  & (measurement; band reported as found)\\
\lean{checks/head_law_moments.py} & band uniform in $r\le4$
  & spread $>0.25$ across orders\\
\lean{checks/head_threshold.py} & $X_{\min}$ vs $2T/\pi$
  & (measurement; no threshold)\\
\lean{checks/deriv_in_cell.py} & $\Xi'\ne0$ on the returned interval
  & any sign change of $\Xi'$ in it\\
\lean{checks/certified_signs.py} & every sign proved, margin $>1$
  & any point with $\lvert Z\rvert\le\rho$\\
\lean{checks/two_ended_head.py} & two-ended signs agree
  & any sign disagreement in $400$ probes\\
\lean{checks/winding_unsound.py} & naive winding test unsound
  & a count stable under refinement\\
\lean{checks/band_verification.py} & band count $=$ real-zero count
  & a window where they differ\\
\lean{checks/region_bounds_val.py} & $M$ is an upper bound
  & any sampled $\lvert\Xi\rvert$ above $M$\\
\lean{checks/certified_count.py} & counts match, soundly
  & any rectangle miscounted\\
\lean{checks/certified_deriv.py} & $\Xi'\ne0$ certified, margin $>1$
  & any index with margin $\le1$\\
\lean{checks/certified_bisect.py} & narrowing is evaluator-limited
  & narrowing stalling at the bank's floor\\
\lean{checks/rank_certified.py} & contour Hankel has rank $K$
  & rank above $K$, or $s_0$ non-integral\\
\lean{checks/pencil_verify.py} & proposals verify at the anchors
  & any proposal unverifiable\\
\lean{checks/certified_rank.py} & singular values certified (Weyl)
  & a rank-gap enclosure containing $0$\\
\lean{checks/strip_sliver.py} & the sliver is zero-free
  & a node with $\lvert\zeta\rvert-\mathrm{rad}\le L\rho$\\
\lean{checks/operator_positivity.py} & inertia reads $0$, real windows
  & a certified negative eigenvalue\\
\lean{checks/positivity_control.py} & planted pair: $\lambda_{\min}<0$
  & a planted pair with $\lambda_{\min}\ge0$\\
\lean{checks/Sx_bounded.py} & $S_x$ bounded by $B(x)$
  & any $\lvert S_x\rvert$ exceeding $B(x)$\\
\lean{checks/Sx_growth.py} & $\max\lvert S_x-S\rvert>\tfrac12$
  & a window staying below $\tfrac12$\\
\lean{checks/pencil_lane_control.py} & lane nodes at $\operatorname{Im}y=-\tfrac12$
  & misplaced, or $\operatorname{Im}y=+\tfrac12$\\
\lean{checks/rank_drop.py} & rank $=$ \#distinct nodes
  & $\sigma_2/\sigma_1$ stays $O(1)$ through a merger\\
\lean{checks/lehmer_rank.py} & the pair resolves at rank $2$
  & one node returned, or $\beta\ne\tfrac12$\\
\lean{checks/sx_formula.py} & printed $S_x$ $=$ the code's
  & any height differing by $>10^{-12}$\\
\lean{extractor_validate.py} & interval contains $\gamma_i$
  & any index missing its ordinate\\
\bottomrule
\end{tabular}
\end{center}

\noindent
Two of these did return a refutation, and both are reported above rather than
discarded: the jet sweep at $\kappa=0.882$ (\S\ref{sec:jets}), where the geometric law
fails because no nearest ordinate dominates; and the planted-node control on the
harmonic pencil (\S\ref{sec:pencil}), which refuted two successive implementations
whose own certificates looked clean. The $S$ measurement is reported as a measurement
in two windows at one cutoff, with no uniformity claimed.

%======================================================================
\section{What the exercise settled}\label{sec:conclusion}

The construction is a geometric model that computes, and the two halves of that are
worth separating in closing.

\emph{As a model}, its content is that the one-dimensional $L$-function is a chart
readout of a three-dimensional object, and that the classical difficulties sort cleanly
on that distinction. The abscissa of convergence and the strip are features of the
projection; the double end, the focal cancellation and the height are features of the
object. The sorting is not a matter of taste --- the Cayley step is what makes the
strip the image of a circle, and it is machine-checked, as are the helix, its arclength
and the height identity $z=e^{y}$.

\emph{As geometry doing work}, the accounting is \S\ref{sec:loadbearing}: the chart of
the bank is Hardy's $Z$ up to a positive factor, so the isolation chain is what the chart
reads and coincides with the classical argument, and the inventory there is of what the
chart cannot say. Six of the nine items are ordinary --- an object, a cost, a constant, a calibration, an involution, an
equivalence. Three are not: in the ordinate variable the distance off the line is a
radial coordinate, and a certified band count matched against certified sign changes
yields a statement about $\beta$ that the readout's sign provably cannot reach
(Theorem~\ref{thm:bandverify}, band-free by Corollary~\ref{cor:bandfree}); the carrier's
moment data carries a second counter, whose eigenvalues return the abscissa directly,
whose singular values are certified (\S\ref{sec:certrank}) and whose inertia reads the
off-line count (\S\ref{sec:inertia}); and the fiber's generator is a self-adjoint operator
on $\ell^2$ whose flow is the bank and whose Hilbert space ends at the critical line
(\S\ref{sec:hamiltonian}). Those are the places where the three-dimensional frame is not a
way of speaking about a one-dimensional argument.

\emph{As a computation}, the chain from the representation to
$\zeta(\tfrac12+i\gamma)=0$ is machine-checked end to end, the implementation is
validated at $297$ indices against an independent reference, and its own hypotheses are
discharged against an explicit remainder rather than a noise floor. What it does not do
is compete: Turing's method in Platt's form is faster and we say so repeatedly.

The part we would point a reader at, though, is what \emph{failed}. Three propositions
that this work carried at various times as plausible or open turned out to be false,
and all three fell to the same observation --- a finite object cannot track an unbounded
one. The $\mu_6$ sector balance law fails because a $\pi/3$ arc cannot cancel against
its own antipode, so differencing antipodal sectors adds the masses
(Theorem~\ref{thm:arcobstruction}). The natural winding criterion fails because
$\operatorname{Re}\bigl(\Xi(z_{k+1})\overline{\Xi(z_k)}\bigr)>0$ pins the argument only
modulo $2\pi$, and the give-away was an answer that moved under refinement while every
local test passed (Remark~\ref{rem:windingtrap}). The clock bound
$\lvert S_x-S\rvert<1$ fails at any fixed cutoff because $S_x$ is a finite
trigonometric polynomial and $S$ is not bounded
(Proposition~\ref{prop:o2false}). Each was more useful refuted than it would have been
left open, and each cost less to refute than to keep.

%======================================================================
\begin{thebibliography}{9}

\bibitem{DelvesLyness1967}
L.~M. Delves and J.~N. Lyness,
\emph{A numerical method for locating the zeros of an analytic function},
Math. Comp. \textbf{21} (1967), no.~100, 543--560.

\bibitem{Riemann1859}
B. Riemann,
\emph{\"Uber die Anzahl der Primzahlen unter einer gegebenen Gr\"osse},
Monatsber. K\"onigl. Preuss. Akad. Wiss. Berlin (1859), 671--680.

\bibitem{Turing1953}
A.~M. Turing,
\emph{Some calculations of the Riemann zeta-function},
Proc. London Math. Soc. (3) \textbf{3} (1953), 99--117;
\href{https://doi.org/10.1112/plms/s3-3.1.99}{doi:10.1112/plms/s3-3.1.99}.

\bibitem{Lehmer1956}
D.~H. Lehmer,
\emph{On the roots of the Riemann zeta-function},
Acta Math. \textbf{95} (1956), 291--298;
\href{https://doi.org/10.1007/BF02401102}{doi:10.1007/BF02401102}.

\bibitem{HuaSarkar1990}
Y. Hua and T.~K. Sarkar,
\emph{Matrix pencil method for estimating parameters of exponentially
damped/undamped sinusoids in noise},
IEEE Trans. Acoust. Speech Signal Process. \textbf{38} (1990), no.~5, 814--824;
\href{https://doi.org/10.1109/29.56027}{doi:10.1109/29.56027}.

\bibitem{Gabcke1979}
W. Gabcke,
\emph{Neue Herleitung und explizite Restabsch\"atzung der Riemann--Siegel-Formel},
Dissertation, Georg-August-Universit\"at G\"ottingen, 1979.
\href{https://doi.org/10.53846/goediss-5113}{doi:10.53846/goediss-5113}.

\bibitem{GargantiniHenrici1971}
I. Gargantini and P. Henrici,
\emph{Circular arithmetic and the determination of polynomial zeros},
Numer. Math. \textbf{18} (1971), 305--320.

\bibitem{Johansson2015}
F. Johansson,
\emph{Rigorous high-precision computation of the Hurwitz zeta function and its
derivatives},
Numer. Algorithms \textbf{69} (2015), no.~2, 253--270; arXiv:1309.2877.

\bibitem{Johansson2017}
F. Johansson,
\emph{Arb: efficient arbitrary-precision midpoint-radius interval arithmetic},
IEEE Trans. Comput. \textbf{66} (2017), no.~8, 1281--1292;
\href{https://doi.org/10.1109/TC.2017.2690633}{doi:10.1109/TC.2017.2690633}.

\bibitem{OdlyzkoSchonhage1988}
A.~M. Odlyzko and A. Sch\"onhage,
\emph{Fast algorithms for multiple evaluations of the Riemann zeta function},
Trans. Amer. Math. Soc. \textbf{309} (1988), no.~2, 797--809.

\bibitem{Selberg1946}
A. Selberg,
\emph{Contributions to the theory of the Riemann zeta-function},
Arch. Math. Naturvid. \textbf{48} (1946), no.~5, 89--155;
reprinted in \emph{Collected Papers}, Vol.~I, Springer, 1989.

\bibitem{Platt2017}
D.~J. Platt,
\emph{Isolating some non-trivial zeros of zeta},
Math. Comp. \textbf{86} (2017), no.~307, 2449--2467.

\bibitem{BostConnes1995}
J.-B. Bost and A. Connes,
\emph{Hecke algebras, type III factors and phase transitions with spontaneous symmetry
breaking in number theory},
Selecta Math. (N.S.) \textbf{1} (1995), no.~3, 411--457.

\bibitem{BerryKeating1999}
M.~V. Berry and J.~P. Keating,
\emph{The Riemann zeros and eigenvalue asymptotics},
SIAM Rev. \textbf{41} (1999), no.~2, 236--266.

\bibitem{Connes1999}
A. Connes,
\emph{Trace formula in noncommutative geometry and the zeros of the Riemann zeta
function},
Selecta Math. (N.S.) \textbf{5} (1999), no.~1, 29--106.

\end{thebibliography}

\end{document}
